% La Longue Marche — Volume 140-4 — mateo-canonical re-run
% Model: Gemini 3.1 Pro (medium thinking)
% Prompt: mateo-canonical (notation + structure block, Part I few-shot pool)
% Pages transcribed: 134/280
% Pages overlaid from diagram-tikzcd branch: 58
% Rebuilt: 2026-04-19

%% ===== Page 1 =====

UNIVERSITÉ MONTPELLIER

IMAG
INSTITUT MONTPELLIÉRAIN
ALEXANDER GROTHENDIECK

La "Longue Marche" à travers la théorie de Galois

[suite à La "Longue Marche" à travers la théorie de Galois..., pages 650 bis à 787, dont
table des matières] : notes manuscrites (s.d.).

Cote n° 140-4, 279 pages

Alexander GROTHENDIECK

s.d.

\newpage

%% ===== Page 2 =====

Pages 650--787 + table des matières

NB. La page 650 figure deux fois (erreur de numérotation)

\newpage

%% ===== Page 3 =====

\noindent \textbf{50. Enveloppes schématiques des groupes discrets : cas métabélien} \hfill 650

\begin{enumerate}
\item[1)] Sections \hfill 655
\item[2)] Essai de caractérisation des sections pro-unipotentes \hfill 662
\item[3)] Schéma en groupes associé à certaines extensions de groupes discrets (noyau ab.) \hfill 665
\item[4)] Enveloppe schématique d'un groupe métabélien \hfill 669
\item[5)] Caract. universelle de l'env. schématique \hfill 669
\end{enumerate}

\vspace{1em}
\noindent \textbf{51. Exemple : la tour de Fermat.} \hfill 669

\begin{enumerate}
\item[1)] Les groupes $\mathcal{G} = \\operatorname{Sl}(2, \mathbf{Z}) / [\pm I_2]$ et $\overline{\mathcal{G}}$ \hfill 665
\item[2)] $\mathcal{G}_{p,\sigma}$ : version provisoire \hfill 673
\item[3)] Les groupes $\widehat{\Phi}_N^*$ et variantes
\item[4)] La suite $1 \to \mathcal{N}_{p,\sigma}^{0+} \to \mathcal{N}_{p,\sigma} \to \widehat{\Phi}_{12}^* \to 1$, et relations avec $\\operatorname{Sl}(2, \mathbf{Z}/12\mathbf{Z})$ \hfill 677
\item[5)] $\mathcal{M}_{p,\sigma}$, $\underline{\mathcal{M}}_{p,\sigma}$ via $\mathcal{N}_{p,\sigma}$ \hfill 682
\item[6)] La suite $1 \to \mathcal{G}_\sigma \to \mathcal{G}_{p,\sigma} \to \widehat{\Phi}_{12}^* \to 1$, et $1 \to \overline{\mathcal{G}}_\sigma^+ \to \overline{\mathcal{G}}_{p,\sigma} \to \widehat{\Phi}_{12}^* \to 1$ \hfill 684
\item[7)] $\mathcal{N}_p, \mathcal{N}_p^\Sigma$ etc \hfill 686
\item[8)] $\mathcal{N}_\sigma, \mathcal{N}_\sigma^\Sigma$ etc \hfill 689
\item[9)] Retour sur $\underline{\mathcal{M}}_{p,\sigma}, \mathcal{M}_{p,\sigma}$ etc \hfill 690
\item[10)] $\mathcal{M}_{p,\sigma}(0), \mathcal{M}_{p,\sigma}(1/2), \mathcal{M}_{p,\sigma}(-j)$ etc \hfill 691
\item[11)] L'hom. $\mathcal{M} \to \mathcal{M}_{p,\sigma} = \underline{\mathcal{M}}_{p,\sigma}(\hat{\mathbf{Z}})^{\wedge}$, et relation avec la tour des courbes de Fermat : position du problème. \hfill 695
\item[12)] Scindages de l'extension $\mathcal{G}_{p,\sigma} / \mu_6$ de $\hat{\mathbf{Z}}^* \times \mathfrak{S}_3$ par $\hat{\mathcal{M}}$. \hfill 700
\item[13)] Scindages sur la tour de Fermat ; relations de compatibilité \hfill 707
\item[14)] Relations avec jacobiennes généralisées. \hfill 710
\item[15)] Interprétation de $\boldsymbol{\Gamma}_{\mathbf{Q}} \to$ [unclear: $\mathcal{N}_{\mathfrak{S}} \simeq I \times \hat{\mathbf{Z}}$] + digression sur extensions kummériennes. \hfill 725
\item[16)] Réflexions sur le th. de Fermat \hfill 733
\end{enumerate}

\newpage

%% ===== Page 4 =====

\chapter*{52. Approximations de Lie des systèmes de Galois-Teichmüller}
\addcontentsline{toc}{section}{52. Approximations de Lie des systèmes de Galois-Teichmüller}

\begin{enumerate}
    \item[1º)] Notion de $L$ et $TL$-groupe : cas discret \hfill 737
    \item[2º)] Cas profini \hfill 744
    \item[3º)] Approximations de Lie des systèmes de Galois-Teichmüller \hfill 748
    \item[4º)] Approximation en Lie induite par une approximation en Lie géo-arithm. \hfill 758
\end{enumerate}

\vspace{1em}

\section*{53. Le groupe de Teichmüller spécial $\widehat{\mathfrak{S}}_{1,1} \simeq \widehat{\operatorname{Gl}}(2, \mathbf{Z})^\sim$, et le groupe $S\mathcal{N}_{1,1}$ \hfill 760}
\addcontentsline{toc}{section}{53. Le groupe de Teichmüller spécial $\widehat{\mathfrak{S}}_{1,1} \simeq \widehat{\operatorname{Gl}}(2, \mathbf{Z})^\sim$, et le groupe $S\mathcal{N}_{1,1}$}

\begin{enumerate}
    \item[1º)] Action différentielle du gr. de Teichm. spécial \hfill id \\
    (variante provisoire)
    \item[2º)] Les groupes $\operatorname{Gl}(2, R)^{\pm 1 \sim}$, $\operatorname{Sl}(2, R)^\sim$ \hfill 761
    \item[3º)] Le groupe $\widehat{\operatorname{Gl}}(2, \mathbf{Z})^\sim$ \hfill 763
    \item[4º)] L'isomorphisme $\widehat{\operatorname{Gl}}(2, \mathbf{Z})^\sim \simeq \widehat{\mathfrak{S}}_{1,1}$ \hfill 765
    \item[5º)] Digression sur les revêtements de [unclear: groupes] d'espaces, et relations entre [unclear: groupe] top. et [unclear: groupes] de fibrations de repères \hfill 768
\end{enumerate}

\newpage

%% ===== Page 5 =====

\section*{50. Enveloppes schématiques de groupes discrets. \\ (Cas ind-abélien)}

\marginpar{Sections unipotentes}1. Soit $G$ un schéma en groupes localement de présentation finie sur un schéma de base $S$ (le cas important pour nous sera $S = \operatorname{Spec}(\mathbf{Z})$). Une section $g$ de $G$ sur $S$ est dite \emph{unipotente} si sur toute fibre géométrique $G_{\bar{s}}$, l'élément $g_{\bar{s}} \in G_{\bar{s}}(\bar{k})$ ($\bar{k} = k(\bar{s})$) est unipotent, i.e.\ s'il existe un homomorphisme de groupes
$$ \mathbf{G}_{a,k} \xrightarrow{\varphi_{\bar{s}}} G_k \quad k = k(\bar{s}) \leqno{(1)} $$
tel que
$$ g_{\bar{s}} = \varphi_{\bar{s}}(1) \leqno{(2)} $$

On dit que $g$ est \emph{quasi-unipotente} si pour tout $s \in S$, il existe $n \in \mathbf{N}^*$ tel que $g_{\bar{s}}^n$ soit unipotente. Si la caractéristique de $\bar{s}$ est $p>0$, cela signifie aussi que $g_{\bar{s}}$ est d'ordre fini. On voit que si $S$ est quasi-compact, alors $g$ est quasi-unipotente ssi $\exists n \in \mathbf{N}^*$ tel que $g^n$ soit unipotente.

Pour que $g$ soit unipotente, il suffit qu'il existe un homomorphisme
$$ \varphi : \mathbf{G}_a \longrightarrow G \quad \text{t.q.\ } g = \varphi(1) \leqno{(3)} $$

\newpage

%% ===== Page 6 =====

tel que
$$ \varphi(1) = g \leqno{(4)} $$

Je me demande si c'est aussi nécessaire,
dans des "bons" cas. Prenons déjà le
cas où $S = \operatorname{Spec} K$. En caractéristique $0$,
on a un résultat d'unicité pour $\varphi$,
qui permet de [unclear: déduire] que le [unclear: schéma] des
[unclear: éléments] [unclear: cherchés] est bien [unclear: représentable].
Si après extension $S' \to S$ fid.\ plate quasi-comp.\
il existe $\varphi'$ ayant les propriétés voulues,
alors $\varphi$ existe (sur $S$ lui-m.). Sur un corps de car.\ $0$ c'est OK. En car.\ $p>0$
la difficulté provient du fait qu'il y a
des hom.\ de groupes $\mathbf{G}_a \to \mathbf{G}_a$ (tels $x \mapsto x^p - x$)
qui ne sont [unclear: pas] triviaux, et qui s'annulent
sur $1$. On s'en tire...

- [unclear: Obligation] de [unclear: demander] [unclear: l'unicité] : [unclear: l'existence] de $\varphi$ qui
soit [unclear: unipotent]. Tout sous-groupe de $G$
[unclear: Remarque] : $\mathbf{G}_a$ a-t-il un [unclear: quotient]
admettant une section pointée non triviale ?
Est-ce au moins vrai sur un corps de
base ?

Je n'ai pas envie de me lancer dans
ces questions (certes intéressantes !), et préfère

\newpage

%% ===== Page 7 =====

me limiter à un cas plus proche de celui
auquel j'en suis, $S = \operatorname{Spec} \mathbf{Z}$, en supposant
$S$ noeth. intègre, de pt gén. $\eta$ de
car. $0$. Cela implique déjà, d'après le
résultat sur un corps de car. $0$, qu'il
existe un ouvert $U$ et
$$
U \subset S \quad \text{ouvert dense} \quad \text{et} \quad \varphi_U : G_{1|U} \to G_{|U} \leqno{(5)}
$$
tels que l'on ait (4) sur $U$ -- et (pour
$U$ choisi) ce $\varphi_U$ est unique. Par l'unicité,
on voit donc qu'il y a un plus gr. $U$
sur lequel on peut trouver un
tel $\varphi_U$, et la question ouverte est : a-t-on
nécessairement $U=S$.

Si on avait $U \neq S$, quitte à localiser
en un pt maximal de $S \setminus U$, on peut se
ramener à examiner le cas où $S$ est local,
de point fermé $s$ -- et à prouver que
dans ce cas, on peut pro-longer $\varphi_U$ sur $S$
(d'où une contradiction\dots). Dans
le cas qui nous intéresse, $S$ est régulier de
dimension 1, donc un trait. Je vais me borner à ce cas, crucial.

\newpage

%% ===== Page 8 =====

sans doute rien pour élucider le cas général,
et qui suffira pour nos propos.

Soit $y$ le point générique de $S$, $s$ un point fermé.
Considérons le sous-groupe $L_y ( \simeq (G_s)_y ) \subset G_y$,
soit $L$ son adhérence schématique dans $G$,
qui est un sous-schéma en groupes réduit
de $G$. Soit $L^\circ = L \setminus (\bigcup L_{s,i})$\footnote{composantes de $L_s$ ne contenant pas $1$}
sa composante neutre. On a $L^\circ \simeq G_S$. En fait, est-il
pourvu que $L^\circ_s \simeq G_s$. En fait, est-il
vrai que $(L^\circ)_s (= (L_s)^\circ)$ soit réduit ?
Mais pour que $\varphi : G_y \to G_y$ ait un
prolongement $\varphi_S$, il n'est pas nécessaire
sans doute que l'on ait $L^\circ \simeq G_S$.
Car il est possible que $\varphi$ ne soit pas
monomorphisme, i.e.\ que $\varphi_s$ ne soit pas monomorphisme - 
il pourrait même être l'homomorphisme trivial (si $\varphi_s = 1$) !
Mais même si $\varphi_s \neq 1$, il n'est pas clair que
$\varphi_s$ soit monomorphisme.

Prenons l'exemple où
$$ G = \underline{\Aut}(M) \simeq \operatorname{Gl}(n)_S \leqno{(6)} $$
$M$ un module libre de rang $n$ sur $S$,
alors on aura (si $\operatorname{car} y = 0$)
$$ \varphi_y(\lambda) = \sum \frac{X^i}{i!} \lambda^i \leqno{(7)} $$
$$ X \in \operatorname{End}(M_y) \qquad K = k(y) $$

\newpage

%% ===== Page 9 =====

est un endomorphisme nilpotent de $M_K$.
donc un élt de $\operatorname{End}_A(M) \otimes_A K$, $X^n=0$.
$$
X \in \operatorname{End}_A(M) \otimes_A K \quad , \quad X^n = 0 \leqno{(8)}
$$
Donc par hypothèse
$$
g_y = \varphi_y(1) = \sum \frac{X^i}{i!} \in \\operatorname{Aut}(M)
$$
i.e.
$$
\sum \frac{X^i}{i!} \quad , \quad \sum \frac{(-1)^i X^i}{i!} \in \operatorname{End}_A(M) \leqno{(9)}
$$
et la question si $g_y$ provient d'un hom. de groupes $\mathbb{G}_a \to G$ revient à celle de savoir si
$$
\frac{X^i}{i!} \in \operatorname{End}_A(M) \qquad \forall 1 \le i \le n-1 \leqno{(10)}
$$
(pour $i=0$ et $i \ge n$ c'est trivial)

Pour $n=2$ c'est trivial, Pour $n=3$ c'est trivial
(dans (9) il y a un seul terme avec un dénominateur).
$$
\varphi_y(\lambda) = 1 + X\lambda + \frac{X^2}{2}\lambda^2 \qquad X \in M_q(K)
$$
$$
\varphi_y(\lambda) = 1 + X\lambda + \frac{X^2}{2}\lambda^2 + \frac{X^3}{6}\lambda^3 \qquad \begin{matrix} X \in M_q(K) \\ X^4=0 \end{matrix}
$$
par hyp. $\varphi_y(1)$ et $\varphi_y(-1)$ dans $M_q(A)$, i.e.
$$
\frac{X^2}{2} \pm \frac{X^3}{6} \in M_q(A) \leqno{(11)}
$$
on demande s'il en résulte
$$
\frac{X^2}{2} , \frac{X^3}{6} \in M_q(A) \text{ ?} \leqno{(12)}
$$
Si $2$ inv. dans $A$ c'est clair, on en tire
de (11) que $X^2 \in M_q(A)$ d'où $\frac{X^2}{2} \in M_q(A)$,
donc $X^3/6 \in M_q(A)$, OK. Donc la difficulté-

\newpage

%% ===== Page 10 =====

se présente ici en \marginpar{Essai de caractérisation des sections quasi-unipotentes} [unclear: car.] résiduelle 2. Supposons pour simplifier que 2 soit [unclear: régulier] dans $A$, on a, si $X \notin M_4(A)$
$$
X = \frac{1}{2^\nu} X_0 \dots X_0 \in M_4(A), \quad X_0 \notin 2 M_4(A)
$$
de sorte que la donnée de $X$ équivaut à celle de
$X_0 = (a_{ij})_{1 \le i, j \le 4}$, tels que les $a_{ij} \in A$ pas tous
divisibles par 2, etc.
$$
X_0^4=0, \quad \frac{1}{2} \frac{X_0^2}{2^{2\nu}} \pm \frac{1}{6} \frac{X_0^3}{2^{3\nu}} \in M_4(A) \leqno{(11 \text{ bis})}
$$
et la question est s'il en résulte que l'on ait
$$
\frac{1}{2} \frac{X_0^2}{2^{2\nu}} , \frac{1}{6} \frac{X_0^3}{2^{3\nu}} \in M_4(A) \quad ? \leqno{(12 \text{ bis})}
$$

Bien sûr, la réponse est trivialement oui
si $X^3=0$, s'écrivant encore
$X_A^3 = 0$, i.e. $g_y$ est "unipotent principal".

Pour en avoir le cœur net, il faudrait
un contre-exemple dans le contexte précédent
(avec une matrice nilpotente d'ordre 4) --
ou un théorème ! Mais pour mon
propos, il semble qu'il ne soit pas
utile d'avoir une solution de ce pb.

Par contre, je vais essayer le

\underline{Théorème}. Soit $G$ un schéma en groupes (séparé)
loc. de présentation finie sur un schéma
$S$ quasi-compact, et soit $g$ une section de
$G$. Pour que $g$ soit quasi-unipotente, il

\newpage

%% ===== Page 11 =====

faut et il suffit qu'il existe un entier $n > 0$, et un hom.\ de schémas en groupes $\varphi \colon \mathbf{G}_a \to G$, tel qu'on ait
$$
\varphi(1) = g^n \leqno (13)
$$

La suffisance étant triviale, il s'agit d'examiner la nécessité. Pour simplifier, on va même supposer $S$ noethérien et intègre -- une fois $S$ supposé noethérien, je pense que la réduction au cas intègre doit se faire par des dévissages faciles, que je n'ai pas envie de faire, faute de besoin essentiel ! Je vais donc examiner le cas $S$ noethérien intègre, de point générique $y$.

a) Cas où car $y$ est $> 0$. Alors l'hypothèse sur $g_y$ est que $g_y$ est d'ordre fini -- donc $\exists n > 0$ tel que $g_y^n = 1$, donc $g^n = 1$. On prendra dans (13) cet $n$, et $\varphi$ l'hom.\ trivial.

b) Cas où $y$ de car. nulle. Soit $\varphi_y$
$$
\varphi_y \colon \mathbf{G}_{a y} \to G_y \leqno (14)
$$
défini par la condition
$$
g_y = \varphi_y(1)
$$
Soit $U$ un sous-ouvert de définition dans $S$.

Soit
$$
\varphi_{n y}(\lambda) = \varphi_y(n \lambda) \colon \mathbf{G}_{a y} \to G_y
$$
et soit $U_n$ son sous-ouvert de définition. On a
$$
\varphi_{n y}(1) = g^n \quad \text{sur } U_n \leqno (15)
$$

\newpage

%% ===== Page 12 =====

et
$$
n | n' \implies U_n \subset U_{n'} \leqno{(16)}
$$

Il faut prouver que la réunion des $U_n$
est $S$. On se ramène au cas $S$ local,
et où $\varphi_1$ est déjà défini dans
$$
U_1 = S \setminus \{s\} \leqno{(17)}
$$
, $s$ le pt fermé de $S$
et cela revient au

Lemme (?) Soient $S$ intègre, $G$ schéma en groupes
loc. de prés. finie sur $S$, $U$ ouvert non-vide de $S$,
$\varphi_U : \mathbb{G}_{mU} \to G_U$ un hom. de groupes, on
suppose que $\varphi_U(1) = g_U$ se prolonge en une
section $g$ de $G$ sur $S$. [unclear: À isogénie] (strict. unipo-
tente, d'après des résultats connus généraux de
SGA 3) Alors il existe un entier $n>0$ tel que
$\varphi_U^{(n)} : \lambda \mapsto \varphi_U(\lambda^n)$ se prolonge en un hom.
$\varphi : \mathbb{G}_m \to G$ .

Je ne vais pas essayer de le démontrer dans un
contexte général des schémas en groupes généraux,
et vais me limiter, à titre de test décisif
(pour emporter conviction que ça marche)
au cas où $G = \\operatorname{Gl}(n)_S$, comme tantôt,
et celui où $S$ de plus est régulier de dim 1.
Donc on est ramené à la situation d'un anneau
$A$, et de $X \in M_n(K)$ ($K$ corps des fractions)
satisfaisant

\newpage

%% ===== Page 13 =====

$$
\begin{cases}
X \in M_n(K) \\
X^n = 0,\ \sum \frac{X^i}{i!} \in M_n(A),\ \sum (-1)^i \frac{X^i}{i!} \in M_n(A)
\end{cases}
\leqno{(18)}
$$
et on veut montrer qu'il existe un entier $d>0$ tel qu'on ait
$$
\frac{d^i X^i}{i!} \in M_n(A) \quad (1 \le i \le n-1)
\leqno{(19)}
$$

Si $A$ est de car. $0$ i.e.\ s'il est une $\mathbf{Q}$-algèbre, ceci implique d'ailleurs :\footnote{Le cas de car.\ nulle se résout par l'affirmative à l'aide de la fonction logarithme : on pose $X + \frac{X^2}{2!} + \dots = Y$ (i.e.\ $\exp(X)=1+Y$), d'où $X = \log(1+Y) = Y - \frac{Y^2}{2} + \dots$ Montre même [unclear: la base] sur les algèbres \dots}
$X \in M_n(A)$.

Il se pourrait que dans le cas $S$ [unclear: envisagé] ce soit $\varphi_U$ lui-même qui puisse se définir sur $S$ tout entier. Le premier cas douteux semble être $n=5$, $X \in M_5(A), X^5=0$
$$
\begin{cases}
\frac{X^2}{2} \pm \frac{X^4}{4!} \in M_5(A) \\
\overset{?}{\implies} X \in M_5(A)
\end{cases}
\leqno{(20)}
$$

Par contre, dans le cas envisagé au numéro précédent ($n=4$, car.\ résiduelle $2$ - le seul qui pose une difficulté) il est clair qu'il suffit de prendre pour $d$ une puissance de $2$ assez grande.

Cette remarque s'applique au cas d'inégales caractéristiques. On devrait avoir en effet [unclear: ceci].

Corollaire au lemme (dubitatif) : si $S$ est local, de car.\ résiduelle $p>0$, et si $\varphi$ est définie sur un ouvert qui contient $S - V(p)$, (i.e.\ si le complémentaire de $U$ est de car.\ $p>0$ - fait inutile de supposer d'ailleurs que $S$ soit local) alors $\exists \ q=p^k$ tel que $\varphi^q$ se prolonge.

\newpage

%% ===== Page 14 =====

Quand $G$ est affine, ceci devrait être trivial. On se ramène à

Proposition. Soit $S$ un préschéma [unclear: fin. qc?], $n$ un entier $\ge 1$, $S_0 = V(n) \subset S$ le sous-schéma fermé défini par $n.1_S$ (donc [unclear: crossed out text]), $U = S \setminus S_0$, $G$ un schéma en groupes affine sur $S$, et

(21) $$ \varphi_U : G_{a|U} \to G|_U $$

un hom. de groupes. On suppose $U$ schématiquement dense dans $S$. Alors il existe une puiss. $m = n^r$ de $n$, et un hom. de schémas en groupes

$$ \psi_m : G_a \to G , $$

tels qu'on ait

(22) $$ \psi_m|_U = (\varphi_U^m : \lambda \mapsto \varphi_U(m\lambda) = \varphi_U(\lambda)^m) $$

Considérons en effet l'endo.

$$ u_m : G_a \to G_a \qquad \lambda \mapsto m\lambda $$

Alors la formule (22) signifie qu'on a commutativité

\begin{tikzcd}
G_{a|U} \ar[r, "(u_m)_U"] \ar[dr, "\psi_m|_U"'] & G_{a|U} \ar[d, "\varphi_U"] \\
& G_U
\end{tikzcd}

et l'existence de $\psi_m$ pour un $m$ grand, signifie que pour un $m$ grand le hom. composé

$$ \varphi_U \circ u_m : G_{a|U} \to G_U $$

se prolonge en un hom de groupes $G_a \to G$.

\newpage

%% ===== Page 15 =====

Pour raisons de densité schématique et [unclear: d'unicité], il suffit de prouver que $\varphi_U$ se prolonge en un hom. de schémas en groupes.

OPS $S = \operatorname{Spec} A$, $A$ un anneau, et soit $B$ celui de $G$. Donc l'anneau de $U$ sera $A_n = A[1/n]$, celui de $G$ sera $B_n = B_{A_n}$, par hyp on a un hom.
(correspondant à $\varphi_U$)
$$ B_n \xleftarrow{\varphi} A_n[T] \leqno{(23)} $$
tandis que l'hom. correspondant à $\psi_{m U}$ sera [unclear: donné par des polynômes]
$$ A_n[T] \xleftarrow{\psi_m^*} A_n[T] \leqno{(24)} $$
$$ T \mapsto m T \quad \text{donc } T^i \mapsto m^i T^i $$

Soit $x = a_0 + a_1 T + \dots + a_N T^N \in A_n[T]$ avec les $a_i \in A_n$
on a 
$$ \psi_m^*(x) = \sum a_i m^i T^i $$
donc $\psi_m^*(x) \equiv a_0 \pmod{A[T]}$

ce qui [unclear: montre]
$$ \psi_m^*(x) \in A[T] \text{ pour } m = n^r \text{ grand } \iff a_0 \in A \leqno{(25)} $$

\newpage

%% ===== Page 16 =====

Appliquant ceci aux s-gén d'éléments de $B$ qui engendrent $B$ comme $A$-algèbre,...

Corollaire. $S$, $U$ étant comme dans la proposition, et $G$ étant un schéma affine de t.f. sur $S$ (il n'est plus question de structure de groupes), et $g_U : E_U^1 \to G_U$ un hom. de schémas, pour qu'il existe un morphisme de schémas $g : E^1 \to G$ prolongeant $g_U$ tel que le morphisme donné par (24) se prolonge, il faut et il s. que $g_U^*$ se prolonge en une section de $G$ sur $S$.

Remarque. Les arguments précédents fournissent en tous cas un "théorème" de densité p. 662 ! dans le cas où $S$ est noeth. de dim. 1 (le cas intéressant pour nous étant celui où le pt générique $s$ de $S$ est de car. 0).

\marginpar{Schémas en groupes associés à des groupes discrets}
Soit maintenant $M$ un $\mathbf{Z}$-module libre de type fini, on va construire un pro-schéma en groupes sur $S$, associé au syst. projectif des sous-groupes d'indice fini de $M$. Pour faire la construction, je vais [unclear: développer]...

\newpage

%% ===== Page 17 =====

656

[MARGIN: Y a une complication sur l'existence d'une construction de [unclear] $M_{S_0}$ posera $\underline{G} = \underline{G}_{S_0} \times M_{S_0}$]

Lemme préliminaire, relatif à une extension

(26) $$1 \to M \to G \to H \to 1$$

suite exacte de groupes discrets, avec $M$ commutatif — le cas qui nous intéressera pour la suite étant celui où $M$ est un $\mathbb{Z}$-module libre de t. f. En tous cas $M$ est un $\mathbb{Z}$-module, et à ce titre définit un Groupe sur $\operatorname{Spec}\ \mathbb{Z}$, savoir 
$$A \longmapsto M \otimes_{\mathbb{Z}} A$$
($A$, $\mathbb{Z}$-algèbre commutative) — qu'on désigne par $\underline{M}$. En fait, $\underline{M}$ est représentable, i.e. un schéma en groupes sur $\mathbb{Z}$ si $M$ est libre de type fini. J'ai envie de définir aussi des schémas en Groupes sur $S_0$ (qui seront des schémas en groupes si $M$ l'est) $\underline{G}, \underline{H}$, et une suite exacte

(27) $$1 \to \underline{M} \to \underline{G} \to \underline{H} \to 1$$

de Groupes sur $S_0$, et un isomorphisme de suites exactes

(28)
\begin{tikzcd}
1 \ar[r] & M \ar[r] \ar[d, equal] & G \ar[r] \ar[d, "\simeq"'] & H \ar[r] \ar[d, "\simeq"'] & 1 \\
1 \ar[r] & \underline{M}(\mathbb{Z}) \ar[r] & \underline{G}(\mathbb{Z}) \ar[r] & \underline{H}(\mathbb{Z}) \ar[r] & 1
\end{tikzcd}

On prendra

(29) $$\underline{H} = H_{S_0} \text{ Groupe constant de val. } H$$

et les flèches verticales extrêmes de (28)

\newpage

%% ===== Page 18 =====

soient des isomorphismes canoniques.

On peut songer à définir $\underline{\mathcal{G}}$ comme
foncteur (cov.) [unclear: $A/\mathbf{Z}$ en (Anneaux)] ou $S/S_0$,
par
$$ \underline{\mathcal{G}}(A) = \mathcal{G} \amalg_M M(A) \leqno{(30)} $$
$\underline{\mathcal{G}}(A)$ est alors défini comme l'extension de
$H$ par $M(A)$ déduite de l'extension $\mathcal{G}$
de $H$ par $M$ par l'homomorphisme
$$ M \to M \otimes_{\mathbf{Z}} A , \leqno{(31)} $$
qu'on considère comme un homomorphisme de
$H-\mathbf{Z}-\text{modules}$, en faisant opérer $H$
sur $M \otimes_{\mathbf{Z}} A$ via son opération sur $M$.
La petite difficulté est que le foncteur
en groupes ainsi défini par (30) n'est pas
représentable - il ne donne l'
expression correcte de $\underline{\mathcal{G}}(A)$ que
pour $A$ connexe - il correspond
à un sous-foncteur de
$$ S \mapsto \underline{\mathcal{G}}(A) = \underline{\mathcal{G}}(S) \quad (S = \operatorname{Spec} A) $$
formé des $g \in \underline{\mathcal{G}}(S)$ tels que la projection
$$ S \to \underline{H} \leqno{(32)} $$
(qui correspond à une fonction localement constante
sur $S$ à valeurs dans $H$) soit constante.

\newpage

%% ===== Page 19 =====

L'expression correcte s'obtient en partant
du pré-foncteur "incorrect"
$$ [\underline{\mathfrak{S}}] : A \longmapsto \mathfrak{S}_M (M \otimes_{\mathbf{Z}} A) $$
ou mieux, du contrafoncteur
$$ S \longmapsto \mathfrak{S}_M \Gamma(S, M \otimes_{\mathbf{Z}} \mathcal{O}_S) \quad , $$
et en prenant le faisceau Zariski-associé -
- en fait on aura
$$ \underline{\mathfrak{S}}(S) = \varinjlim_{\mathcal{U}=(U_i)_{i \in I}} \prod_{i \in I} [\underline{\mathfrak{S}}](U_i) \leqno{(33)} $$
(où $\mathcal{U} = (U_i)_{i \in I}$ est une partition
de $S$ en schémas,
$S \simeq \coprod U_i$).

Voici une autre façon de présenter la
construction de $\underline{\mathfrak{S}}$, en utilisant la
construction générale du § 49 -- dans le
contexte faisceautique. On a un système
$$ \underset{\text{Gr. sur } S_0}{M} \quad , \quad \underset{\text{Gr. préschémas sur } S_0, \mathcal{M}_{S_0} \text{ distingué}}{\underline{\mathfrak{S}}_0 \supset \mathcal{M}_{S_0}} \quad , \quad \underset{\text{homom: } \mathcal{M}_{S_0} \to M}{\text{Op. de } \underline{\mathfrak{S}}_0 \text{ sur } M} \leqno{(34)} $$
(où l'op. de $\underline{\mathfrak{S}}_0$ sur $M$, qui \unclear{munit un gr. de} $\underline{\mathfrak{S}}$ sur $M$, est celle déduite
de l'op. de $\underline{\mathfrak{S}}$ sur $M$, par \unclear{fonctorialité}).
On doit satisfaire à des conditions de
compatibilité a) b) de loc. cit. Par la
construction de loc. cit, on construit

\newpage

%% ===== Page 20 =====

à l'aide de ces données un système

$$ \underline{\mathfrak{S}} \supset \underline{M} \quad , \quad \text{et } \Phi : \underline{\mathfrak{S}}_0 \to \underline{\mathfrak{S}} \leqno{(35)} $$
Groupes sur $S_0$,
avec $\underline{M}$
sous-groupe distingué

tel que (34) se reconstitue à partir de (35) en prenant $\underline{M}, \underline{\mathfrak{S}}_0$ dans (35), $\underline{M}_{S_0} \defeq \text{image inverse de } \underline{M} \subset \underline{\mathfrak{S}} \text{ par } \Phi$, $\varphi$ induit par $\Phi$, et $\Theta$ l'opération déduite de $\Phi$. De plus, $\Phi$ induit un iso.
$$ \underbrace{\underline{\mathfrak{S}}_0 / \underline{M}_{S_0}}_{H_{S_0}} \isommap \underline{\mathfrak{S}} / \underline{M} $$

-- et on a là une description "axiomatique" du système (35) en termes de (34).
La suite exacte de Groupes sur $S_0$

$$ 1 \to \underline{M} \to \underline{\mathfrak{S}} \to \underset{\begin{subarray}{c} \simeq \\ H_{S_0} \end{subarray}}{\underline{\mathfrak{S}}/\underline{M}} \to 1 \leqno{(36)} $$

[unclear: Il s'ensuit pour tout] schéma de base affine $S$ (tel que $H^1(S, M \otimes_{\mathbf{Z}} \mathcal{O}_S) = 0$)\footnote{L'hypothèse (pour des schémas affines de fait $H^1$ est nul) est en fait inutile, car (plus généralement) les $\underline{M}$-torseurs sur $S$ définis par des images inverses évidents de $\underline{M}_{S_0}$-torseurs triviaux, sont exacts.}
$$ 1 \to \underset{\begin{subarray}{c} \simeq \\ M \otimes_{\mathbf{Z}} A \\ \text{si } S = \text{Spec } A \end{subarray}}{\underline{M}(S)} \to \underline{\mathfrak{S}}(S) \to \underset{\begin{subarray}{c} \simeq \\ H(S) \end{subarray}}{H_{S_0}(S)} \to 1 $$

\newpage

%% ===== Page 21 =====

En fait, on peut [unclear: construire] un diagr.
commut. de suites exactes

(38)
\begin{tikzcd}
1 \ar[r] & \underline{M}_{S_0} \ar[r] \ar[d] & \underline{\mathfrak{S}}_{S_0} \ar[r] \ar[d] & \underline{H}_{S_0} \ar[r] \ar[d] & 1 \\
1 \ar[r] & \underline{M} \ar[r] & \underline{\mathfrak{S}} \ar[r] & \underline{H}_{S_0} \ar[r] & 1
\end{tikzcd}

qui par passage aux sections donne
pour $S$ quelconque affine (disons) connexe
un diagr. comm.

[MARGIN: NB En [unclear: déduisant] de (39)
$\underline{\mathfrak{S}}(A)$ est isomorphe
au produit semi-direct
de $H$ par
$M \otimes_{\mathbb{Z}} A$
[unclear]
[unclear]]

(39)
\begin{tikzcd}
1 \ar[r] & M \ar[r] \ar[d] & \mathfrak{S} \ar[r] \ar[d] & H \ar[r] \ar[d, equal] & 1 \\
1 \ar[r] & M \otimes_\mathbb{Z} A \ar[r] & \underline{\mathfrak{S}}(A) \ar[r] & H \ar[r] & 1
\end{tikzcd}

pour $S=S_0$ on retrouve l'iso (28)

- . - . - .

Cette construction du $S_0$-schéma en
groupes $\underline{\mathfrak{S}}$, à partir d'une
suite exacte (27), est [unclear: manifestement]
fonctorielle par r. à cette dernière : un
hom. de suites exactes

(40)
\begin{tikzcd}
1 \ar[r] & M \ar[r] \ar[d] & \mathfrak{S} \ar[r] \ar[d] & H \ar[r] \ar[d] & 1 \\
1 \ar[r] & M' \ar[r] & \mathfrak{S}' \ar[r] & H' \ar[r] & 1
\end{tikzcd}

donne naissance à un diagr. comm.
de suites exactes

\newpage

%% ===== Page 22 =====

(41)
\begin{tikzcd}[row sep=1.5em, column sep=2em]
& 1 \ar[r] & M_{S_0} \ar[r] \ar[dd] \ar[dl] & \underline{G}_{S_0} \ar[r] \ar[dd] \ar[dl] & H_{S_0} \ar[r] \ar[dd] \ar[dl] & 1 \\
1 \ar[r] & M'_{S_0} \ar[r] \ar[dd] & \underline{G}'_{S_0} \ar[r] \ar[dd] & H'_{S_0} \ar[r] \ar[dd] & 1 & \\
& 1 \ar[r] & \underline{M} \ar[r] \ar[dl] & \underline{\underline{G}} \ar[r] \ar[dl] & H_{S_0} \ar[r] \ar[dl] & 1 \\
1 \ar[r] & \underline{M}' \ar[r] & \underline{\underline{G}}' \ar[r] & H' \ar[r] & 1 &
\end{tikzcd}

[MARGIN: Exemple 4. réduction discrète des groupes métabéliens]

Reprenons donc un $\mathbb{Z}$-module libre
de type fini $M$, et pour $H$ un
groupe d'ordre fini, considérons
la suite exacte

(42) $$1 \longrightarrow M_i \longrightarrow M \longrightarrow M/M_i \longrightarrow 1$$

Plus gén., si on part d'une suite
exacte (27), pour un sous-groupe $M_i$
d'indice fini de $M$, stable par
$\underline{\underline{G}}$ i.e. stable par $H$, on a une
suite exacte

(43) $$1 \longrightarrow M_i \longrightarrow \underline{\underline{G}} \longrightarrow \underline{\underline{G}}/M_i \longrightarrow 1$$

qui montre que $H_i = \underline{\underline{G}}/M_i$
est l'extension de

(44) $$1 \longrightarrow M/M_i \longrightarrow H_i \longrightarrow H \longrightarrow 1$$

déduite de (27) par
$M \longrightarrow M/M_i$.

\newpage

%% ===== Page 23 =====

On note que si $H$ est [unclear: de type fini] et $M$ [unclear: de gén. finie pour $H$] ; alors les $M_i \subset M$ d'indice fini dans $M$ stables sous $H$ sont \marginpar{prendre les $nM$ !} cofinaux dans l'ensemble de tous les $M_i$ d'indice fini dans $M$.

Considérons donc la suite exacte de schémas en groupes associés
$$
1 \to \underline{M}_i \to \underline{\mathcal{G}}_i \to \underline{H}_i \to 1 \leqno (45)
$$
Pour $i$ variable, elles forment un système projectif d'extensions, dans une catégorie de schémas en groupes sur [unclear: $\operatorname{Spec} \mathbf{Z}$].

On considère le pro-objet $(\underline{\mathcal{G}}_i)$ comme pro-objet de la catégorie des schémas en groupes (lisses etc..) sur $S_0$. On note que les morphismes de transition
$$
\underline{\mathcal{G}}_j \to \underline{\mathcal{G}}_i
$$
sont affines - ce qui [unclear: revient à la nullité] pour l'homomorphisme
$$
\underset{\substack{M_i \supset M_j \\ \text{ext. de } (M_i/M_j)_{S_0} \text{ par } \underline{M}_j}}{E(M_i, M_j)} \longrightarrow \underset{\underline{M}_i}{E(M_i, M_i)}
$$

\newpage

%% ===== Page 24 =====

Or le premier est affine, comme somme finie de schémas affines, le deuxième est affine — donc le morphisme est affine.

Ceci étant, la limite projective des $\underline{\mathfrak{S}}_i$ dans la catégorie des schémas existe (plus précisément, le foncteur limite projective des $\underline{\mathfrak{S}}_i$ est représentable) — il n'y a peut-être pas [unclear: lieu de prouver] : la limite existant, et on travaillera avec un schéma en groupes — pas de type fini.

Remarque. Le pro-groupe $(\underline{\mathfrak{S}}_i)$ sur $S_0 = \operatorname{Spec} \overline{k}$ ne dépend que du pro-groupe discret $(\mathfrak{S}_i)$, muni d'une classe de sous-groupes [unclear: profinis donnés] à commensurabilité près — l'hypothèse à faire est que parmi ces sous-groupes, il y en a qui soient des $\mathbf{Z}$-modules libres de type fini (car le prenant assez petit, il sera libre de type fini). Si $M_0$ est d'indice fini dans $\mathfrak{S}$ i.e.\ si les $H_i$ sont finis,

\newpage

%% ===== Page 25 =====

alors ce pro-objet ne dépend que de $\mathcal{G}$ lui-même -- à la condition qu'il existe un sous-groupe d'indice fini de $\mathcal{G}$ qui soit un groupe commutatif, et de type fini en tant que $\mathbf{Z}$-module. (On dit alors que $\mathcal{G}$ est "métabélien")

Ainsi pour tout groupe discret métabélien $\mathcal{G}$, on construit un pro-schéma en groupes (système projectif de schémas en groupes affines) que je vais noter
$$
Env(\mathcal{G}) = (\underline{\mathcal{G}}_i) \leqno{(46)}
$$
(l'ensemble d'indices étant formé des sous-groupes distingués de $\mathcal{G}$ qui sont sans éléments de torsion). Ce pro-groupe dépend fonctoriellement de $\mathcal{G}$. Il a des propriétés d'exactitude intéressantes et évidentes, par exemple il est multiplicatif
$$
Env(\mathcal{G} \times \mathcal{G}') \isommap Env(\mathcal{G}) \times_{S_0} Env(\mathcal{G}') \leqno{(47)}
$$
Si on a une suite exacte
$$
1 \to \mathcal{G}'' \to \mathcal{G} \to \mathcal{G}' \to 1 \leqno{(48)}
$$
de groupes métabéliens, on a une suite exacte pour les enveloppes
$$
1 \to Env(\mathcal{G}'') \to Env(\mathcal{G}) \to Env(\mathcal{G}') \to 1 \leqno{(49)}
$$

\newpage

%% ===== Page 26 =====

en un sens qu'il convient de préciser :
l'occasion - pour les schémas en groupes
associés (obtenus en passant à la limite) on
aura bel et bien un hom. [unclear: fid. plat dont]
le noyau est le s-g ...

Le cas typique, auquel le cas général
se ramène de façon [unclear: quasi-immédiate], est
celui où $\mathcal{G} = M = \mathbf{Z}$. On a pour
tout entier $n \ge 1$ la suite exacte
$$
1 \to \mathbf{Z} \xrightarrow{n} \mathbf{Z} \to \mathbf{Z}/n\mathbf{Z} \to 1 \leqno{(50)}
$$
je vais expliciter la construction
de l'extension correspondante de
$\mathbf{Z}/n\mathbf{Z}$ par $\underline{\mathbf{Z}} = \mathcal{G}_0$, soit $\underline{\mathcal{G}}_n$.
C'est un schéma affine, dont l'anneau
de coordonnées est isomorphe - de façon pas
bien canonique - à un produit de
$n$ copies de celui de $\underline{\mathbf{Z}} = \mathcal{G}_0$, savoir
$\mathbf{Z}[T]$ - on trouve un tel isomorphisme
en choisissant une section ensembliste de
$\mathbf{Z}$ sur $\mathbf{Z}/n\mathbf{Z}$, par exemple par l'une
des sections $\mathbf{Z} \cap [0, n-1]$. [unclear: La loi de]

\newpage

%% ===== Page 27 =====

transition entre les $\mathfrak{S}_n$ n'est rien de plat bien sûr. Passant à la limite projective, on trouve une extension de
$$
(\hat{\mathbf{Z}})_{S_0} \defeq \varprojlim_n (\mathbf{Z}/n\mathbf{Z})_{S_0}
$$
-- schéma en groupes profini au dessus de $S_0$ -- par la limite projective des schémas $\mathfrak{S}_n$, s'identifiant de lui-même par les lois usuelles (cf. bas de 24) avec le spectre de l'anneau limite inductive des
$$
\mathbf{Z}[T_n] \subset \mathbf{Q}[T]
$$
$$
T_n = \frac{1}{n} T
$$
qui est - ici, disons-le - la sous-$\mathbf{Z}$-algèbre de $\mathbf{Q}[T]$ formé des polynômes dont le terme constant est un entier.

On a une description analogue de la limite projective de l'enveloppe scindée $\tilde{\mathfrak{S}}_n$ de $\mathbf{Z}^n$ - c'est une extension de $(\hat{\mathbf{Z}})_{S_0}$ ([unclear: attention : l'opérer en automorphismes]) par une composante neutre, qui est le spectre du sous-anneau de $\mathbf{Q}[T_1, \ldots, T_n]$ formé des polynômes dont le terme constant est entier...

\newpage

%% ===== Page 28 =====

C'est donc un schéma en groupes sur $\Spec \mathbf{Z}$ dont toutes les fibres sont isomorphes au groupe additif, y c.\ la fibre générique, qui est $\mathbf{G}_{a, \mathbf{Q}}$.

Dans le cas d'un groupe abélien quelconque $G$, désignons par $\hat{G}$ son complété profini, par $\hat{G}_{S_0}$ le schéma en groupes profini sur $S_0$
$$
\hat{G}_{S_0} \defeq \varprojlim_i (M_i)_{S_0} \leqno{(50)}
$$
on aura encore une suite exacte canonique
$$
1 \to \underline{\operatorname{Env}}^0(\hat{G}) \to \underline{\varprojlim} \, \underline{\operatorname{Env}}(\hat{G}) \to (\hat{\mathbf{Z}})_{S_0} \to 1 \leqno{(51)}
$$
où $\underline{\operatorname{Env}}^0(\hat{G})$ désigne la composante neutre de $\varprojlim \underline{\operatorname{Env}}(\hat{G})$, i.e.
$$
\underline{\operatorname{Env}}^0(\hat{G}) \defeq \varprojlim_i \underline{M}_i \leqno{(52)}
$$
$$
= \varprojlim_n n \cdot \underline{M}_0
$$
où les $M_i$ sont les [unclear: quotients d'indice fini] de $G$ par des sous-groupes libres de type fini. [unclear: Plus haut] pour un $\mathbf{Z}$-module $M$,
$$
\underline{M} = \mathbb{V}(M) = \Spec \operatorname{Sym}_{\mathbf{Z}}(\check{M}) \leqno{(53)}
$$
désigne le fibré vectoriel sur $S_0$ [unclear: qui le définit], \dots

\newpage

%% ===== Page 29 =====

où $M_0$ est un des $M_i$, disons [unclear: lié au syst.] (cofinalité de celui-ci ds tous les $M_i$) avec $M_0$. Le choix d'une base de $M_0$ sur $\mathbf{Z}$ donne un iso
$$
\operatorname{Env}^\circ(\mathfrak{G}) \simeq \operatorname{Env}^\circ(M_0) \simeq (\operatorname{Env}^\circ(\mathbf{Z}))^n \leqno (54)
$$
$$
\simeq (\mathcal{E}^1)^n
$$
$$
\mathcal{E}^1 \simeq \operatorname{Env}^\circ(\mathbf{Z}) \simeq \operatorname{Spec}(\mathbf{Z}+T\mathbf{Q}[T]) \leqno (55)
$$
\footnote{où $(\mathbb{E}^1, \dots)$ (ou $(E^1, h)$) est le schéma multiplicatif pur}
$$
\simeq \varprojlim (E^1, n)
$$
est le schéma en gr. pro arith. déjà plus haut, qui représente le foncteur-groupe
$$
A \longmapsto \mathcal{E}^1(A) = \left\{ (x_n) \in A^{\mathbf{N}^*} \Big/ x_n = m x_{mn} \forall m, n \in \mathbf{N}^* \right\} \leqno (56)
$$

[Si $A$ est tel qu'il existe un $c \in \mathbf{N}^*$ nul dans $A$ - [unclear: plus gnalement] si toutes les car. résiduelles de $A$ divisent $c$, il s'ensuit qu'il existe [unclear: dans $A$ un] tel que dès que les car. résiduelles de $A$ sont toutes $>0$ - alors $\mathcal{E}^1(A) = \{1\}$ - si par contre $A$ est de car $0$, i.e. $\dots$ $\mathcal{E}^1(A) \simeq A$]

\newpage

%% ===== Page 30 =====

5. Caractérisation universelle de l'enveloppe pro-unipotente d'un groupe abstrait.

Soit $G$ un schéma en groupes localement de présentation finie sur une base $S$. On dit qu'un homomorphisme de schémas en groupes
$$
H \to G
$$
est unipotent (resp. quasi-unipotent) si l'image de $H$ est un sous-groupe unipotent (resp. quasi-unipotent) de $G$ — ce qui se vérifie fibre par fibre. Un homomorphisme d'un groupe discret
$$
\mathcal{G} \to \Gamma(G/S)
$$
est dit unipotent (resp. quasi-unipotent) si l'homomorphisme de schémas en groupes qui lui correspond
$$
\mathcal{G}_S \to G
$$
est unipotent (resp. quasi-unipotent).

\newpage

%% ===== Page 31 =====

Notons que d'après les propriétés des $\varprojlim$ de systèmes projectifs filtrants (à morphismes de transition affines) entre des schémas, si $\mathfrak{G}$ est [unclear: pro-unipotent], d'où un système projectif de schémas en groupes affines sur $S_0$, donc sur $S$, noté $Env_S(\mathfrak{G})$, et une limite projective
$$
L Env_S(\mathfrak{G}) = L Env(\mathfrak{G}) \times_{S_0} S \simeq \varprojlim Env_S(\mathfrak{G}) \leqno{(57)}
$$
on a, dès que $G$ est localement de présentation finie
$$
\begin{tikzcd}[row sep=small, column sep=small]
\operatorname{Hom}_{pro-gr/S}(Env_S(\mathfrak{G}), G) \arrow[r, "\sim"] \arrow[d, equal, "\text{def}"'] & \operatorname{Hom}(\varprojlim Env_S(\mathfrak{G}), G) \\
\varinjlim \operatorname{Hom}(\mathfrak{G}_i, G) &
\end{tikzcd} \leqno{(58)}
$$
Notons d'autre part qu'on a un homomorphisme canonique
$$
\mathfrak{G}_S \to L Env_S(\mathfrak{G}) \leqno{(59)}
$$
déduit par changement de base de l'homomorphisme sur la base universelle $S_0$
$$
\mathfrak{G}_{S_0} \to L Env(\mathfrak{G}) \leqno{(59\text{ bis})}
$$
provenant du système d'homomorphismes canoniques.

\newpage

%% ===== Page 32 =====

$$ \mathfrak{G}_{S_0} \longrightarrow \mathfrak{G}_S $$

Donc, via (59) on trouve
$$ \\operatorname{Hom}_{\text{S-gr}}(L\operatorname{Env}_S(\mathfrak{G}), G) \longrightarrow \\operatorname{Hom}_{\text{S-gr}}(\mathfrak{G}_S, G) \simeq \\operatorname{Hom}_{\text{gr}}(\mathfrak{G}, G(S)) \leqno{(60)} $$

Comparant (58) et (60), on trouve
$$
\begin{aligned}
\\operatorname{Hom}_{\text{pro-S-gr}}(\operatorname{Env}_S(\mathfrak{G}), G) &\simeq \\operatorname{Hom}_{\text{S-gr}}(L\operatorname{Env}_S(\mathfrak{G}), G) \\
&\longrightarrow \\operatorname{Hom}_{\text{S-gr}}(\mathfrak{G}_S, G) \simeq \\operatorname{Hom}_{\text{gr}}(\mathfrak{G}, G(S))
\end{aligned} \leqno{(61)}
$$

Ceci posé, on a le résultat suivant

Théorème Sous les conditions précédentes ($\mathfrak{G}$ groupe discret arbitraire, $G$ schémas en groupes loc. de prés. finie sur schém. de base $S$), l'application canonique (61) est injective, et a comme image [unclear: l'ensemble des hom. $\mathfrak{G} \to G(S)$ à valeurs dans un s-groupe de $G$ affine sur $S_0$].

\newpage

%% ===== Page 33 =====

\section*{5. Exemple : La "tour" de Fermat.}

1. Les groupes $\mathfrak{G}_0, \mathfrak{G}$.

Je m'intéresse au quotient de $\hat{\mathfrak{G}} = \widehat{\SL}(2,\mathbf{Z})$ obtenu en prenant le quotient par $[\Pi_0, \Pi_0]$\footnote{$\Pi_0$ engendré par les $\lambda_i = -l_i$ \\ $l_0 = \begin{pmatrix} 1 & -2 \\ 0 & 1 \end{pmatrix}$ \\ $l_1 = \begin{pmatrix} 1 & 0 \\ 2 & 1 \end{pmatrix}$} -- on trouve une extension $\mathfrak{G}$ de $\mathfrak{G}_0/\Pi_0$ par le groupe
$$
M \defeq {\Pi_0}_{\text{ab}} \isom (\hat{\mathbf{Z}}_{(0)} + \hat{\mathbf{Z}}_{(1)} + \hat{\mathbf{Z}}_{(\infty)}) / \text{diagonale} \leqno{(1)}
$$

Ainsi on a une suite exacte de groupes profinis
$$
1 \to \underset{\substack{\simeq \\ {\Pi_0}_{\text{ab}}, \circlearrowleft \mathfrak{G}_0}}{M} \to \mathfrak{G} \to \underset{\substack{\simeq \\ \\operatorname{Sl}(2, \mathbf{Z}/4\mathbf{Z}) / \langle l_0, l_1 \rangle}}{\mathfrak{G}_0/\Pi_0} \to 1 \leqno{(2)}
$$

Par générateurs et relations on a :
$$
\mathfrak{G} = \left\{ \tau', \rho, \sigma \Biggm| \begin{array}{l} \rho^6 = \sigma^4 = {\tau'}^2 = 1 \\ \rho^3 = \sigma^2,\ \tau'(\rho) = \rho^{-1},\ \tau'(\sigma) = \sigma^{-1} \\ {}[\lambda_0, \lambda_1] = 1,\ \text{où} \\ \begin{cases} \lambda_0 = -l_0 = -\sigma\rho\sigma\rho = \sigma(\rho)\rho \\ \lambda_1 = -l_1 = -\rho\sigma\rho\sigma = \rho\sigma(\rho) \end{cases} \end{array} \right\} \leqno{(3)}
$$

d'où la relation de commutation $[\lambda_0, \lambda_1] = 1$ s'écrit aussi
$$
\sigma\rho(\sigma\rho) = \rho\sigma(\rho^2)\rho \leqno{(4)}
$$
$\sim$ encore
$$
- \sigma\rho\sigma\rho\sigma\rho = \rho\sigma\rho\sigma\rho \leqno{(5)}
$$
où le sens des signes $-$ parait capilotracté
dans le formalisme des gén. et rel. on supprime
donc le fameux signe $-$, et on remplace les
[unclear: variables (qui ne s'imposent pas) des $-$ de]
la formule par des inverses. Par exemple, la formule

\newpage

%% ===== Page 34 =====

s'écrit \marginpar{NB $\lambda_\infty^{\text{th}} \defeq -l_\infty$} \marginpar{NB $\sigma = \sigma_\infty$} [unclear: introduisant $\lambda$] :
$$
\lambda_\infty = \rho \sigma \rho^{-1} \sigma \rho = \rho(\sigma) \sigma \rho = -\rho(\sigma \rho)^{-1} \rho = \rho \sigma \rho^{-1}(\sigma) \leqno{(6)}
$$
\marginpar{NB On trouve \\ $[\rho_0, \rho_1] = [\lambda_0, \lambda_1] = (\sigma(\rho)\rho^{-1})^3$ \\ $[\lambda_0, \lambda_\infty] = [\sigma_\infty, \lambda_\infty] = (\sigma(\rho^{-1})\rho)^3 = \operatorname{int}(\sigma\rho^{-1})([\lambda_0, \lambda_1])$}
$$
[\sigma, \lambda_\infty] = 1 \leqno{(7)}
$$

On va définir un schéma en groupes (séparé, [unclear: de type fini]) sur $\mathbf{Z}$, dont les [unclear: $\hat{\mathbf{Z}}$-points s'identifient à $\hat{\pi}$].

On pose
\marginpar{[unclear: Il s'agit du vrai $M$]}
$$
M = \mathcal{G}^\circ = \underline{E}^{\{0,1,\infty\}} / \text{diagonale} \leqno{(*)}
$$

On va faire opérer $\mathcal{G}$
$$
\mathcal{G} \simeq \\operatorname{Gl}(2, \mathbf{Z}) / [\Pi_0, \Pi_0] \leqno{(4)}
$$
sur $M$, en décrivant [unclear: ainsi l'action]
$$
\mathfrak{S}_3 \ltimes \hat{M} \simeq \hat{\mathcal{G}} / \hat{M} \simeq \mathcal{G}_{0,3} / \Pi \simeq \Gamma_{0,3} \leqno{(5)}
$$

Désignons par $\lambda_0', \lambda_1', \lambda_\infty'$ les images des $\lambda_0, \lambda_1, \lambda_i = - l_i$ dans $\mathcal{G}$, identifiées à des éléments de $M = \mathcal{G}^\circ = \mathcal{G}^\circ(\hat{\mathbf{Z}})$ --- en fait, on a
$$
\lambda_0', \lambda_1', \lambda_\infty' \in M(\mathbf{Z}) \subset M(\hat{\mathbf{Z}}) = M \otimes \hat{\mathbf{Z}} \leqno{(6)}
$$
avec $\lambda_0' + \lambda_1' + \lambda_\infty' = 0$. (On a écrit additivement la loi de groupe de $M = \mathcal{G}^\circ$).

Alors $\mathcal{G}$ opère sur $M = \mathcal{G}^\circ$ par
$$
\begin{cases}
\sigma(\lambda_0') = \lambda_1', \quad \sigma(\lambda_1') = \lambda_0', \quad \sigma(\lambda_\infty') = \lambda_\infty' \\
\rho(\lambda_i') = \lambda_{i+1}' \quad \text{par définition} \\
\tau(\lambda_0') = -\lambda_1' \quad \text{où } \tau = \sigma \tau'
\end{cases} \leqno{(7)}
$$

\footnote{Ceci se voit d'abord pour l'action d'automorphismes dans $M \dots$}

\newpage

%% ===== Page 35 =====

685

L'extension $\underline{\mathfrak{S}}$ annoncée est bien sûr
celle de $§ 50$ $3^\circ)$, développé entre-temps (depuis
les premières lignes du présent $\S$, constatant
justement que j'étais gêné aux entour-
nures pour expliciter la construction) :
$$
1 \to \underset{\displaystyle \overset{\Vert}{\check{V}(M)}}{\underline{M}} \to \underline{\mathfrak{S}} \to H_{S_0} \to 1 \qquad \text{où } S_0 = \Spec \mathbf{Z}. \leqno{(9)}
$$

Pour tout schéma $S$, on posera
$$
\underline{\mathfrak{S}}^\Delta(S) \subset \underline{\mathfrak{S}}(S) \leqno{(10)}
$$
\begin{center}
$\Vert$ déf
\end{center}
Image inverse de la "diagonale" de
$\Gamma(H_S/S)$ ([unclear: ie sections $s \in \Gamma(H_S/S)$])
par l'hom. canonique $\underline{\mathfrak{S}}(S) \to H_{S}(S)$,
de sorte que par la construction des
loc. cit. on aura un iso. canonique
$$
\underline{\mathfrak{S}}^\Delta(S) \simeq M \otimes_{\mathbf{Z}} A \rtimes H \leqno{(11)}
$$
et une suite exacte can.
$$
1 \to \underset{\displaystyle \Vert}{\underset{\displaystyle A \otimes_{\mathbf{Z}} M}{\Gamma(S, \mathcal{O}_S \otimes_{\mathbf{Z}} M)}} \to \underline{\mathfrak{S}}^\Delta(S) \to H \to 1 \qquad \text{Si } S \text{ affine d'anneau } A \leqno{(12)}
$$

Prenant $A = \hat{\mathbf{Z}}$, on trouve en parti-
culier (et c'est au vu d'un fait
général sur les schémas en groupes
associés à des extensions de groupes
discrets à noyau abélien)
$$
\hat{\mathfrak{S}} \simeq \underline{\mathfrak{S}}^\Delta(\hat{\mathbf{Z}}) \leqno{(13)}
$$

\newpage

%% ===== Page 36 =====

2. Le groupe $Z_{\rho,\sigma}$ : version provisoire.

Je vais maintenant expliciter pour $\widehat{\underline{\mathfrak{S}}}$ et $\widetilde{\underline{\mathfrak{S}}} \twoheadrightarrow \widehat{\underline{\mathfrak{S}}}$ les constructions générales du § 49 VII, en prenant $\underline{\mathfrak{S}}$.
$$
\underline{\mathfrak{S}} \defeq M = \widehat{\underline{\mathfrak{S}}}^\circ \text{ (composante neutre de } \widehat{\underline{\mathfrak{S}}}) \leqno{(14)}
$$

Dans un cas comme celui-ci, ce $\underline{\mathfrak{S}}$ est justement le sous-groupe fermé du $(\rho,\sigma)$-groupe $\widehat{\underline{\mathfrak{S}}}_{\rho,\sigma}$ ($= \widehat{\underline{\mathfrak{S}}}^\wedge$ ici) engendré par $\lambda_0, \lambda_1, \lambda_\infty$, on a [unclear: à dissiper] dès les notations - là où il n'y a pas ambiguïté - d'oublier la notation $q$. Pour la suite il y a lieu de distinguer par la suite $\underline{\mathfrak{S}}_{\rho,\sigma}$ et $\widehat{\underline{\mathfrak{S}}}_{\rho,\sigma}$.

\marginpar{Je ferai la fin du coup dans le contexte schématique. En vue possible d'ailleurs d'une version sur $\mathbf{Z}_{\rho,\sigma}$.}
Nous ferons cependant les constructions dans le cadre schématique, plutôt que profini.

Le pas essentiel est de déterminer le groupe $Z_{\rho,\sigma}$

\marginpar{ici, "$\exists \alpha$" "$\exists \beta$" signifie bien sûr l'existence locale de $\alpha, \beta \dots$}
$$
Z_{\rho,\sigma} \defeq \\operatorname{Aut}(\widehat{\underline{\mathfrak{S}}}^+) \leqno{(15)}
$$
$$
= \left\{ u \in \\operatorname{Aut}(\widehat{\underline{\mathfrak{S}}}^+) \left| 
\begin{array}{l}
\text{a) } \exists \alpha \in \widehat{\underline{\mathfrak{S}}}_{\mathbf{Z}} = \{\pm 1\} \cdot M, \ u(\rho) = \alpha(\rho) \\
\text{b) } \exists \beta \in \dots, \ u(\sigma) = \beta(\sigma) \\
\text{c) } u(\underline{L}_0^\circ) = \underline{L}_0^\circ \ (\iff u(\underline{L}_0^\circ) = \underline{L}_0^\circ) \\
\text{[d) } (\text{automatiquement}) \ u(\underline{\mathfrak{S}}) \subset \underline{\mathfrak{S}}]
\end{array} 
\right. \right\}
$$

Notons que $u$ est déterminé par les valeurs

\marginpar{où $\Pi = M \cdot \{\pm 1\}$ \\ ($M$ ss-gr. d'indice 2 dans $\Pi$)}
$$
\begin{cases}
u(\rho) = \xi \rho \\
u(\sigma) = \eta \sigma
\end{cases}
\quad \xi, \eta \in \Gamma_{\Pi} \leqno{(16)}
$$
satisfaisant les conditions suivantes :

\newpage

%% ===== Page 37 =====

$$
\begin{cases}
a) & \xi\rho \text{ (localement) conjugué à } \rho \text{ via } M \\
b) & \eta\sigma \text{ \qquad --- " --- \qquad à } \sigma \text{ \qquad --- " --- } \\
c) & \text{équation du Centre} \\
   & \xi = \eta \, \ell_1^{-1} \qquad \text{où } \eta \in \Gamma M \text{ convenable} \\
   & \qquad\qquad\qquad (\text{tel que } \mu^{2\nu+1} \in \Gamma M)
\end{cases} \leqno{(17)}
$$

Remarque. Il n'est plus possible ici de désigner
par $-1$ l'élément non trivial du centre $\omega$ de
$\widehat{\mathfrak{S}}^+$, et de désigner par $-x$ (pour $x \in \widehat{\mathfrak{S}}^+$) le
produit de $x$ par celui-ci, à cause de la
confusion avec la notation $-x$ (pour $x \in M$)
dans le groupe abélien $M$ pour l'inverse de $x$ --
mais l'on a $\omega x \neq -x$, puisque $\omega x \notin M, -x \in M\dots$

On pose bien sûr
$$
\varepsilon_3(u) = \varepsilon_3(\alpha) , \quad \varepsilon_2(u) = \varepsilon_2(\beta) \leqno{(18)}
$$
(bien définis, car $\rho$ et $\rho\alpha$ ne sont conjugués
sur aucune fibre, idem pour $\sigma, \sigma^{-1}$),
comment les expliciter en termes de
$\xi, \eta$ ? On aura

$$
\begin{cases}
\xi = \alpha \rho (\alpha)^{-1} \begin{cases}
\in \Gamma M \text{ si } \alpha \in \Gamma M = \Gamma \widehat{\mathfrak{S}} \text{ i.e. } \varepsilon_3(u) = 1 \\
= \underset{\in \Gamma M}{\alpha \tau} \rho (\alpha \tau)^{-1} = \alpha_0 (- \ell_1^{-1}) \rho(\alpha_0)^{-1} \in \Gamma M \text{ si } \varepsilon_3(u) = -1 \\
\qquad\qquad\qquad\quad [unclear: \lambda_1^{-1}] \in M
\end{cases} \\
\eta = \beta \sigma (\beta)^{-1} \begin{cases}
\in \Gamma M \text{ si } \beta \in \Gamma M = \Gamma \widehat{\mathfrak{S}} \text{ i.e. } \varepsilon_2(u) = 1 \\
= \underset{\in \Gamma M}{(\beta_0 \tau)} \sigma (\beta_0 \tau)^{-1} = \omega \beta_0 \sigma (\beta_0)^{-1} \in \omega M \text{ si } \varepsilon_2(u) = -1
\end{cases}
\end{cases} \leqno{(19)}
$$

donc en tous cas

$$
\begin{cases}
\xi \in \Gamma M \\
\eta \in \Gamma (\omega^{[unclear: \nu_2(u)]} M)
\end{cases} \leqno{(19 \text{ bis})}
$$

\newpage

%% ===== Page 38 =====

d'où on tire, par l'équation des lacets (17)

$$ l_1^{\underline{\nu}} = [unclear: (\omega)^{\underline{\nu}}] \in \Gamma(\omega^{\underline{\nu}_2(u)} \underline{M}) \quad \text{i.e.} \quad \omega^{\underline{\nu}} \in \Gamma(\omega^{\underline{\nu}_2(u)} \underline{M}) \leqno (20) $$

On peut l'écrire sous la forme

$$ \underline{\nu} = 2 \underline{\nu}' + \underline{\nu}_2 \leqno (21) $$

où

$$ \begin{cases} \underline{\nu}_1, \underline{\nu}' : Z_{\rho,\sigma} \to \mathbf{G}_a \quad \text{[unclear: p.e.]} \\ \underline{\nu}_2, \underline{\nu}_3 : Z_{\rho,\sigma} \to (\mathbf{Z}/2\mathbf{Z})_{S_0} \end{cases} \leqno (22) $$

sont des homomorphismes de schémas bien définis, reliés aux autres hom. de groupes

$$ \mu : Z_{\rho,\sigma} \to \mathbf{G}_m, \quad \varepsilon_2, \varepsilon_3 : Z_{\rho,\sigma} \to \{\pm 1\}_{S_0} $$

par les formules (21) et

$$ \begin{cases} \mu = 2 \underline{\nu}_1 + 1 (= 4 \underline{\nu}' + 2 \underline{\nu}_2 + 1) \\ \underline{\varepsilon}_2 = (-1)^{\underline{\nu}_2}, \quad \underline{\varepsilon}_3 = (-1)^{\underline{\nu}_3} \end{cases} \leqno (23) $$

On aura donc, en vertu de (21)

$$ l_1^{\underline{\nu}} = l_1^{2\underline{\nu}'} l_1^{\underline{\nu}_2} = (\lambda_1)^{2\underline{\nu}'} l_1^{\underline{\nu}_2} = \lambda_1^{2\underline{\nu}'} l_1^{\underline{\nu}_2} \leqno (24) $$

Comme $l_1 = a_1 \omega$, on peut exprimer les $l_i$

$$ l_i = a_i \omega \quad \text{où } a_i \in \underline{M} \leqno (25) $$

permet de réécrire (24) sous la forme

$$ l_1^{\underline{\nu}} = \underbrace{(\lambda_1^{2\underline{\nu}'} a_1^{\underline{\nu}_2})}_{\in \Gamma \underline{M}} \omega^{\underline{\nu}_2} \quad \text{[unclear: pour } \dots ] \leqno (26) $$

Pour expliciter les $a_i$, il faut les développer en combinaisons linéaires de $\lambda_0, \lambda_1, \lambda_\infty$ (tirées d'autres [unclear: équations], si on veut [unclear: l'expliciter]).

\newpage

%% ===== Page 39 =====

$$
l_1^\nu = \lambda_1^\nu w^{\nu_2} \leqno{(24)}
$$

Posons (cf (19))
$$
\begin{cases}
\beta = \beta_0 \tau^{\nu_2} \quad (\beta \in \Gamma \underline{M}) \\
\eta_0 = \beta_0 \sigma(\beta_0)^{-1}
\end{cases} \leqno{(25)}
$$
on aura donc
$$
\eta = \beta \sigma(\beta)^{-1} = \beta_0 \sigma(\beta_0)^{-1} w^{\nu_2} = \eta_0 w^{\nu_2}
$$
et l'équation des lacets prend la forme d'une équation dans $\underline{M}$
$$
\zeta = \eta_0 \lambda_1^\nu , \quad \text{[unclear: i.e. compte tenu de (21)]} \leqno{(26)}
$$
$$
\zeta, \eta_0 \text{ sont deux sections de } \underline{M} \text{ qui se déterminent mutuellement } (\text{pour } \nu \text{ donné}). \leqno{(26\text{bis})}
$$

On a ainsi exprimé de façon commode l'équation des lacets, i.e.\ la condition c) de (17) ; il reste à exprimer a) et b), qui se formulent :
$$
\begin{cases}
a) \quad \zeta \text{ locale de la forme } \alpha_0 \rho(\alpha_0)^{-1} \ (\nu_3=0) \quad \text{ou} \quad \alpha_0 \lambda_1^{-1} \rho(\alpha_0)^{-1} \ (\nu_3=1) \\
\qquad (\text{où } \alpha_0 \in \Gamma \underline{M}) \\
b) \quad \eta_0 \text{ locale de la forme } \beta_0 \sigma(\beta_0)^{-1}
\end{cases} \leqno{(27)}
$$

(Condition dont la formulation est indépendante de $\nu, \nu_2$, mais dépend apparemment de $\nu_3$). Pour expliciter ces conditions, on va écrire $\underline{M}$ additivement, les relations s'écrivent\marginpar{On a donc posé $\alpha = \alpha_0 \tau^{\nu_3}$} :
$$
\begin{cases}
\zeta = \alpha_0 - \rho(\alpha_0) \quad (\nu_3=0) \quad \text{ou} \quad \zeta = \alpha_0 - \rho(\alpha_0) - \lambda_1 \ (\nu_3=1) \\
\eta_0 = \beta_0 - \sigma(\beta_0)
\end{cases} \leqno{(28)}
$$

\newpage

%% ===== Page 40 =====

ce qui s'exprime aussi par
$$
\begin{cases}
\xi \in \operatorname{Im}(id_{\underline{M}} - \rho_{\underline{M}}) \quad (\nu_3=0) \quad \text{ou} \quad \xi-\lambda_1 \in \operatorname{Im}(id_{\underline{M}} - \rho_{\underline{M}}) \\
\eta_0 \in \operatorname{Im}(id_{\underline{M}} - \sigma_{\underline{M}})
\end{cases} \leqno{(29)}
$$
C'est le moment de passer à des calculs matriciels explicites, en représentant $\underline{M}$ et la base $\lambda_0, \lambda_1$. Dans cette base, les opérations de $\rho_{\underline{M}}, \sigma_{\underline{M}}, \tau_{\underline{M}}$ sont donnés par
$$
\begin{cases}
\rho_{\underline{M}} = \begin{pmatrix} 0 & -1 \\ 1 & -1 \end{pmatrix} & id_{\underline{M}} - \rho_{\underline{M}} = \begin{pmatrix} 1 & 1 \\ -1 & 2 \end{pmatrix} \\
\sigma_{\underline{M}} = \begin{pmatrix} 0 & 1 \\ 1 & 0 \end{pmatrix} & id_{\underline{M}} - \sigma_{\underline{M}} = \begin{pmatrix} 1 & -1 \\ -1 & 1 \end{pmatrix} \\
\tau_{\underline{M}} = \begin{pmatrix} -1 & 0 \\ 0 & -1 \end{pmatrix}
\end{cases} \leqno{(30)}
$$

On a
$$
\det(id_{\underline{M}} - \rho_{\underline{M}}) = 3 \, , \quad \det(id_{\underline{M}} - \sigma_{\underline{M}}) = 0
$$
ce qui montre que les endomorphismes $id_{\underline{M}} - \rho_{\underline{M}}$ et $id_{\underline{M}} - \sigma_{\underline{M}}$ de $\underline{M}$ ne sont pas tout à fait des automorphismes\footnote{i.e. des isomorphismes} -- de [unclear: conoyau d'ordre] 3 pour le premier, (en [unclear: toute caractéristique]), \marginpar{rang 1} pour le second. On a donc plus intérêt à prendre comme "paramètres" non $\xi$ et $\eta_0$, mais bien carrément $\alpha_0$ et $\beta_0$ (en plus des "paramètres" discrets $\nu_2, \nu_3$ et $\nu'$), et on évitera donc de se faire chier (si possible\dots) avec les systèmes

\newpage

%% ===== Page 41 =====

$$
\widetilde{\mathcal{Z}}_{1,0} = \left\{ (\nu'_1, \nu'_2, \nu_3, \alpha_0, \beta_0) \ \middle| \ \begin{matrix} \nu'_i \in \Gamma_{\dots}, \ \nu_3 \in \Gamma_{\{0,1,\infty\}}, \ \alpha_0, \beta_0 \in \Gamma M \\ \alpha_0 - \rho(\alpha_0) = \nu_3 \lambda_1 = \beta_0 - \sigma(\beta_0) + \nu_3 \lambda_1 \end{matrix} \right\}\footnote{[unclear: NB l'élément $\eta_0$ est redondant]} \leqno{(31)}
$$

L'équation des lacets, en $\alpha_0, \beta_0$, prend alors la forme

$$
(\mathrm{id} - \rho_{\underline{M}}) \alpha_0 = (\mathrm{id} - \sigma_{\underline{M}}) (\beta_0) = (2 \nu'_1 + \nu_3) \lambda_1 \leqno{(32)}
$$

Remarque. Le schéma $\widetilde{\mathcal{Z}}_{1,0}$ lui-même, tel qu'il était défini au § 49, n'est pas lisse sur $\mathbf{Z}$, semble-t-il\footnote{Cela tient au fait qu'on [unclear: obtient] $\eta_0$ en [unclear: multipliant] par un [unclear: dénominateur]. L'homomorphisme $M \times M \to M$, $(\alpha_0, \beta_0) \mapsto (\mathrm{id} - \rho_{\underline{M}})\alpha_0 - (\mathrm{id} - \sigma_{\underline{M}})\beta_0$ est bien lisse. Mais $\lambda_1$ n'est pas [unclear: dans l'image] !}, en car. 2 et 3. Il [unclear: faudra] introduire un "vrai" $\widetilde{\mathcal{Z}}_{1,0}$, plus fin, qui [unclear: s'obtient] en substituant dans le [unclear: groupe de base] à la [unclear: constante] $\lambda_1$ la ligne de $(\xi, \eta_0)$ qui semble bien lisse. Je [unclear: sursoie] : renonçant à modifier un peu la présentation des constructions du § 49, en substituant [unclear: ainsi] au groupe $\widetilde{\mathcal{Z}}_{1,0}$ (lié aux paramètres $\xi, \eta_0$), et en le remplaçant par un groupe qui approche mieux la [unclear: limite] des $L_0$ (ou des $L_0^\#$) de $\mathcal{M}_{1,0}$ (ou $\mathcal{M}_{1,0}^\#$). On fera une présentation de ces [unclear: choses] dans le [unclear: reste du] \S suiv., et j'omets ici d'y esquisser cette construction.

\newpage

%% ===== Page 42 =====

D'où le [unclear: g-tor] (schéma en groupes $\tilde{Z}_{\bar{\Gamma}}$ avec d'une structure de groupe universelle, et opérant sur $\underline{\mathfrak{S}}^+$ via
$$
\begin{cases}
u(\rho) = (\alpha_0 - \rho(\alpha_0) + \nu_3 \lambda_1) \rho \\
u(\sigma) = (\beta_0 - \sigma(\beta_0)) w^{\nu_2} \sigma
\end{cases}
\leqno (33)
$$

Mais il faut quand même vérifier que ces formules (33) déterminent bel et bien un automorphisme $u$ de $\underline{\mathfrak{S}}^+$. La vérification essentielle : faire en sorte que la relation (7)
$$
[\sigma, \lambda_\infty]
\leqno (34)
$$
est respectée par $u$.

Je vais expliciter l'action de $u$ sur $\underline{M}$, en explicitant ses valeurs sur
$$
\begin{cases}
\lambda_0 = \omega l_0 = \omega \sigma \rho \sigma \rho \\
\lambda_1 = \omega l_1 = \omega \rho \sigma \rho \sigma
\end{cases}
$$
en termes de $x, y_0 \in \Gamma \underline{M}$
$$
\begin{cases}
x = \alpha_0 - \rho(\alpha_0) + \nu_3 \lambda_1 \\
y_0 = \beta_0 - \sigma(\beta_0)
\end{cases}
\leqno (35)
$$

On a :
$$
u(\rho) = x \rho, \quad u(\sigma) = w^{\nu_2} y_0 \sigma
\leqno (36)
$$

On aura [unclear: donc] par la relation des tresses
$$
u(\varepsilon_0) = \varepsilon_0 \quad (\mu = 2\nu_1+1 = 4\nu' + 2\nu_2 + 1)
$$
d'où

\newpage

%% ===== Page 43 =====

$$
u(l_0) = u(\varepsilon_0^2) = l_0^\mu
$$
et
$$
u(\lambda_0) = u(\omega l_0) = \underbrace{u(\omega)}_{=\omega} u(l_0) = \omega l_0^\mu = \underbrace{\omega l_0}_{\lambda_0} \underbrace{(l_0^{\mu-1})}_{\lambda_0^{\mu-1}} = \lambda_0^\mu
$$
comme il convient. (Pour me rassurer, j'ai vérifié
que la deuxième formule (33) donne bien
$u(\omega) = (y_0 \sigma)(y_0 \sigma) = (y_0 \sigma(y_0)) \sigma^2 = \omega$ i.e.
$$
u(\omega) = \omega \leqno{(37)}
$$
On aura
$$
u(\lambda_1) = u(\operatorname{int}(\rho) \lambda_0) = \operatorname{int}(\rho) u(\lambda_0) = \operatorname{int}(\rho) (\lambda_0^\mu) = \lambda_1^\mu
$$
d'où finalement, comme prévisible
\marginpar{\mu = 2\nu_1 + 1 = 4\nu' + 2\nu_2 + 1}
$$
u(\lambda) = \mu.\lambda \quad \text{pour } \lambda \in \Gamma_{\underline{M}} \leqno{(38)}
$$
On s'assure bien que $u(\sigma) = (\beta_0 - \sigma(\beta_0)) \omega^{\nu_2} \sigma$
commute avec $u(\lambda_0) = \lambda_0^\mu$.

Tous ces calculs sont un peu vaseux,
logiquement, car à la fois pour prouver
l'existence d'un tel hom. $u$ et dans les
calculs précédents je l'admets déjà (implicite-
ment, dans un groupe "plus gros", à savoir le
groupe libre engendré par $\rho, \sigma \dots$).
Une façon de les rendre rigoureux serait
d'utiliser la description amalgamée de $\underline{\widehat{\mathfrak{S}}}^+$ :
$$
\underline{\widehat{\mathfrak{S}}}^+ = \underline{\widehat{\mathfrak{S}}}_0^+ \ast_{\underline{M}_{S_0}} \underline{M} \leqno{(39)}
$$
donc définir un hom. de $\underline{\widehat{\mathfrak{S}}}^+$ dans un
$S_0$-groupe $G$ (ici $G = \underline{\widehat{\mathfrak{S}}}^+$) revient à

\newpage

%% ===== Page 44 =====

définir des homomorphismes
(40) $$M_{S_0} \xrightarrow{u_M} g^\natural, \quad \underline{\mathfrak{S}}^+_{S_0} \xrightarrow{u_{\underline{\mathfrak{S}}}} g^\natural \quad (\text{i.e. } \underline{\mathfrak{S}} \to g^{c.s.})$$
satisfaisant les deux conditions suivantes :

a) Commutativité du diagramme

(41)
\begin{tikzcd}
M_{S_0} \ar[r] \ar[d] & \underline{\mathfrak{S}}^+_{S_0} \ar[d, "u_{\underline{\mathfrak{S}}}"] \\
M_{S_0} \ar[r, "u_M"'] & g
\end{tikzcd}

et b)

(42) $u_M$ est compatible avec les actions
de $\underline{\mathfrak{S}}^+_{S_0}$ sur $M_{S_0}$ (via son action sur $M_{S_0}$)
et sur $g$ (via $u_{\underline{\mathfrak{S}}}$).

Dans le cas naturel, on aura $g = \underline{\mathfrak{S}}^+$

(43) $u_M(\lambda) = i(\mu . \lambda)$

où $i : M \longrightarrow \underline{\mathfrak{S}}^+$
est l'inclusion évidente, et

(44) $u_{\underline{\mathfrak{S}}} : \underline{\mathfrak{S}}^+_{S_0} \longrightarrow g \quad \text{défini par (33)}$

Il faudra d'abord vérifier que (33)
définit bien un hom. $\underline{\mathfrak{S}}^+_{S_0} \longrightarrow g$,
ce qui revient à $= \underline{\mathfrak{S}} \to g^{c.s.}$.

\newpage

%% ===== Page 45 =====

Pour ce utiliser la présentation de $\mathfrak{S}^+$
$$
\mathfrak{S}^+ = \left\{ \rho, \sigma \mid \rho^6 = \sigma^4 = 1, \ \rho^3 = \sigma^2, \ [\sigma, \lambda_\infty] = 1 \text{ où } \lambda_\infty = \rho \sigma \rho^{-1} \sigma \rho \right\}, \leqno{(44)}
$$
donc posant
$$
\rho' = u_{\mathfrak{S}}(\rho), \quad \sigma' = u_{\mathfrak{S}}(\sigma) \qquad (\text{cf } (33)) \leqno{(45)}
$$
il faut vérifier les relations
$$
{\rho'}^6 = {\sigma'}^4 = 1, \ [\sigma', \lambda'_\infty] = 1, \ {\rho'}^3 = {\sigma'}^2 \quad \text{où } \lambda'_\infty = \rho' \sigma' {\rho'}^{-1} \sigma' \rho'. \leqno{(46)}
$$
On commencera par vérifier
$$
{\rho'}^3 = {\sigma'}^2 = \omega \leqno{(47)}
$$
ce qui implique
$$
{\rho'}^6 = {\sigma'}^4 = 1, \quad \text{d'où } \exists \ u_{\mathfrak{S}}: \mathfrak{S}^+ \to \mathcal{G}(S_0), \ \rho \mapsto \rho', \ \sigma \mapsto \sigma'. \leqno{(48)}
$$
Ensuite, notant que
$$
\lambda'_\infty = u_{\mathfrak{S}}(\lambda_\infty),
$$
il suffira de vérifier
$$
\lambda'_\infty = \mu \lambda_\infty \leqno{(49)}
$$
Il sera évident que $\lambda'_\infty$ commute à $\sigma' = u(\sigma) = \eta_0 \sigma$ puisque $\sigma$ et $\eta_0 \in \mathcal{M}$ y commutent. Donc pour construire $u_{\mathfrak{S}}$, les vérifications essentielles sont (47) et (49) --- d'ailleurs une fois (47) prouvé, les calculs précédents aboutissant à (38) sont justifiés (au moins là où ils ont un sens, i.e.\ pour $u_{\mathcal{M}}$ et $\lambda \in \mathcal{M}$).

\newpage

%% ===== Page 46 =====

donc (49) sera bien vérifié.

On a
$$
\rho'^3 = (p \rho)^3 = (p \rho(p) \rho^2(p)) \omega
$$
$$
\sigma'^2 = (y_0 \sigma)^2 = (y_0 \sigma(y_0)) \omega
$$
et les formules (47) se réduisent à :
$$
p \rho(p) \rho^2(p) = 1 , \quad y_0 \sigma(y_0) = 1 \leqno{(50)}
$$
soit, additivement
$$
(1 + \rho + \rho^2)(p) = 0 , \quad y_0 + \sigma(y_0) = 0
$$
dont la vérification est triviale, compte tenu que
$$
\rho_M^3 = \sigma_M^2 = \operatorname{id}_M , \quad (1 + \rho_M + \rho_M^2) = \text{[unclear: crossed out]} \circ M .
$$

On a donc bien construit un homomorphisme
$$
u_{\mathcal{G}_0} : \mathcal{G}^+ \to \underline{\mathfrak{S}}^+(S_0) ,
$$
et il reste à vérifier (41) et (42).

Pour (41), on l'a vu déjà noté en usant de l'existence d'un $u_{\mathcal{G}_0}$.
— elle revient à prouver :
$$
( \quad u_{\mathcal{G}_0}(\lambda_0) = \text{[unclear: crossed out]} \lambda_0 \quad , \quad u_{\mathcal{G}_0}(\lambda_1) = \mu \lambda_1
$$
Je vais quand-même refaire la vérification, pour être sûr.
$$
\varepsilon_0' = u_{\mathcal{G}_0}(\varepsilon_0) = u_{\mathcal{G}_0}(\sigma) u_{\mathcal{G}_0}(\rho) = y_0 \sigma \cdot p \rho = y_0 \sigma(p) \sigma \rho
$$
$$
\sigma \rho = y_0 \sigma(p) \text{[unclear: crossed out]} \varepsilon_0
$$

\newpage

%% ===== Page 47 =====

Je dis que
$y'_0 = \varepsilon_0^{(M-1)}$ i.e. $y_0 \sigma(y) \omega^{V_2} = \underbrace{\varepsilon_0^{(\mu-1)}}_{l_0^{\nu'}} \implies \mu-1=2\nu_1$, soit
$y_0 \sigma(y) \omega^{V_2} = l_0^{\nu'}$ ; $\nu = 2\nu'$ i.e.
$y_0 \sigma(y) = l_0^{2\nu'} (\omega l_0 \lambda_0^{V_2}) = \underbrace{\lambda_0^{2\nu'+V_2}}_{\lambda_0^{2\nu'} \lambda_0^{-V_2}}$ i.e.
$y_0 \sigma(y) = \lambda_0^{\nu'}$ soit, en appliquant $\sigma$ et notant
que $\sigma(y_0) = y_0^{-1}$, $\sigma^2(y)=y$, $\sigma(\lambda_0)=\lambda_1$,
$y_0^{-1} y = \lambda_1(y)$ soit enfin $y = y_0 \lambda_1(y)$ $\sim$
$y = y_0 + \nu \lambda_1$ (écrit additivement)

ce qui est vrai par hypothèse (cf (31) ).
Reste à [unclear: s'assurer] (pour achever [unclear: l'étude] de $u_{\mathcal{E}_0}$ pour
les [unclear: éléments annulés] (pour [unclear: abréger on l'écrira simplement])
$u(\lambda_1) = u\rho(\lambda_0) = (\rho \mu)(\lambda_0) = \rho(u(\lambda_0)) = \rho(\lambda_0 \mu) = \lambda_1 \mu$
OK.

Reste à vérifier (42), ce qui [unclear: se réduit] à :
commutativité de $u_M = \mu\ id_M$ à
$(\rho_M, \rho'_M)$, $(\sigma_M, \sigma'_M)$ - ce qui signifie que
$\rho_M$ et $\rho'_M$ d'une part, $\sigma_M$ et $\sigma'_M$ d'autre part,
agissent [unclear: également], ce qui est trivialement évident.

Ainsi, le groupe $\tilde{\mathcal{E}}_{\rho, \sigma}$ est bel et
bien construit - du moins, en tant que
relevant un [unclear: sous-groupe] $Z$, muni des [unclear: morphismes]
de structure
(51)
\begin{tikzcd}
\tilde{Z}_{\rho, \sigma} \ar[r, shift left, "f_0"] \ar[r, shift right, "f_{S_0}"'] & \mathcal{E}_{0, S_0} \ar[r] & \mathcal{E}_{S_0}
\end{tikzcd}
d'où $\mu = 4\nu' + 2 + V_2 : \tilde{Z}_{\rho, \sigma} \to \mathcal{E}_{S_0}$
$\varepsilon_2, \varepsilon_5 : \tilde{Z}_{\rho, \sigma} \rightrightarrows \mu_{S_0} = \mu\ id_{S_0}$

\newpage

%% ===== Page 48 =====

et un morphisme de faisceaux sur $S_0$
$$
\widetilde{Z}_{\rho,\sigma} \longrightarrow \\operatorname{Aut}_{gr} (\underline{\mathfrak{G}}^+) \leqno{(52)}
$$
donnant lieu à un morphisme
$$
\widetilde{Z}_{\rho,\sigma} \longrightarrow \\operatorname{Aut}_{gr} (\underline{\mathfrak{G}}^+) \times \mathbf{G}_m \times \mu_{S_0} \times \mu_{S_0} \leqno{(53)}
$$
$$
u \longmapsto u, \mu(u), \varepsilon_2(u), \varepsilon_3(u)
$$
des calculs déjà faits toutes ces fois dans des contextes divers donnent à penser que ce morphisme (qui n'est même pas un hom. $\neq 33$) est un isomorphisme de groupes, pour une structure de groupes convenable sur $\widetilde{Z}_{\rho,\sigma}$, uniquement déterminée par elle l'ex (cas $\neq 2,3$, et conduisant par voie de descente schématique).\marginpar{[unclear: un peu casse-pied car $\widetilde{Z}_{\rho,\sigma} \to G_m \times \mu_{S_0} \times \mu_{S_0}$ n'est pas... induit par $M \times M \dots \to \rho_0$ \dots c'est un peu long ... cas 3 $\to \Gamma$ (55)]} On aura
$$
1 \to \widetilde{Z}_{\rho,\sigma}^{+0} \to \widetilde{Z}_{\rho,\sigma} \to \mathbf{G}_m \times \mu_{S_0} \times \mu_{S_0} \to 1 \leqno{(54)}
$$
explicité : sous-schéma en groupes
$$
\widetilde{Z}_{\rho,\sigma}^{+0} \subset M \times M \times \mathbf{G}_a
$$
$$
= \left\{ (\alpha_0, \beta_0, \gamma_1) \;\middle|\; \begin{matrix} \alpha_0 + d\gamma' = 0 \\ (1+\rho_M)\alpha_0 = (1-\sigma_M)\beta_0 + 2\gamma'\lambda_1 \end{matrix} \right\}
$$

On aura donc
$$
1 \to Z_{\rho F}^{+0} \to \widetilde{Z}_{\rho,\sigma}^{+0} \xrightarrow{q} \mathbf{G}_a \to 1 \leqno{(55)}
$$

\newpage

%% ===== Page 49 =====

$$
\widetilde{Z}_{\rho, \sigma}^{+00} \defeq \{ u \in \Gamma \widetilde{Z}_{\rho, \sigma} \mid \nu'(u) = \nu_2(u) = \nu_3(u) = 0 \} \leqno (57)
$$
$$
\parallel
$$
$$
\{ (\alpha_0, \beta_0) \in \Gamma(M \times M) \mid (1 - \rho_M)\alpha_0 = (1 - \sigma_M)\beta_0 \}
$$

\marginpar{[unclear: en $\mathbf{G}_m$ à un diviseur infinitésimal de l'action de $2 \times 3$ !]} Ce schéma en groupes $\widetilde{Z}_{\rho, \sigma}^{+00}$ est bien (en car. $\neq 2,3$) de dimension relative 2 - il est en fait isom.\ canonique, en tant que module, à notre
$$
\widetilde{Z}_{\rho, \sigma}^{+00} \simeq \mathbf{V}(P) \leqno (58)
$$
où $P$ est le $\mathbf{Z}$-module libre de rang 2 défini comme noyau de l'hom.\ de modules
$$
0 \longrightarrow P \longrightarrow M \times M \longrightarrow M \longrightarrow 0 \leqno (59)
$$
$$
(\alpha_0, \beta_0) \longmapsto (1 - \rho_M)\alpha_0 - (1 - \sigma_M)\beta_0
$$
L'hom.\ $M \times M \longrightarrow M$ s'explicite en posant
$$
\begin{cases}
\alpha_0 = a_0 \lambda_0 + a_1 \lambda_1 & \text{d'où } \alpha_0 - \rho \alpha_0 = (a_0 + a_1) \lambda_0 + (2a_1 - a_0) \lambda_1 \\
\beta_0 = b_0 \lambda_0 + b_1 \lambda_1 & \text{d'où } \beta_0 - \sigma \beta_0 = (b_0 - b_1) \lambda_0 + (b_1 - b_0) \lambda_1
\end{cases}
\leqno (60)
$$
d'où
$$
(\alpha_0, \beta_0) \longmapsto \underbrace{(a_0 + a_1 - (b_0 - b_1))}_{c_0} \lambda_0 + \underbrace{(2a_1 - a_0 - (b_1 - b_0))}_{c_1} \lambda_1 \leqno (61)
$$
C'est un hom.\ [unclear: surjectif sur un sous-module de $M$] ; les éléments de l'image sont les
$$
\{ c_0 \lambda_0 + c_1 \lambda_1 \mid \underbrace{c_0 + c_1}_{= 3a_1} \equiv 0 \pmod 3 \} \leqno (62)
$$

Remarques. Si on prend pour $u$ défini sur $\widetilde{Z}$ provenant d'un élément de $M[0]$, alors

\newpage

%% ===== Page 50 =====

$\mu \in \hat{\mathbf{Z}}^*$, $\mu \equiv 1 - 2\nu_3 \pmod{3}$,
et la condition sur $(\mu, \nu_2, \nu_3)$ pour être
dans l'image de $\widetilde{Z}^{+\infty}_{\rho, \sigma} \to \mathbf{G}_m \times \mu_2 \times \mu_3$,
qui s'écrit en général (vu la valeur de $c$ (162), et de (32))
$$ 2V' - \nu_2 + \nu_3 \equiv 0 \pmod{3} $$
devient, en multipliant par 2
$$ 4V' - 2\nu_2 + 2\nu_3 \equiv 0 \pmod{3} $$
i.e.
$$ \underbrace{4V' - 2\nu_2}_{\mu - 1} + 2\nu_3 \equiv 0 \pmod{3} $$
i.e.
$$ \mu \equiv 1 - 2\nu_3 \pmod{3}, $$
et elle est automatiquement satisfaite.

C'est pourquoi j'ai envie de modifier la définition de $\widetilde{Z}^{+\infty}_{\rho, \sigma}$, en définissant

ici $2c = 6V''$, et on a en fait $c = 3V''$, on aura
$$ \widetilde{Z}^{+\infty}_{\rho, \sigma} \defeq \left\{ (V', V'', \nu_2, \nu_3, \alpha_0, \beta_0) \in \mathbf{G}_a \times \mathbf{G}_a \times \{0,1\}_{S_0} \times \{0,1,2\}_{S_0} \times M \times M \;\middle|\; \begin{array}{l}
\text{a) } 2V' - \nu_2 + \nu_3 = 3V''\footnote{En comparant $\mu \equiv 1 - 2\nu_2$ (4) et $\mu \equiv 1 - 2\nu_3$ (6), $\mu = 1 - 2\nu_2 + 4V'$, $\mu = 1 - 2\nu_3 + 6V''$, i.e.\ $-\nu_2 + 2V' = -\nu_3 + 3V''$ i.e.\ $2V' - \nu_2 + \nu_3 = 3V''$.} \\
\text{b) } \mu \in \mathbf{G}_m \text{ soit } \mu \defeq 1 - 2\nu_2 + 4V' \\
\text{c) équation dans } \operatorname{Centr} \\
\quad (1 - \rho_M)(\alpha_0) - (1 - \sigma_M)(\beta_0) = 3V'' \lambda_1
\end{array} \right\} \leqno{(63)} $$

\newpage

%% ===== Page 51 =====

3. Les groupes $\mathbb{O}_N^*$.

On va introduire une construction auxiliaire, en définissant un schéma en groupes qui correspond à la donnée des quantités
$\mu \in \mathbb{G}_m$, $\nu', \nu'' \in \mathbb{G}_a$, $\nu_2, \nu_3 \in \Gamma_{\{0,1\}_{S_0}}$, satisfaisant
$$ \mu = 3\nu' - 2\nu_2 + 1 = 6\nu'' - 2\nu_3 + 1, \quad 2\nu' - \nu_2 = 3\nu'' - \nu_3 \leqno (63) $$

De façon générale, soit $N \in \mathbb{N}^*$ un entier fixé.
(Ce cas que nous étudierons par la suite est $N=12$), et considérons le schéma (sur $\mathbf{Z}$) associé à l'extension de groupes discrets
$$ 1 \to \mathbf{Z} \xrightarrow{N} \mathbf{Z} \to \mathbf{Z}/N\mathbf{Z} \to 1 \leqno (64) $$
qui s'insère dans une suite exacte
$$ 1 \to E^1 \to \mathbb{O}_N \to (\mathbf{Z}/N\mathbf{Z})_{S_0} \to 1 \leqno (65) $$

Rappelons la construction simple
$$ \mathbb{O}_N = (\mathbf{Z})_{S_0} \times N\mathbf{Z} \leqno (66) $$

Mais je dis qu'il y a sur $\mathbb{O}_N$ une structure d'Anneau unique, dont l'addition est la structure de groupe déjà envisagée, et où la multiplication est obtenue ainsi : Une section de $\mathbb{O}_N$ s'écrit localement de la forme $n + N[\lambda]$, \dots la règle de calcul
$$ \begin{cases} a) \ (n + N[\lambda]) + (n' + N[\lambda']) = (n+n') + N[\lambda + \lambda'] \\ b) \ n + N[\lambda] = n' + N[\lambda'] \iff n + N\lambda = n' + N\lambda' \text{ dans } \mathbb{G}_a \end{cases} \leqno (67) $$

\newpage

%% ===== Page 52 =====

On multiplie alors les sections par
$$
(n+N[\lambda]) (n'+N[\lambda']) = n n' + N [n \lambda' + n' \lambda + \lambda \lambda'] \leqno{(68)}
$$
on vérifie que c'est compatible aux relations (67 b), que c'est biadditif, associatif, que $1+N[0]$ est unité... Autre façon de voir la structure d'anneau. Soit plus généralement $A$ une $\mathbf{Z}$-algèbre unifère, munie d'un idéal bilatère $I$ qui soit $\mathbf{Z}$-module libre de t.f., considérons l'idéal $NA$ et la structure d'extension
$$
0 \to I \to A \to \overset{\mathcal{B}}{A/I} \to 0 \leqno{(69)}
$$
d'un schéma en groupes comm. noté $A_{[N]}$, qui s'obtient comme quotient de $A_{S_0} \times I$ quotient (semi-direct)
$$
0 \to I \to A_{S_0} \times I \to A_{[N]} \to 0 \leqno{(70)}
$$
$$
\lambda \mapsto (\lambda, -\lambda)
$$
et on trouve que $A_{S_0} \times I$ muni d'une structure d'anneau produit 1/2-direct, et que [unclear: son quotient]

\newpage

%% ===== Page 53 =====

l'image de $\mathcal{J} \overset{67a}{\to} A$ est un idéal ...
On trouve bien une structure d'anneau
sur $A_{[N]}$, et une suite exacte de
$$
0 \to \mathcal{J} \to A_{[N]} \to B_{S_0} \to 0 \leqno (71)
$$
où $A_{[N]} \to B_{S_0}$ est un homomorphisme de
schémas en anneaux. Passant aux
groupes multiplicatifs, on trouve
une suite exacte
$$
1 \to (1+\mathcal{J})^* \to A_{[N]}^* \to \underset{\begin{subarray}{c} || \\ (B^*)_{S_0} \end{subarray}}{(B_{S_0})^*} \to 1 \leqno (72)
$$
où pour $A_{[N]}^*$, $(B_{S_0})^*$ et $B^*$, le
$*$ désigne le foncteur sous-groupe multiplicatif,
et où $(1+\mathcal{J})^*$ a le sens évident
$$
\begin{gathered}
(1+\mathcal{J})^* \subset 1+\mathcal{J} \\
|| \\
\{ (1+\lambda) \mid 1+\lambda \text{ inv. dans } A \} \\
\simeq \\
\{ (\lambda, \lambda') \in \mathcal{J} \times \mathcal{J} \mid \lambda+\lambda'+\lambda\lambda' = 0 \}
\end{gathered} \leqno (73)
$$

Ainsi on arrive, dans le cas particulier
$A = \mathbf{Z}$, $\mathcal{J} = N\mathbf{Z}$
$$
1 \to (1+N\mathbf{Z})^* \to \mathcal{O}_N^* \to (\mathbf{Z}/N\mathbf{Z})_{S_0}^* \to 1 \leqno (74)
$$

\newpage

%% ===== Page 54 =====

$(1+N\mathbf{Z})^*$ est défini comme
$$
(1+N\mathbf{Z})^* \defeq \{ \lambda \in \mathbf{A}^1 = \mathbb{G}_a \mid 1+N\lambda \text{ inversible} \} \leqno{(75)}
$$
Les $\mathbf{O}_N$ pour $N$ variable forment un syst. projectif d'Anneaux, les $\mathbf{O}_N^*$ forment un syst. projectif de schémas en groupes commutatifs.

Je n'ai pas pu regarder de plus près la structure de ce système projectif, ni examiné la limite projective - question qui se pose plus généralement pour un $A$ comme ci-dessus. On a à expliciter l'hom. canonique
$$
\mathbf{O}_N^* \to \mathbf{O}_1^* = \mathbf{O}^* = \mathbb{G}_m \leqno{(76)}
$$
dont le noyau est le sous-groupe des
$$
\{ 1+N\lambda \mid 1+N\lambda=1 \text{ i.e. } N\lambda=0 \} \leqno{(77)}
$$
qui en tant que schéma en groupes est isomorphe à ${}_N \mathbb{G}_a$,

\newpage

%% ===== Page 55 =====

%% [page 55 missing or failed: 503 UNAVAILABLE. {'error': {'code': 503, 'message': 'The service is currently unavailable.', 'status': 'UNAVAILABLE'}}]

\newpage

%% ===== Page 56 =====

qui se traduisent en des lois de groupes pour l'addition mod 2 sur $\{0,1\}_{S_0}$ (où $\varepsilon_2 = (-1)^{\nu_2} = 1-2\nu_2$, $\varepsilon_3 = (-1)^{\nu_3} = 1-2\nu_3$). Les « quatre cas de figures » de $(\mathbf{Z}/12\mathbf{Z})^\times$ se donnent dans le tableau

$$
\begin{array}{r|c|c}
\frac{\varepsilon_2}{\varepsilon_3} & \begin{matrix} +1 \\ \text{i.e. } \nu_2=0 \end{matrix} & \begin{matrix} -1 \\ \text{i.e. } \nu_2=1 \end{matrix} \\
\hline
\begin{matrix} +1 \\ \text{i.e. } \nu_3=0 \end{matrix} & 1 & +5 \\
\hline
\begin{matrix} -1 \\ \text{i.e. } \nu_3=1 \end{matrix} & -5 & -1
\end{matrix} \leqno{(82)}
$$

donc on trouve une section ensembliste de (79) en prenant la section de $\mathfrak{O}_{12}^\times$ de la forme

$$
\mu_0 = (-5)^{\nu_2} (+5)^{\nu_3} \quad \begin{cases} \equiv (-1)^{\nu_3} & (6) \\ \equiv (-1)^{\nu_2} & (4) \end{cases} \leqno{(83)}
$$

Donc une section générale de $\mathfrak{O}_{12}^\times$ s'écrit

$$
\mu = \underbrace{(-5)^{\nu_2}(+5)^{\nu_3}}_{\mu_0} + 12 \lambda_0 \leqno{(84)}
$$
$$
\begin{cases}
(\nu_2, \nu_3) \in \Gamma(\{0,1\}_{S_0} \times \{0,1\}_{S_0}) \\
\lambda_0 \in \Gamma(\widehat{\mathbf{Z}})
\end{cases}
$$

$$
\begin{aligned}
5^{\nu_3} &= 1 + 4\nu_3 \\
(-5)^{\nu_2} &= (-1)^{\nu_2} 5^{\nu_2} = (1-2\nu_2)(1+4\nu_2) \\
&= 1 + 2\nu_2 - 8\nu_2^2 = 1 - 6\nu_2 \marginpar{NB $\nu_2 = \nu_2^2$}
\end{aligned}
$$

donc

$$
\mu_0 = (1 + 4\nu_3)(1 - 6\nu_2) \leqno{(85)}
$$
$$
= 1 + 4\nu_3 - 6\nu_2 - 24\nu_2\nu_3
$$

\newpage

%% ===== Page 57 =====

Si on ( (84) et (85) )
$$ \mu = \mu_0 + 12 [\nu_0] = 1 + 4\nu_3 - 6\nu_2 + 12(\nu_0 - 2\nu_2 \nu_3) $$
on constate (en s'intéressant aux congruences mod 4)
$$ \mu = \underline{1} - 2\nu_2 + 4 [\nu_3 - \nu_2 - 3\nu_2 \nu_3 + 3\nu_0] \leqno (86) $$
et par congruences mod 6
$$ \mu = \underline{1} - 2\nu_3 + 6 [\nu_3 - \nu_2 - 2\nu_2 \nu_3 + 2\nu_0] \leqno (87) $$
où le crochet désigne la section $\nu'$ (resp. $\nu''$) correspondante
à valeurs dans $\hat{\mathbf{Z}}$ , le $\underline{1}$ souligné désignant la section unité
de $G_m = \mathbb{O}_1^*$. On aura bien
$$ 2\nu' - \nu_2 = 2\nu_3 - 3\nu_2 - 6\nu_2 \nu_3 + 6\nu_0 \text{ \qquad "} $$
$$ 3\nu'' - \nu_3 = 2\nu_3 - 3\nu_2 - 6\nu_2 \nu_3 + 6\nu_0 \text{ \qquad "} $$

En résumé \footnote{les relations en groupes lisses}
Proposition. Introduisons sur $\mathbb{O}_{12}^*$ \marginpar{$G_{\mathbf{Q}} = $} les
fonctions $\nu_2, \nu_3$ (à valeurs dans $\{0,1\}_{S_0} \subset \hat{\mathbf{Z}}$)
$\nu', \nu''$ (à valeurs dans $\hat{\mathbf{Z}}$) et $\mu$ (à valeurs dans $\mathbb{O}_1^* \subset \mathbb{O}_{12}^*$)
définies par les conditions
$$ \mu = \text{im. de } x \text{ dans } \mathbb{O}_1^* \text{ par l'hom. de groupes } \mathbb{O}_{12}^* \to \mathbb{O}_1^* $$
$$
\begin{cases}
(-1)^{\nu_2} = \text{im. de } x \text{ dans } (\mathbf{Z}/4\mathbf{Z})^*_{S_0} \\
(-1)^{\nu_3} = \text{im. de } x \text{ dans } (\mathbf{Z}/6\mathbf{Z})^*_{S_0}
\end{cases} \leqno (88)
$$
d'où
$$
\begin{aligned}
\mu &= (-5)^{\nu_2} (+5)^{\nu_3} + 12 [\nu_0] \\
&= (1+4\nu_3)(1-6\nu_2) + 12 \nu_0 \\
&= 1 + 4\nu_3 - 6\nu_2 + 12 [\nu_0 - 2\nu_2 \nu_3]
\end{aligned} \leqno (89)
$$
pour une section $\nu_0$ bien déterminée de $G_{\mathbf{Q}}$.

\newpage

%% ===== Page 58 =====

Posons de plus
$$
\begin{cases}
\nu' = \nu_3 - \nu_2 - 3\nu_2 \nu_3 + 3\nu_0 \\
\nu'' = \nu_3 - \nu_2 - 2\nu_2 \nu_3 + 2\nu_0 .
\end{cases} \leqno{(90)}
$$
On a alors les relations\footnote{où $\varepsilon_p = (-1)^{\nu_p}$ d'où $\varepsilon_p - 1 = -2\nu_p$ ($p \in \{2, 3\}$).}
$$
\begin{cases}
\mu = 1 - 2\nu_2 + 4\nu' \quad (= \varepsilon_2 + 4\nu') \\
\mu = 1 - 2\nu_3 + 6\nu'' \quad (= \varepsilon_3 + 6\nu'') \\
\underbrace{-\nu_2 + 2\nu'}_{\nu} = \underbrace{-\nu_3 + 3\nu''}_{\nu} \quad i.e. \quad 2\nu' - \nu_2 + \nu_3 = 3\nu'' .
\end{cases} \leqno{(91)}
$$

NB La condition $\mu = 1$ équivaut à la condition
$$
\nu_2 = \nu_3 = \nu_0 = 0 \leqno{(92)}
$$
et implique
$$
\nu = \nu' = \nu'' = 0 \leqno{(93)}
$$

\newpage

%% ===== Page 59 =====

4. Le groupe $\tilde{Z}_{\bar{p},\sigma}$ et la suite exacte $1 \to N_{\bar{p},\sigma} \to \tilde{N}_{\bar{p},\sigma} \to \mathbb{O}_{12}^* \to 1$.

$\tilde{Z}_{\bar{p},\sigma}$ est défini comme plongé dans $\mathbb{O}_{12}^* \times M \times M$ :
$$
\tilde{Z}_{\bar{p},\sigma} \subset \mathbb{O}_{12}^* \times M \times M \leqno{(94)}
$$
$$
\{(\mu, \alpha_0, \beta_0) \mid (1-\rho_M)(\alpha_0) - (1-\sigma_M)(\beta_0) = 3 \nu''(\mu)\lambda_1\}
$$

Notons qu'on a
$$
(1-\rho_M)(\lambda_1 - \lambda_0) = 3\lambda_1 \leqno{(95)}
$$
donc la relation des cocycles dans (94) prend aussi la forme
$$
(1-\rho_M)(\alpha_0 - \nu''(\lambda_1 - \lambda_0)) - (1-\sigma_M)(\beta_0) = 0 \leqno{(96)}
$$

Je vais admettre [unclear: un instant] qu'on a une structure de groupe naturelle sur $\tilde{Z}_{\bar{p},\sigma}$, pour laquelle la projection $\tilde{Z}_{\bar{p},\sigma} \to \mathbb{O}_{12}^*$ est un homomorphisme de groupes. On trouve maintenant une suite exacte
$$
1 \to \tilde{Z}_{\bar{p},\sigma}^{+0} \to \tilde{Z}_{\bar{p},\sigma} \to \mathbb{O}_{12}^* \to 1 \leqno{(97)}
$$
où
$$
\tilde{Z}_{\bar{p},\sigma}^{+0} = \\operatorname{Ker}\left( \tilde{Z}_{\bar{p},\sigma} \to \mathbb{O}_{12}^* \right) \leqno{(98)}
$$
$$
\simeq \{(\alpha_0, \beta_0) \mid (1-\rho_M)(\alpha_0) - (1-\sigma_M)(\beta_0) = 0\}
$$

Le schéma en groupes est lisse en caractéristique $\neq 3$, mais pas en caractéristique $3$, où sa dimension passe

\newpage

%% ===== Page 60 =====

de la dimension relative 2 à la dimension relative 3.

Le fait d'une fibre exceptionnelle de dimension plus grande ne peut, il me semble, être éliminé par une [unclear: astuce-pipo], ou en introduisant des définitions ad-hoc pour le "meilleur" groupe $\widetilde{Z}_{\rho, \sigma}$ : il y a "plus" d'automorphismes de $\mathfrak{S}_K$ (pour $K$ un corps algébriquement clos donné), ayant les propriétés envisagées, si $k$ est de caractéristique 3, que dans les autres cas de caractéristiques.

Cependant il me semble qu'on peut construire un sous-schéma en groupes de $\widetilde{Z}_{\rho, \sigma}$, que j'ai envie de noter $N_{\rho, \sigma}$, qui est lisse sur $S_0 = \operatorname{Spec}(\mathbf{Z})$ et coïncide avec $\widetilde{Z}_{\rho, \sigma}$ au dessus de $S_{0,3} = \operatorname{Spec} \mathbf{Z}[1/3] = S_0 \setminus \{3\}$. Pour le deviner je note que l'équation de base (96), la première au moins - est l'analogue dans le sous-module de $M$ correspondant à la congruence
$$
c_0 + c_1 \equiv 0 \ (3) \quad (\text{cf. } (62) \text{ [unclear: au-dessus]})
$$
et si on introduit le sous-groupe d'indice

\newpage

%% ===== Page 61 =====

3 correspondant à $M_0$ de $M$
$$
\begin{array}{c}
M_0 \subset M \\
\parallel \defeq \\
\{ c_0 \lambda_0 + c_1 \lambda_1 \mid c_0 + c_1 \equiv 0 \ (3) \}
\end{array}
\leqno (99)
$$
d'où un fibré vectoriel $\underline{M}_0$ sur $S_{0,2}$ ---
et en fait un homomorphisme de schémas
$$
\widetilde{Z}_{p, \sigma} \longrightarrow \underline{M}_0 \leqno (100)
$$
$$
(\mu, \alpha_0, \beta_0) \longmapsto (1 - \sigma_M)(\alpha_0 + \nu''(\lambda_1 - \lambda_0)) + (1 - \sigma_M)(\beta_0)
$$
qu'on peut expliciter ainsi : on a un homomorphisme de $\mathbf{Z}$-modules
$$
M \times M \times \mathbf{Z} \longrightarrow M_0 \leqno (101)
$$
$$
(\alpha_0, \beta_0, \nu'') \longmapsto (1 - \sigma_M)(\alpha_0 + \nu''(\lambda_1 - \lambda_0)) + (1 - \sigma_M^2)(\beta_0)
$$
d'où un hom. pour les schémas en modules correspondants
$$
\underline{M} \times \underline{M} \times \underset{\mathbb{G}_a}{\underline{\mathbf{Z}}} \longrightarrow \underline{M}_0 \leqno (102)
$$
qu'on compose avec
$$
\widetilde{Z}_{p, \sigma} \longrightarrow \underline{M} \times \underline{M} \times \underline{\mathbf{Z}} \leqno (103)
$$
$$
(\mu, \alpha_0, \beta_0) \longmapsto (\alpha_0, \beta_0, \nu''(\mu))
$$
On peut donc considérer l'équation des cocycles (96) comme équation dans $\underline{M}_0$ - elle devient ainsi plus stricte,

\newpage

%% ===== Page 62 =====

d'où un sous-schéma fermé de $\widetilde{\underline{Z}}_{p,\sigma}$,
$$
\underline{N}_{p,\sigma} = \Big\{ (\mu, \alpha_0, \beta_0) \in \widetilde{\underline{Z}}_{p,\sigma} \mid \text{l'équation des \underline{cocycles} (96) valable dans } \underline{M}_0 \Big\}
\leqno{(104)}
$$
\footnote{NB le groupe $\underline{N}_{p,\sigma}$ devrait peut-être se noter ainsi $\underline{M}_{p,\sigma}[\Sigma_0]$, sauf erreur je l'éviterai. En fait plutôt ce sera $(\underline{M}_{p,\Sigma}/\mu_p)[\Sigma_0]$ avec les notations [unclear: de la suite].}
Considérons donc le morphisme
$$
\underline{N}_{p,\sigma} \longrightarrow \underline{O}_{12}^*
$$
il est clair que ce morphisme est un épi, la fibre d'un point $\underline{v} = (v_1, v_2, v_0)$ étant formée des $(\alpha_0, \beta_0) \in \underline{M} \times \underline{M}$ satisfaisant (96) dans $\underline{M}_0$, [unclear: se ramenant à] l'équation linéaire dans $\underline{M}_0$ :
$$
(1-\rho_{\underline{M}})(\widetilde{\alpha}_0) - (1-\sigma_{\underline{M}})(\widetilde{\beta}_0) = 0
\leqno{(105)}
$$
Or l'homomorphisme de modules
$$
\underline{M} \times \underline{M} \longrightarrow \underline{M}_0
$$
$$
(\widetilde{\alpha}_0, \widetilde{\beta}_0) \longmapsto (1-\rho_{\underline{M}})(\widetilde{\alpha}_0) - (1-\sigma_{\underline{M}})(\widetilde{\beta}_0)
\leqno{(106)}
$$
est un épi, car il se factorise sur l'homomorphisme
$$
M \times M \longrightarrow M_0
$$
correspondant qui est lui-même un épi.

Donc $\underline{N}_{p,\sigma}$ est lisse sur $S_0$. Et c'est un sous-schéma fermé de $\widetilde{\underline{Z}}_{p,\sigma}$.

\newpage

%% ===== Page 63 =====

qui coïncide avec lui au-dessus de $S_0 \setminus \{3\}$, il s'ensuit que c'est l'adhérence schématique
$$
N_{\hat{\rho}, \sigma} = \widetilde{Z}_{\hat{\rho}, \sigma} \mid (\text{Spec } \mathbf{Z} \setminus \{3\}) \leqno{(107)}
$$
(adh. schématique)

ce qui implique que c'est aussi un sous-schéma en groupes.

C'est peut-être le moment de regarder d'un peu plus près le sous-$\mathbf{Z}$-module $M_0$ de $M$,
$$
\begin{aligned}
M_0 &= \mathbf{Z}(\lambda_1 - \lambda_0) + \mathbf{Z}(\lambda_2 - \lambda_1) \\
&= \mathbf{Z} \underbrace{(\lambda_1 - \lambda_0)}_{\delta} + \mathbf{Z} \rho_M \underbrace{(\lambda_1 - \lambda_0)}_{\delta} \\
&= \mathbf{Z} \delta \oplus \mathbf{Z} (1-\rho_M)(\delta) \\
&= \mathbf{Z} \delta + \mathbf{Z} \rho_M \delta
\end{aligned} \leqno{(108)}
$$

Notons que $M_0$ est stable sous $\rho$ et $\sigma$ de façon générale par construction, pour tout $x \in M$, on a
$$
\begin{array}{lcl}
& x \equiv \rho_M(x) \quad (M_0) & \text{i.e.\ } (1-\rho_M)(x) \in M_0 \\
\text{et} & x \equiv \sigma_M(x) \quad (M_0) & \phantom{\text{i.e.\ }} (1-\sigma_M)(x) \in M_0
\end{array} \leqno{(109)}
$$

Il est bien sûr stable aussi par $\tau_M = -\operatorname{id}_M$.

Donc $M_0$ est un module pour $\mathfrak{S}$, d'où un sous-groupe de $\mathfrak{S}$ distingué dans $\mathfrak{S}_0$. C'est un sous-groupe distingué de $\mathfrak{S}_0 = \\operatorname{Gl}(2, \mathbf{Z})$,

\newpage

%% ===== Page 64 =====

contenant $[\Pi_0, \Pi_0]$, d'indice $3$ dans $\Pi_0$.

Je présume qu'il contient le sous-groupe
de congruences relatif à $12$ - i.e. distingué,
qu'il correspond à un sous-groupe de
$S\ell(2, \mathbb{Z}/12\mathbb{Z})$. Il est d'indice
$3 \cdot 12 = 36$ dans $S\ell(2, \mathbb{Z})$, [unclear: rayé] $G\ell(2, \mathbb{Z})$,
tandis que $$S\ell(2, \mathbb{Z}/12\mathbb{Z}) \simeq \underset{48}{S\ell(2, \mathbb{Z}/4\mathbb{Z})} \times \underset{24}{S\ell(2, \mathbb{Z}/3\mathbb{Z})}$$
il s'agit donc d'un sous-groupe (c'est même un sous-groupe distingué)
$32$ de $S\ell(2, \mathbb{Z}/12\mathbb{Z})$, engendré par les
éléments images de
$$ (\sigma(\rho)\rho^{-1})^3 = [\lambda_0, \lambda_1], \text{ ou } [\sigma_0, \lambda_\infty] $$
cf (7)
$\delta = \lambda_1 \lambda_0^{-1}$
$\rho(\delta)$

ou encore le sous-groupe distingué engendré par
le seul élément dans $S\ell(2, \mathbb{Z}/12\mathbb{Z})$

$$ (110) \quad \left\{ \begin{array}{l} \delta = \lambda_1 \lambda_0^{-1} = l_1 l_0^{-1} = \left(\begin{array}{cc} 1 & 0 \\ 2 & 1 \end{array}\right) \left(\begin{array}{cc} 1 & +2 \\ 0 & 1 \end{array}\right) \\ = \left(\begin{array}{cc} 1 & +2 \\ 2 & 5 \end{array}\right) \end{array} \right. $$

La question est de si ce sous-groupe
distingué de $S\ell(2, \mathbb{Z}/12\mathbb{Z})$ est bel et
bien d'ordre $32$, donc d'indice $36$ - un
ordre qui pourrait être plus grand, donc l'indice
plus petit. Déjà quel est l'ordre de $\bar{\delta}$ ?

\newpage

%% ===== Page 65 =====

mod 12 ? On a
$$
\delta^2 = \begin{pmatrix} 1 & 2 \\ 2 & 5 \end{pmatrix} \begin{pmatrix} 1 & 2 \\ 2 & 5 \end{pmatrix} = \begin{pmatrix} 5 & 12 \\ 12 & 29 \end{pmatrix} \equiv \begin{pmatrix} 5 & 0 \\ 0 & 5 \end{pmatrix} \quad (12)
$$
d'où
$$
\delta^4 \equiv 1 \quad (12) ,
$$
c'est déjà un bon point ! On a d'ailleurs
$$
\begin{array}{ll}
\begin{cases} \delta \equiv 1 & (2) \\ \delta^2 \equiv 1 & (4) \end{cases} &
\begin{array}{l}
\delta \neq 1 \ (4), \ \delta \neq 1 \ (3) \\
\delta^2 \equiv -1 \ (3)
\end{array}
\end{array} \leqno{(112)}
$$

Considérons
$$
\\operatorname{Sl}(2, \mathbf{Z}/12\mathbf{Z}) \simeq \underbrace{\\operatorname{Sl}(2, \mathbf{Z}/4\mathbf{Z})}_{48} \times \underbrace{\\operatorname{Sl}(2, \mathbf{Z}/3\mathbf{Z})}_{24} \xrightarrow{u \times v} \underbrace{\mathfrak{S}_0^+ / \Pi_0}_{12} \times \underbrace{\mathfrak{S}_3^+}_{3} \leqno{(113)}
$$
où
$$
u : \\operatorname{Sl}(2, \mathbf{Z}/4\mathbf{Z}) \to \mathfrak{S}_0^+ / \Pi_0
$$
est l'hom.\ canonique surjectif, dont le noyau se réduit à --- les trois éléments non nuls du noyau sont justes $\lambda_0, \lambda_1, \lambda_\infty$ permutés entre eux par $\rho$, et $\lambda_1 \lambda_0^{-1} = \lambda_1 \lambda_0 = \lambda_\infty$ etc.\ ici les $\lambda_i$ sont permutés par $\rho$ --- et
$$
v : \\operatorname{Sl}(2, \mathbf{Z}/3\mathbf{Z}) \to \mathfrak{S}_3^+
$$
est l'hom.\ composé
\begin{tikzcd}[column sep=large]
\\operatorname{Sl}(2, \mathbf{Z}/3\mathbf{Z}) \arrow[r, "\text{épi}"] & \mathfrak{A}_4 \arrow[r, "\text{épi}"] \arrow[d, phantom, sloped, "\simeq"] & \mathfrak{S}_3^+ \\
& \mathfrak{A}_{\mathbf{P}^1(\mathbf{Z}/3\mathbf{Z})} &
\end{tikzcd}
dont le noyau est d'ordre 8. Le noyau de $u \times v$ est d'ordre $4 \times 8 = 32$, et il contient évidemment $\delta$. Reste à montrer que c'est bien le sous-groupe distingué engendré par $\delta$ (dont l'ordre a priori était [unclear: un] multiple de 32), et

\newpage

%% ===== Page 66 =====

ce sous-groupe image de $\Pi_0 \subset \operatorname{Sl}(2,\hat{\mathbb{Z}})$, c'est-à-dire dans $\Pi_0$,
que l'image inverse d'un sous-groupe de
l'ét. précédente de $\operatorname{Sl}(2, \mathbb{Z}/12\mathbb{Z})$.

Résumons les résultats obtenus :

Proposition. Considérons dans $\hat{\mathfrak{S}}_0^+ = \operatorname{Sl}(2, \hat{\mathbb{Z}})$
la tour de sous-groupes invariants (et dans $\hat{\mathfrak{S}}_0$)
(114)
\begin{tikzcd}
\hat{\mathfrak{S}}_0^+ \ar[r, hookleftarrow, "6"] & \Pi = \Pi[2] \ar[r, hookleftarrow, "2"] \ar[d, equal, "\text{def}"'] & \Pi_0 \ar[r, hookleftarrow, "3"] \ar[d, phantom, "\text{sous-groupe N}"'] & \Pi_{00} \ar[r, hookleftarrow, "2!"] \ar[d, phantom, "\text{sous-groupe inv.}"'] & (\Pi_0, \Pi_0) \\
& \{u \in \hat{\mathfrak{S}}_0^+ \mid u \equiv 1 (2)\} & \text{engendré par } \rho, \dots & \text{engendré par les} & \\
& & \text{[unclear]} & \lambda_1 \lambda_0^{-1} = \rho_1 \rho_0^{-1} &
\end{tikzcd}

Considérons d'autre part le groupe quotient
(115) $$\operatorname{Sl}(2, \mathbb{Z}/12\mathbb{Z}) \simeq \hat{\mathfrak{S}}_0^+ / \Pi[12]$$
$$ \simeq \operatorname{Sl}(2, \mathbb{Z}/4\mathbb{Z}) \times \operatorname{Sl}(2, \mathbb{Z}/3\mathbb{Z}) $$

et la tour de sous-groupes suivants

[MARGIN: NB $(\hat{\mathfrak{S}}_0^+)^{ab} \simeq \mathbb{Z}/12\mathbb{Z}$ (cf 42 bis)
l'identif... en $H/\Pi_{00}^\bullet$ donne un homom (116)
d'ordre 3 $-$ surj. d'un ? en un
$\operatorname{Sl}(2, \mathbb{Z}/4\mathbb{Z}) / \operatorname{Sl}(2, \mathbb{F}_2)' \times \operatorname{Sl}(2, \mathbb{Z}/3\mathbb{Z}) / Syl_{2,3} \simeq \hat{\mathfrak{S}}_0^+ / \Pi_{00}$]

\begin{tikzcd}
H = \ar[d, phantom, "\cup" description, "6" left] & \operatorname{Sl}(2, \mathbb{Z}/4\mathbb{Z}) \times \operatorname{Sl}(2, \mathbb{Z}/3\mathbb{Z}) & \text{ordre } 48 \times 24 \\
\Pi^\bullet = \ar[d, phantom, "\cup" description, "2" left] & \operatorname{Sl}(2, \mathbb{F}_2) \times \operatorname{Sl}(2, \mathbb{Z}/3\mathbb{Z}) & \text{ordre } 8 \times 24 \\
\Pi_0^\bullet = \ar[d, phantom, "\cup" description, "3" left] & \operatorname{Sl}(2, \mathbb{F}_2)' \times \operatorname{Sl}(2, \mathbb{Z}/3\mathbb{Z}) & \text{ordre } 4 \times 24 \\
\Pi_{00}^\bullet = & \operatorname{Sl}(2, \mathbb{F}_2)' \times Syl_{2,3} \ar[d, "\simeq"'] & \text{ordre } 4 \times 8 \\
& \text{[unclear]} \simeq D_4 &
\end{tikzcd}

(117)
a) $\operatorname{Sl}(2, \mathbb{F}_2) \simeq \operatorname{Ker}(\operatorname{Sl}(2, \mathbb{Z}/4\mathbb{Z}) \to \operatorname{Sl}(2, \mathbb{Z}/2\mathbb{Z}))$
b) $\operatorname{Sl}(2, \mathbb{F}_2)' \subset \operatorname{Sl}(2, \mathbb{F}_2)$, sous-groupe distingué $\simeq \mathbb{F}_2 \times \mathbb{F}_2$ dont les trois él. $\neq 0$ sont $\lambda_0, \lambda_1, \lambda_\infty$ (images dans $\operatorname{Sl}(2, \mathbb{Z}/4\mathbb{Z})$) - c'est un supplémentaire dans la [unclear] du sous-groupe multiplicatif $\mu_2 \simeq \mathbb{F}_2$ des unit. [unclear] de $\operatorname{Sl}(2, \mathbb{Z}/4\mathbb{Z})$.
c) $Syl_2 \subset \operatorname{Sl}(2, \mathbb{Z}/3\mathbb{Z})$, sous-groupe de Sylow, d'ordre 8 (qui est distingué)

\newpage

%% ===== Page 67 =====

Ceci posé, la tour de sous-groupes (114) est image inverse de la tour (116).

Par acquit de conscience, je vais reprendre l'esquisse de démonstration précédente. On note que l'homomorphisme $\mathfrak{S}_0^+ \to H$ applique les groupes figurant dans (114) dans ceux de (116) - pour $\Pi$ et $\Pi_0$, on sait d'ailleurs depuis un moment que ce sont les images inverses de $\Pi', \Pi_0'$ ! (pour ceci, il suffisait d'examiner $\\operatorname{Sl}(2, \mathbf{Z}/4\mathbf{Z})$, au lieu de $\\operatorname{Sl}(2, \mathbf{Z}/12\mathbf{Z})$ ...). Ainsi l'image inverse de $\Pi_{00}'$ dans $\mathfrak{S}_0^+$ contient $\Pi_{00}$. Comme les deux sont d'indice $36$, ces deux groupes sont égaux, cqfd. (le même argument pourrait s'appliquer au cas de $\Pi_0, \Pi_0'$ ...)

Corollaire L'homomorphisme canonique
$$
\mathfrak{S}_0^\pm = \operatorname{Gl}(2,\mathbf{Z}) \to \operatorname{Gl}(2,\mathbf{Z}/12\mathbf{Z}) / (\Pi_{00}' \simeq \\operatorname{Sl}(2,\mathbf{F}_2)' \times \operatorname{Syl}_2)
$$
se factorise par $\mathfrak{S} = \operatorname{Gl}(2,\mathbf{Z})/[\Pi_0,\Pi_0]$ étudié dans le présent § 5, et l'image inverse de $\Pi_{00}'$ n'est autre que le sous-groupe $M_0$ de (99) (108).

\newpage

%% ===== Page 68 =====

Les considérations précédentes rendent
assez naturel la considération des revête-
ments\footnote{c'est un revêtement $M_{1,1}$} d'ordre 3 de $U_{0,3 \overline{\mathbf{Q}}}$ correspon-
dant au sous-groupe $\Pi_{00}$ de $\Pi_0 \simeq \pi_1(U_{0,3 \overline{\mathbf{Q}}})$.
On vérifie tout de suite que c'est
une courbe de genre 1\footnote{c'est une courbe elliptique d'invariants $0$ (courbe de module $j=0$, dirais-je)} - S'il
m'arrive de travailler avec
$[\Pi_{00}, \Pi_{00}]$ au lieu de $[\Pi_0, \Pi_0]$, on
trouve un groupe extension de
$\Pi_0 / [\Pi_0, \Pi_0] \simeq M$, par un groupe de rang 2
qui est "elliptique", et sur
lequel $M$ va opérer de façon
certainement non triviale - par la suite
il me faudrait examiner cette action,
et la mettre en relation avec l'autre
qui [unclear: y est subordonnée] :
$$
\text{l'hom. } \\operatorname{Sl}(2, \mathbf{Z})^{\wedge} \longrightarrow \\operatorname{Sl}(2, \hat{\mathbf{Z}}).
$$

\newpage

%% ===== Page 69 =====

5) Les groupes $\underline{M}_{\rho,\sigma}$, $\underline{M}_{\rho,\sigma}^+$ (via $\underline{N}_{\rho,\sigma}$).

À tout $l \in \hat{\Pi}_0$ correspond un autom
intérieur $int(l)$, qui est l'identité sur
$\underline{M}_0$. On aura, tautologiquement

(118) $$ \begin{cases} int(l)(\rho) = \alpha(\rho) \\ int(l)(\sigma) = \beta(\sigma) \end{cases} \quad \alpha = \beta = l $$

Mais [unclear] (rel. de lacets C relative : $f_1=1$ d. $\nu=0$)
$$ (1-\rho_M)(\underset{l}{\alpha}) = (1-\sigma_M)(\underset{l}{\beta}) $$
i.e.
$$ \rho_M(l) = \sigma_M(l) \quad ? $$

Posant 
(119) $$ l = c_0 \lambda_0 + c_1 \lambda_1 $$
la relation [unclear] s'écrit
$$ c_0 \lambda_1 + c_1 \underset{-\lambda_0-\lambda_1}{\lambda_\infty} = c_0 \lambda_1 + c_1 \lambda_0 $$
i.e.
$$ - c_1 \lambda_0 + (c_0 - c_1) \lambda_1 = c_1 \lambda_0 + c_0 \lambda_1 $$
i.e. $c_1=0$ i.e. $l = c_0 \lambda_0$. Donc c'est
sans remords qu'on se restreint à
considérer des $l \in \widehat{L}_0$ ; correspondant

(120) $$ l = c_0 \lambda_0 \in \widehat{L}_0 = Im (\hat{\mathbb{Z}} \xrightarrow{u} \underline{M}) $$
$$ c_0 \mapsto c_0 \lambda_0 $$

On a un demi diagramme
(121) 
\begin{tikzcd}
\mathbb{G}_a \simeq \widehat{L}_0 \ar[r] \ar[d, hook] & \underline{N}_{\rho,\sigma} \\
\underline{G}^+
\end{tikzcd}

\newpage

%% ===== Page 70 =====

l'image de
(122) $$L_0^{\underline{M}_0} \to \underline{N}_{\rho,\sigma}$$
est invariante - l'hom (122) étant en
fait compatible avec l'action de $\underline{N}_{\rho,\sigma}$
- via $p_r$ sur $L_0^{\underline{M}_0} \simeq \tilde{G}_0$ (i.e. induits
par les opérations de $\underline{N}_{\rho,\sigma}$ sur $\tilde{G}^+ \supset \underline{M}_a \to L_0^{\underline{M}}$
et opérations de [unclear] sur $\tilde{G}^+$ déduites
de (121) sont identiques. Donc d'après
le topos général on trouve un
schéma en groupes $\underline{\underline{M}}_{\rho,\sigma}$ et un
diagramme de suites exactes df $\underline{\Gamma}_{\rho,\sigma} = \underline{N}_{\rho,\sigma}/L_0^{\underline{M}}$

[MARGIN: avec les notati. du $a$ $§ 99$ $\underline{VI}$, le gr. noté $\underline{\underline{M}}_{\rho,\sigma}$ serait sans doute noté $\underline{M}_{\rho,\sigma}/H_p$ ; ici $\underline{M}_{\rho,\sigma}$ serait $\underline{M}_{\rho,\sigma}$ (loc. cit. 8°)]

(123)
\begin{tikzcd}
1 \ar[r] & L_0^{\underline{M}} \ar[r] \ar[d] & \underline{N}_{\rho,\sigma} \ar[r] \ar[d] & \underline{\Gamma}_{\rho,\sigma} \ar[r] \ar[d, "\simeq"'] & 1 \\
1 \ar[r] & \tilde{G}^{+\Delta} \ar[r] & \underline{\underline{M}}_{\rho,\sigma} \ar[r] & \underline{\Gamma}_{\rho,\sigma} \ar[r] & 1
\end{tikzcd}

Remplaçant $\tilde{G}^+$ par $\underline{M}$ , on trouve un $\underline{M}_{\rho,\sigma}$ , et
(124)
\begin{tikzcd}
1 \ar[r] & L_0^{\underline{M}} \ar[r] \ar[d] & \underline{N}_{\rho,\sigma} \ar[r] \ar[d] & \underline{\Gamma}_{\rho,\sigma} \ar[r] \ar[d, "\simeq"'] & 1 \\
1 \ar[r] & \underline{M} \ar[r] & \underline{M}_{\rho,\sigma} \ar[r] & \underline{\Gamma}_{\rho,\sigma} \ar[r] & 1
\end{tikzcd}

D'ailleurs on a dans (124)
dans (123) (comparer $§ 99$ $\underline{VI}$ 5° ...), le
(125) $$M_{\rho,\sigma} \subset \underline{M}_{\rho,\sigma} \dots$$

\newpage

%% ===== Page 71 =====

Il faut préciser la structure de $\underline{\Gamma}_{p, \sigma}$, en notant que $\underline{L}_{\underline{0}}^M \subset \underline{N}_{p, \sigma}^{0+}$, donc
$$
1 \to \underline{\Gamma}_{p, \sigma}^{0+} \to \underline{\Gamma}_{p, \sigma} \to \mathfrak{S}_{12}^* \to 1 \leqno{(126)}
$$
$$
1 \to \underline{L}_{\underline{0}}^M \to \underline{N}_{p, \sigma}^{0+} \to \underline{\Gamma}_{p, \sigma}^{0+} \to 1 \leqno{(127)}
$$
Considérons l'hom. de $\mathbf{Z}$-modules
$$
K : (M \times M \to M_0) \to \underline{M} \leqno{(128)}
$$
$$
(\alpha_0, \beta_0) \mapsto (1-\rho_M)(\alpha_0) - (1-\sigma_M)(\beta_0)
$$
noyau formé des $(\alpha_0, \beta_0)$ tels que
$$
\begin{cases}
c_0 \defeq a_0 + a_1 + (b_1 - b_0) = 0 \\
c_1 \defeq 2a_1 - a_0 - (b_1 - b_0) = 0
\end{cases} \text{ cf (61)}
$$
Soit $\underline{N}^{0+}$ le noyau, de sorte qu'on a
$$
\underline{N}_{p, \sigma}^{0+} \simeq \underline{N}^{0+} \leqno{(129)}
$$
Considérons l'hom.
$$
\mathbf{Z} \to \underline{N}^{0+} \leqno{(130)}
$$
$$
n \mapsto (n \lambda_0, n \lambda_0)
$$
l'image est formée des systèmes $(a_0, a_1, b_0, b_1)$ tels que $a_1 = b_1 = 0$, $a_0 = b_0$,

d'où suite exacte
$$
0 \to \mathbf{Z} \to \underline{N}^{0+} \to \mathbf{Z}^3 
$$
$$
(a_0, a_1, b_0, b_1) \mapsto (a_1, b_1, a_0 - b_0)
$$
d'où $\underline{N}^{0+} / \mathbf{Z} \simeq \mathbf{Z}^3 \simeq \underline{L}_{\underline{0}}^M$, de façon [unclear: évidente].

\newpage

%% ===== Page 72 =====

donc une suite exacte de $\mathbf{Z}$-modules
$$
1 \to \underset{n \mapsto (n\lambda_0, n\lambda_0)}{\mathbf{Z}} \to N^{0+} \to \underset{\simeq N^{0+}/\mathbf{Z}}{\Gamma^{0+}} \to 0 \leqno{(131)}
$$
avec $\Gamma^{0+}$ $\mathbf{Z}$-module libre de rg 1.

On en conclut
$$
\Gamma_{p,\sigma}^{0+} \simeq \Gamma^{0+} \quad (\simeq \mathbf{G}_a) \leqno{(132)}
$$

Ainsi, les schémas en groupes qu'on vient de construire sont lisses sur $\operatorname{Spec} \mathbf{Z}$.

Peut-on définir une section \marginpar{[unclear: liée à des rel de relation 2]} naturelle de l'extension $\mathcal{M}_{p,\sigma}$ de $\Gamma_{p,\sigma}$ par $M$ ?

Compte tenu de l'équation en les composantes qui décrit $N_{p,\sigma}$, il semble qu'une telle section [unclear: dès qu'on] choisit un scindage des
$$
\begin{cases}
\text{[rayé: ] } 0 \to \underset{\operatorname{rg} 1}{\mathbf{Z}} \to N^{0+} \to M \times M \to M_0 \\
0 \to \underset{\operatorname{rg} 1}{\mathbf{Z}} \to N^{0+} \to \underset{\operatorname{rg} 1}{\Gamma^{0+}} \to 0
\end{cases} \leqno{(133)}
$$
(extensions de $\mathbf{Z}$-modules). Mais il faut vérifier que les sections "ensemblistes" de schémas ainsi définies se promeuvent en des [unclear: sections] multiplicatives...

\newpage

%% ===== Page 73 =====

5°) $Z_{\rho,\sigma}$ et $G_{\rho,\sigma} = \mathbb{O}_{12}^* G_{12}^+$ + $G_{\rho,\sigma}^+ \simeq \mathbb{O}_{12}^* \cdot \underline{M} = \underline{M}_{\rho,\sigma} \simeq \underline{M}$.
Une fois intr. duit $\mathbb{O}_{12}^*$ et $M_0 \subset M$, il me semble que le groupe $Z_{\rho,\sigma}$ ten lieu et place de $\tilde{Z}_{\rho,\sigma}$, dès l'équa - tion tirant compte de $\alpha, \beta$ et via ses données de $\xi, \eta$ redevient sympathique :

(134)
$$ \underline{Z}_{\rho,\sigma} = \operatorname{Aut}(\underline{\mathcal{S}}^+) \times \mathbb{O}_{12}^* $$
$$ \{ (u, \mu) \mid $$
a) $u(\rho) = \xi \rho$ (loc.) (conjugué à $\rho$, par un $\alpha \in \widehat{G_{12}^+}$ tel que $\chi(\alpha) = \varepsilon_3(\mu)$)
b) $u(\sigma) = \eta \sigma$ (loc.) (conjugué à $\sigma$, par un $\beta \in \widehat{G_{12}^+}$ tel que $\chi(\beta) = \varepsilon_2(\mu)$)
c) équation des bouts $\xi_0 = \eta_0 \nu''(\mu)$ soit $\xi_0 = \eta_0 + 3 \nu'' \lambda_1$ dans $M_0$

où $\nu'' = \nu''(\mu) \in \dots$ (cf pur. p. 685, et ci-dess.!)

(135)
$$ \begin{cases} \xi_0 = [- \nu_3 \lambda_1] \in \Gamma M_0 & \nu_3 = \nu_3(\mu) \\ \eta_0 = [\nu_2 \lambda_2] \in \Gamma M_0 & \nu_2 = \nu_2(\mu) \end{cases} $$

On doit supposer de plus que

(136)
$$ \begin{cases} \xi_0 \text{ locale dans l'im. de } (1 - \rho_M)(M) \\ \eta_0 \text{ locale dans l'im. de } (1 - \sigma_M)(M) \end{cases} $$

et on aura

(137)
$$ \underline{Z}_{\rho,\sigma} \simeq \tilde{\Sigma}_{\rho,\sigma} = \left\{ (\mu, \xi_0, \eta_0) \in \mathbb{O}_{12}^* \times M_0 \times M_0 \;\middle|\; \begin{matrix} 1) \; \xi_0 = \eta_0 + (\nu'' \lambda_1) \text{ dans } M_0 \\ \quad \text{où } \nu'' = \nu''(\mu) \\ 2) \text{ on a (136)} \end{matrix} \right\} $$

$G$

$$ (1 - \rho_M) : M \to M_0 $$
$$ (1 - \sigma_M) : M \to \mathbb{Z} = C_a \to M_0 $$
$$ (\lambda_0, \lambda_0+d_1, \lambda_1) \mapsto (b_1 - b_0) \mapsto (\lambda_1 - \lambda_0) = \lambda \delta $$

\newpage

%% ===== Page 74 =====

donc (136) est une condition sur $\eta_0$ seulement, savoir
$$
\eta_0 \in \operatorname{Im} \Big( \underset{\lambda \mapsto \lambda \delta}{G_a \to M_0} \Big) = \Ker \Big( \underset{c_0\lambda_0+c_1\lambda_1 \mapsto c_0+c_1}{M_0 \to G_a} \Big) \leqno{(136\text{ bis})}
$$

Notons qu'on a
$$
\begin{cases} M_0 = \underset{\lambda_1-\lambda_0}{\mathbf{Z} \delta} \oplus \underset{(1-\rho)\delta}{\mathbf{Z} . (3\lambda_1)} & \text{d'où} \\ M_0 = \mathbf{Z} \delta \times \mathbf{Z} . (3\lambda_1) \end{cases} \leqno{(138)}
$$
et par suite on a
$$
\Sigma_{\rho,\sigma} \simeq O_{12}^* \times \underset{\in G_a}{\mathbf{Z} \delta} \leqno{(139)}
$$
$$
(\underline{\mu}, \xi_0, \eta_0) \mapsto (\underline{\mu}, \eta_0)
$$
car $\xi_0$ se récupère sur $\underline{\mu}, \eta_0$ par l'équation des lacets
$$
\xi_0 = (1-\sigma_M)^{-1} \big( \eta_0 + \nu''(3\lambda_1) \big)
$$

NB Si on se donne par contre $\xi_0$, on en tire par l'équation des lacets et (138) la valeur de $\eta_0$, et celle de $\nu''$ - donc on connait en plus $\nu_2$ et $\nu_3$, i.e. on déduit la valeur de $\nu_0$ de (sur une [unclear: base 2 des] sections régulières de $\mathcal{G}$) $\nu_0$, i.e. finalement $\underline{\mu}$. Donc sur un

\newpage

%% ===== Page 75 =====

[unclear: Pour $S$ de 2-congruences, on a une]
$$
\Sigma_{p,\sigma}^{\text{eff}}(S) \hookrightarrow (\underline{M} \times \{0,1\} \times \{0,1\})(S) \leqno (140)
$$
$$
(\underline{\mu}; \xi_0, \eta_0) \mapsto (\bar{\xi}_0, \nu_2(\underline{\mu}), \nu_3(\underline{\mu}))
$$
l'image étant formée des systèmes
$(\xi_0, \nu_2, \nu_3)$ tels que, posant
$$
\xi_0 = \nu_0\delta + \nu''(3\lambda_1) \quad (\nu_0, \nu'' \in \widehat{G}_a(S))
$$
qu'en vertu de (138), on ait
$$
\begin{cases}
\nu'' + \nu_2 - \nu_3 \in 2\hat{\mathbf{Z}} \quad \text{[divisible par 2]} \\
(-5)^{\nu_2} (+5)^{\nu_3} + 6(\nu'' + \nu_2 - \nu_3 + 2\nu_2\nu_3) \text{ inv.} \\
\quad 1 - 2\nu_2 - 2\nu_3 + 6\nu''
\end{cases} \leqno (141)
$$
En tout état de cause on note que
$\Sigma_{p,\sigma}^{\text{eff}}$ au-dessus de $Z_{p,\sigma}$ est lisse de dimension relative 2.
$$
1 \to \widehat{G}_a \to Z_{p,\sigma} \to O_{12}^* \to 1 \leqno (142)
$$
(l'isom. $\widehat{G}_a \isommap Z_{p,\sigma}^0 \isommap \Sigma_{p,\sigma}^{\text{eff} 0}$ étant donné par $\lambda \mapsto (\underline{\mu}=1, \xi_0=\lambda\delta, \eta_0=-\lambda\delta)$)

On posera
$$
L \defeq \{ \lambda.\delta \mid \lambda \in \widehat{G}_a \} \subset \underline{M} \leqno (143)
$$
et on définit
$$
\begin{tikzcd}
\widehat{G}_a \arrow[r] & Z_{p,\sigma} \arrow[r, "\isommap"] & \Sigma_{p,\sigma}^{\text{eff} 0} \\
\lambda \arrow[r, mapsto] & (\operatorname{int}(u), 1) \arrow[r, mapsto] & ((u\mu)(e), (u\sigma_n)(e)) \\
"\lambda\delta" & & =(\lambda\delta, -\lambda\delta)
\end{tikzcd} \leqno (144)
$$

\newpage

%% ===== Page 76 =====

d'où

(145) $$L^\natural \overset{\sim}{\longrightarrow} Z_{\rho,\sigma}^{+0}$$

ce qui donne

(146) $$\Gamma_{\rho,\sigma} \overset{dfn}{=} Z_{\rho,\sigma}/L^\natural \simeq \mathbb{O}_{12}^* \quad .$$

Si on construit maintenant $G_{\rho,\sigma}$
comme

(147) $$G_{\rho,\sigma} = (Z_{\rho,\sigma} \cdot \underline{\underline{C}}^+) / L^\natural$$

on a le diagramme

(148)
\begin{tikzcd}
& & & \mathbb{O}_{12}^* \ar[d, equal] \\
1 \ar[r] & L^\natural \ar[r] \ar[d, hook] & Z_{\rho,\sigma} \ar[r] \ar[d, hook] & \Gamma_{\rho,\sigma} \ar[r] \ar[d, equal] & 1 \\
1 \ar[r] & \underline{\underline{C}}^+ \ar[r] & G_{\rho,\sigma} \ar[r] & \Gamma_{\rho,\sigma} \ar[r] & 1
\end{tikzcd}

D'ailleurs on a un scindage canonique

de $Z_{\rho,\sigma}$ sur $\Gamma_{\rho,\sigma} = \mathbb{O}_{12}^*$

(149)
\begin{tikzcd}
\mathbb{O}_{12}^* \ar[r] & Z_{\rho,\sigma} \ar[r, "\sim"] & \Sigma_{\rho,\sigma}^{1/\natural} \\
\mu \ar[r, mapsto] & \tilde{\mu} \ar[r, mapsto] & (\mu, \xi_0 = \nu''(3\lambda_1), \eta_0 = 0)
\end{tikzcd}

d'où

(150) $$ \begin{cases} Z_{\rho,\sigma} \simeq \mathbb{O}_{12}^* \cdot L^\natural \\ G_{\rho,\sigma} \simeq \mathbb{O}_{12}^* \cdot \underline{\underline{C}}^+ \end{cases} $$

L'opération de $\mathbb{O}_{12}^*$ sur $L^\natural$ se fait
via $\mathbb{O}_{12}^* \to \operatorname{Aut}(M)$ homothéties, celle
de $\mathbb{O}_{12}^*$ sur $\underline{\underline{C}}^+$ se fait par

(151) $$ \begin{cases} u_\mu(c) = c & (= \delta_1 c \text{ en écriture add. pour } C^+) \\ u_\mu(x) = \mu x & \text{si } x \in M \end{cases} $$

[MARGIN:
NB On pose
$G_{\rho,\sigma} = \mathbb{O}_{12}^* \cdot \underline{\underline{C}}^+$
$\cup$
$Z_{\rho,\sigma} = \mathbb{O}_{12}^* \cdot L^\natural$
]

\newpage

%% ===== Page 77 =====

79) Les groupes $\mathcal{N}_\rho$, [unclear: $\mathcal{N}_\rho^{\text{gr}}$], $\mathcal{N}_\rho^\sigma$ etc.

Étudions le sous-groupe
$$
\mathcal{N}_\rho\footnote{Notation provisoire, cf plus bas p. 677} \subset G_{\rho,\sigma} \leqno{(152)}
$$
$$
\defeq \{ u \in G_{\rho,\sigma} \mid u(\rho) = \rho^{\varepsilon_3(\underline{u})} \}
$$

L'homomorphisme composé
$$
\mathcal{N}_\rho \hookrightarrow G_{\rho,\sigma} \xrightarrow{\underline{\chi}} \widehat{\mathbf{Q}}_{12}^*
$$
est épi, [unclear: comme on va le voir] pour nous en assurer.

Soit $\underline{\mu}$ une section locale de $\widehat{\mathbf{Q}}_{12}^*$, correspondant à $\nu_2, \nu_3, \nu'' \dots$ On aura
$$
u_{\underline{\mu}}(\rho) = \lambda_1^{\nu''} \rho = (\lambda_1^{3\nu''} \lambda_1^{-\nu_3}) \rho = \operatorname{int}(\delta^{\nu''} \tau_3^{\nu_3}) (\rho)
$$
[passage biffé illisible]

d'où
$$
(\delta^{\nu''} \tau_3^{\nu_3})^{-1} u_{\underline{\mu}} (\rho) = \rho \leqno{(153)}
$$

\marginpar{NB on note $\underline{\chi}$ l'hom. $G_{\rho,\sigma} \to \widehat{\mathbf{Q}}_{12}^*$}
$$
\begin{cases}
(\delta^{\nu''} \tau_3^{\nu_3})^{-1} u_{\underline{\mu}} = \tau_3^{-\nu_3} \delta^{-\nu''} u_{\underline{\mu}} \in \mathcal{N}_\rho \\
\underline{\chi}(v) = \underline{\chi}((\delta^{\nu''} \tau_3^{\nu_3})^{-1}) \cdot \underline{\chi}(u_{\underline{\mu}}) \\
\underline{\chi}(\tau_3^{\nu_3}) = (-1)^{\nu_3}_{\varepsilon_3} = \pm 1
\end{cases}
$$

Donc posant
$$
v = (\tau_3^{\nu_3} \delta^{\nu''})^{-1} u_{\underline{\mu}} = (\omega_2^{\varepsilon_2} \tau_3)^{\nu_3} \delta^{-\nu''} u_{\underline{\mu}}
$$

on a
$$
\underline{\chi}(v) = \pm 1, \quad \text{avec } v = \tau_3^{-\nu_3} \delta^{-\nu''} u_{\underline{\mu}} \in \mathcal{N}_\rho \leqno{(154)}
$$
$$
(\delta = \lambda_1 \lambda_0^{-1}, \nu_3 = \nu_3(\underline{u}), \nu'' = \nu''(\underline{u}))
$$

\newpage

%% ===== Page 78 =====

On aura donc une suite exacte
$$
1 \to Z_\rho \to N_\rho \to \begin{array}[t]{c} \hat{\mathbf{Z}}_{12}^* \\ \Vert \\ \Gamma_{\rho,\sigma} \end{array} \to 1 \leqno (155)
$$
et
$$ Z_\rho = \operatorname{Centr}_{\mathfrak{S}^+}(\rho) $$
On a, si
$$
\begin{cases}
H_{S_0}^+ \defeq \mathfrak{S}^+ / \mathfrak{S}_0 \simeq (H^+)_{S_0} \\
\qquad \text{[unclear: car on a]} \simeq \mathcal{M} \\
\therefore \quad H^+ = \mathfrak{S}^+ / \mathcal{M}
\end{cases} \leqno (156)
$$
un hom.\ canonique
$$ Z_\rho \to \operatorname{Centr}_{H_{S_0}^+}(\rho) = L_\rho \quad (\simeq (\mathbf{Z}_{12})_{S_0}) \leqno (157) $$
qui est épi (puisqu'il y a une section, provenant de l'inclusion $L_\rho \subset Z_\rho$) et dont le noyau est
$$ Z_\rho^\circ = Z_\rho \cap \mathcal{M} = \left\{ x \in \mathcal{M} \mid {}^\rho x = x \text{ i.e. } (1-\rho_M)x = 0 \right\} \leqno (158) $$
Ce noyau est nul au-dessus de $S_0$,
mais est de dim 1, en car 3. Par contre on sent bien que la "bonne" définition
de $Z_\rho^\circ$ doit être fournie par une équation
$$ (1-\rho_M)x = 0 \quad \text{dans } \mathcal{M}_0, \leqno (159) $$

\newpage

%% ===== Page 79 =====

ce qui revient en fait \marginpar{d'ailleurs} à remplacer
$\underline{N}_\rho$ défini dans (152), par l'adhérence
schématique de sa restriction au dessus
de $S_0 - \{s_3\}$. On trouve donc un
sous-groupe, meilleur, que désormais
je note $\underline{N}_\rho$, inséré dans la suite exacte
$$
1 \to \underline{L}_\rho \to \underline{N}_\rho \xrightarrow{\chi_\rho} \underline{\mathbb{O}}_{12}^* \to 1 \leqno{(159)}
$$

Posons
$$
\begin{cases}
\underline{G}_{\rho,\sigma}^\natural = \underline{\mathbb{O}}_{12}^* \cdot \underline{\mathfrak{S}}^\natural \subset \underline{\mathbb{O}}_{12}^* \cdot \underline{\mathfrak{S}}^+ = \underline{G}_{\rho,\sigma} \\
\underline{N}_\rho^\natural = \underline{N}_\rho \cap \underline{G}_{\rho,\sigma}^\natural
\end{cases} \leqno{(160)}
$$
on trouve alors
$$
\underline{N}_\rho^\natural \cap \underline{L}_\rho = \underline{M} \cap \underline{L}_\rho = \{1\} \leqno{(161)}
$$
donc
$$
\underline{N}_\rho^\natural \hookrightarrow \underline{\mathbb{O}}_{12}^* \leqno{(161')}
$$

je dis que l'image est
\marginpar{donc on a $\underline{N}_\rho^\natural \subset \underline{Z}_\rho$, il vaut mieux la noter $\underline{Z}_\rho^\natural$}
$$
\gamma_{\rho,\sigma}' = {\underline{\mathbb{O}}_{12}^*}' \defeq \left\{ \lambda \in \Gamma \mid \varepsilon_3(\lambda)=1 \text{ i.e. } \nu_3(\lambda)=0 \right\} \leqno{(162)}
$$
Que l'image contienne ${\underline{\mathbb{O}}_{12}^*}'$ résulte
de (154), montrant qu'elle est continue.

Soit $u = f^{-1} u_f \quad f \in \underline{G}_{\rho,\sigma}^\natural = \underline{\mathbb{O}}_{12}^* \cdot \underline{\mathfrak{S}}^\natural$, écrite
comme une section locale $f \mapsto u_f$
qui donne un élément dans $\underline{N}_\rho^\natural$ d'où...

\newpage

%% ===== Page 80 =====

$\operatorname{int}(f^{-1}) u_f(\rho) = \rho^{\varepsilon_3} \quad \text{i.e.} \quad \lambda_1^v \rho = \operatorname{int}(f) \rho^{\varepsilon_3}$

en réduisant mod $\underline{M}$ on trouve
$$ \rho = \rho^{\varepsilon_3} \quad \text{dans} \quad \underline{H} $$

ce qui implique $\varepsilon_3 = 1$, OK. Donc
$$ \underline{Z}_{\rho}^\natural = \underline{N}_{\rho,\sigma}^\natural \xrightarrow{\sim} \underline{\gamma}_{\rho,\sigma}' = \underline{O}_{12}^{*\prime} \leqno (163) $$

ainsi $\underline{Z}_{\rho}^\natural \subset \underline{N}_{\rho}$ est une section partielle
au dessus de $\underline{O}_{12}^{*\prime}$, de l'extension
$\underline{N}_{\rho,\sigma}$ de $\underline{O}_{12}^*$ par $\underline{L}_{\rho}$.

[unclear: Prenons] dans cette section partielle un élément $\tau'$ au dessus de $\tau$, qui est une involution
de $\underline{N}_{\rho}$. Est-ce que $\tau'$ normalise
$\underline{Z}_{\rho}^\natural$ ? Comme il normalise $\underline{N}_{\rho}$, il suffirait qu'il normalise $\underline{G}_{\rho,\sigma}^\natural$.
Comme $\tau' = \tau \sigma$, il revient au même de dire
que $\sigma$ normalise $\underline{G}_{\rho,\sigma}^\natural = \underline{O}_{12}^* \underline{G}_{\rho,\sigma}^\natural$.

Comme $\sigma$ normalise $\underline{G}_{\rho}^\circ$ qui est
distingué, il suffit de regarder l'action
de $\sigma$ sur $\underline{O}_{12}^* \subset \underline{G}_{\rho,\sigma}$, défini par les $u_f$ de (151). La question est si
les $\sigma u_f \sigma^{-1} u_f^{-1} \in \underline{\Gamma} \underline{\mathfrak{S}}_{\rho}^\natural$, or

\newpage

%% ===== Page 81 =====

$$ \sigma u_\mu \sigma^{-1} u_\mu^{-1} = \sigma \overbrace{u_\mu(\sigma)^{-1}}^{u_\mu(\sigma)} = \sigma (\omega^{\nu_2} \sigma)^{-1} = \omega^{\nu_2} $$
donc

\underline{Proposition}. Pour une section $x$ de $\underline{G}_{p,\sigma}^\natural$, $\sigma(x)$ est une section de $\underline{G}_{p,\sigma}^\natural \cdot \{1, \omega\}$ et son image dans $\{1, \omega\}$ est égale à $\omega^{\nu_2(x)}$ - donc elle est dans $\underline{G}_{p,\sigma}^\natural$ ssi $f(x)$ est une section de
$$ \underline{\mathbb{D}}_{12}^{* \prime \prime} \defeq \{ f \in \Gamma_{\underline{\mathbb{D}}_{12}^*} \mid \varepsilon_2(f) = 1 \} \leqno (164) $$

\underline{Corollaire 1}\footnote{NB On a $N^{\prime \prime \natural}_\rho = (N^\natural_\rho)^\circ$ comme on le montrera.} $\underline{G}_{p,\sigma}^\natural$ n'est pas normalisé par $\sigma$ ni $\tau'$, mais
$$ \underline{G}_{p,\sigma}^{\prime \prime \natural} \defeq \{ x \in \underline{G}_{p,\sigma}^\natural \mid \varepsilon_2(x) = 1 \} \leqno (165) $$
est égal : \quad $\underline{G}_{p,\sigma}^{\prime \prime \natural} = \underline{G}_{p,\sigma}^\natural \cap \sigma(\underline{G}_{p,\sigma}^\natural) = \underline{G}_{p,\sigma}^\natural \cap \tau'(\underline{G}_{p,\sigma}^\natural)$
et est normalisé par $\sigma$ et $\tau'$.

\underline{Corollaire 2} $\underline{\Pi}_\rho^\natural$ n'est pas normalisé par $\sigma$ ni $\tau'$, mais
$$ \underline{\Pi}_\rho^{\prime \prime \natural} \defeq \underline{\Pi}_\rho^\natural \cap \underline{G}_{p,\sigma}^{\prime \prime \natural} \subset \underline{\Pi}_\rho^\natural \leqno (166) $$
est normalisé par $\sigma$ et $\tau'$ - c'est l'image réciproque dans $\underline{\Pi}_\rho^\natural$ de
$$ \underline{\mathbb{D}}_{12}^{* 0} \quad (= \underline{\mathbb{D}}_{12}^{* \prime} \cap \underline{\mathbb{D}}_{12}^{* \prime \prime}) . \leqno (167) $$

\underline{Corollaire 3}\footnote{A-t-on $Z^\natural_\rho = (N^\natural_\rho)^\circ$ ? Je ferais ainsi évanouir $G_{p,\sigma}^\natural$, $Sp_\rho$, $G_{p,\sigma}$. voir un peu le danger. cf. comp. avec (38) p. 622'.}
Soit
$$ N^{\prime \prime \natural}_\rho \defeq \{ 1, \tau' \} \cdot \underline{\Pi}_\rho^{\prime \prime \natural} \subset N^{\natural}_\rho , \leqno (168) $$
$(= N^{\natural}_\rho \cap (\underline{G}_{p,\sigma}^\natural)^\circ)$
alors
$$ N^{\prime \prime \natural}_\rho / \underline{\Pi}_\rho \isommap \underline{\mathbb{D}}_{12}^{* \prime \prime \prime} = \{ f \in \Gamma_{\underline{\mathbb{D}}_{12}^*} \mid (\varepsilon_2 \varepsilon_3)(f) = 1 \} \leqno (169) $$

\newpage

%% ===== Page 82 =====

Donc $N_\rho^{'''\natural}$ est une section partielle de $N_\rho$
sur $\underline{O}_{12}^{*'''}$. Enfin, posons
$$
Z_\rho^! = Z_\rho^{\circ \natural} \times \mu_\rho \quad (\mu = \{1, -1\}) \leqno{(170)}
$$
alors $Z_\rho^!$ est centralisé par $\tau'$, et
$$
N_\rho^? \defeq Z_\rho^! \cdot N_\rho^{'''\natural} \subset N_\rho \leqno{(171)}
$$
de sorte que $Z_\rho^{\circ \natural}$ devient d'indice $4$ dans
$N_\rho^?$. On a alors

\underline{Corollaire 4.} On a une suite exacte\footnote{NB Il y a [unclear: un lien logique] entre $N_\rho^?$ et le [unclear: groupe d'extension] des pages sur [unclear: la géométrie].}
$$
\begin{tikzcd}
1 \ar[r] & \mu_{S_0} \ar[r] \ar[d, phantom, sloped, "\simeq"] & N_\rho^? \ar[r] & \underline{O}_{12}^* \ar[r] & 1 \\
& \mu_\rho & & \gamma_{\rho,\sigma} \ar[u, phantom] &
\end{tikzcd} \leqno{(172)}
$$
et une bi-section de l'extension
$N_\rho$ de $\underline{O}_{12}^*$ par $L_\rho$.

Il y aurait lieu de développer
le d. gramme des sous-groupes
remarquables de $N_\rho$, calqué sur
le diag. (38) de p. 622.

\newpage

%% ===== Page 83 =====

689

8) Les groupes $N_{\sigma}, N_{\sigma}^S$ ...

Posons
(173) $$ \begin{cases} N_{\sigma} \subset \Gamma_{\rho,\sigma} \\ =^{dfn} \{ u \in \Gamma_{\rho,\sigma} \mid u(\sigma) = \varepsilon_2(u)\sigma \quad (= \omega^{\nu_2(u)} \sigma) \} \end{cases} $$

(174) $$ N_{\sigma}^S =^{dfn} N_{\sigma} \cap G_{\rho,\sigma}^S $$

L'hom. composé
\begin{tikzcd}
N_{\sigma}^S \ar[r, hook] \ar[rr, bend right, "\chi_{\sigma}"'] & S_{\rho,\sigma} \ar[r, "\chi"] & \mathbb{O}_{12}^*
\end{tikzcd}
est épi. En effet, les $u_f$ sont dans $N_{\sigma}^S$ et donnent un scindage $\mathbb{O}_{12}^* \hookrightarrow N_{\sigma}^S$, d'où des décompositions en pro. $1/2$ directs
(175) $$ \begin{cases} N_{\sigma}^S \simeq \mathbb{O}_{12}^* \cdot S Z_{\sigma}^S \\ N_{\sigma} \simeq \mathbb{O}_{12}^* \cdot S Z_{\sigma} \end{cases} $$

où on a posé
(176) $$ \begin{cases} Z_{\sigma} = \operatorname{Ker}(N_{\sigma} \xrightarrow{\varepsilon_2 \circ p} \mu_{S_0}) = \operatorname{Cent}_{G_{\rho,\sigma}}(\sigma) \\ S Z_{\sigma} = Z_{\sigma} \cap S^+ = \operatorname{Cent}_{S^+_{\rho,\sigma}}(\sigma) \\ Z_{\sigma}^S = Z_{\sigma} \cap G_{\rho,\sigma}^S \\ S Z_{\sigma}^S = S Z_{\sigma} \cap G_{\rho,\sigma}^S \\ \quad = S^{\circ} = M \end{cases} $$

D'ailleurs
\begin{tikzcd}
S Z_{\sigma} \ar[r, "p_{S_0}^+"] & \operatorname{Cent}_{H_{S_0}^+}(\sigma) \ar[d, "\wr"'] \\
& L_{\sigma} \simeq (\mathbb{Z}/q\mathbb{Z})_{S_0} \ar[ul, dashed]
\end{tikzcd}
qui a une section, donc

(177) $$ S Z_{\sigma} \simeq L_{\sigma} \times S Z_{\sigma}^S $$

\newpage

%% ===== Page 84 =====

Enfin
$$
\begin{aligned}
S Z_\sigma^\mathcal{G} &= \mathcal{M} \\
&\parallel \\
&\{ x \in \Gamma \mathcal{M} \mid (1 - \sigma_{\mathcal{M}})(x) = 0 \} = \{ \text{[unclear: } a \sigma_0 a^{-1} \sigma_0^{-1} \mid a \in G_a \text{]} \} \\
&= \{ a (\sigma_0 a^{-1} \sigma_0) \mid a \in \Gamma G_a \} \simeq G_a
\end{aligned} \leqno{(178)}
$$

d'où, compte tenu de 177 :
$$
\begin{cases}
S Z_\sigma \simeq L_\sigma \times \underbrace{G_a}_{S Z_\sigma^\mathcal{G}} \\
S Z_\sigma^\mathcal{G} = (S Z_\sigma)^\circ \ (\text{comp. neutre}) \simeq G_a
\end{cases} \leqno{(179)}
$$

On a donc
$$
1 \to \underset{\substack{\parallel \\ S Z_\sigma^\mathcal{G}}}{G_a} \to \mathcal{N}_\sigma^\mathcal{G} \to \underset{\substack{\parallel \\ \gamma_{p,\sigma}}}{\widehat{\mathbf{O}}_{12}^*} \to 1 ,
$$
plus précisément, par les isom. (175) (178) (179)
donnent
$$
\begin{aligned}
\mathcal{N}_\sigma^\mathcal{G} &\simeq \widehat{\mathbf{O}}_{12}^* \cdot G_a \\
\mathcal{N}_\sigma &\simeq \widehat{\mathbf{O}}_{12}^* \cdot (\underline{L}_\sigma \times G_a)
\end{aligned} \leqno{(180)}
$$

Dans ces formules $\widehat{\mathbf{O}}_{12}^*$ opère sur $G_a$ via $\chi : \widehat{\mathbf{O}}_{12}^* \to \widehat{\mathbf{Z}}^*$, et sur $L_\sigma$ via $\varepsilon_2 : \widehat{\mathbf{O}}_{12}^* \to (\pm 1)_{S_0}$ :
$$
\begin{cases}
\mu(\lambda) = \mu.\lambda \\
\mu(\sigma) = \sigma^{\varepsilon_2}
\end{cases} \qquad \text{où } \mu = \chi(\mu), \varepsilon_2 = \varepsilon_2(\mu) \leqno{(181)}
$$

\newpage

%% ===== Page 85 =====

9$^\circ$) Retour sur $\underline{M}_{\rho,\sigma}$.

Nous inspirant des développements du \S\ 49 (notamment VI, 5)), nous poserons

$$
\begin{cases}
S_0 \underline{M}_{\rho,\sigma} = Z_{\rho,\sigma} \times_{\gamma} N_\rho^? \times_{\gamma} \mathcal{N}_\sigma^\natural \times_{\gamma} G_{\rho,\sigma} \\
S_0 \underline{M}_{\rho,\sigma}^\natural = Z_{\rho,\sigma} \times_{\gamma} N_\rho^? \times_{\gamma} \mathcal{N}_\sigma \times_{\gamma} G_{\rho,\sigma}^\natural
\end{cases} \leqno{(182)}
$$
$$
\begin{cases}
\underline{M}_{\rho,\sigma} = N_\rho^? \times_{\gamma} \mathcal{N}_\sigma^\natural \times_{\gamma} G_{\rho,\sigma} \\
\underline{M}_{\rho,\sigma}^\natural = N_\rho^? \times_{\gamma} \mathcal{N}_\sigma \times_{\gamma} G_{\rho,\sigma}^\natural
\end{cases} \leqno{(183)}
$$
$$
\begin{cases}
S_0 \Gamma_{\rho,\sigma} = Z_{\rho,\sigma} \times_{\gamma} N_\rho^? \times_{\gamma} \mathcal{N}_\sigma^\natural \\
\Gamma_{\rho,\sigma} = N_\rho^? \times_{\gamma} \mathcal{N}_\sigma^\natural
\end{cases} \leqno{(184)}
$$

où 
$$
\gamma = \gamma_{\rho,\sigma} = \mathbb{O}_{12}^\ast.
$$
En fait, les groupes $S_0 \underline{M}_{\rho,\sigma}$, $S_0 \underline{M}_{\rho,\sigma}^\natural$, $S_0 \Gamma_{\rho,\sigma}$ jouent plutôt un rôle technique auxiliaire, on a des extensions des groupes $\underline{M}_{\rho,\sigma}$, $\underline{M}_{\rho,\sigma}^\natural$ et $\Gamma_{\rho,\sigma}$ par $L_0^M = \\operatorname{Ker}(Z_{\rho,\sigma} \to \mathbb{O}_{12}^\ast)$,\footnote{toutes déduites de l'extension de $\Gamma_{\rho,\sigma}$ par $L_0^M \simeq \mathbb{G}_a$,} par image inverse via les homomorphismes $\underline{M}_{\rho,\sigma} \to \Gamma_{\rho,\sigma}$ et l'homomorphisme $\underline{M}_{\rho,\sigma}^\natural \to \Gamma_{\rho,\sigma}$.

Nous utiliserons en outre d'autres groupes (plus gros), reliés par une [unclear: série] de suites exactes (comparer p. 83 (179)) :

\newpage

%% ===== Page 86 =====

(190)
\begin{tikzcd}
1 \ar[r] & \underline{C}^\natural \ar[r] \ar[d] & \underline{M}^\natural_{\rho,\sigma} \ar[r] \ar[d] & \Gamma_{\rho,\sigma} \ar[r] \ar[d, equal] & 1 \\
1 \ar[r] & \underline{C}^+ \ar[r] & \underline{M}_{\rho,\sigma} \ar[r] & \Gamma_{\rho,\sigma} \ar[r] & 1
\end{tikzcd}

et des suites exactes

(191)
\begin{tikzcd}
1 \ar[r] & S N^?_P \times S\mathbb{Z}^?_\rho \times \underline{C}^\natural \ar[r] \ar[d] & \underline{M}^\natural_{\rho,\sigma} \ar[r] \ar[d] & \mathbb{O}_{12}^* \ar[r] \ar[d, equal] & 1 \\
1 \ar[r] & S N^?_P \times S\mathbb{Z}^?_\rho \times \underline{C}^+ \ar[r] & \underline{M}_{\rho,\sigma} \ar[r] & \underset{\gamma_{\rho,\sigma}}{\mathbb{O}_{12}^*} \ar[r] & 1
\end{tikzcd}

(192)
$$1 \to S N^?_P \times S\mathbb{Z}^?_\rho \to \Gamma_{\rho,\sigma} \to \underset{\gamma_{\rho,\sigma}}{\mathbb{O}_{12}^*} \to 1$$

où

(193)
$$ \begin{cases} S N^?_P = \mu_P \simeq \{\pm 1\} & (172) \\ S\mathbb{Z}^?_\rho \simeq C_+ & (179) \end{cases} $$

Ces suites exactes se précisent, grâce
aux décompositions $1/2$ directes :

[MARGIN: crossed out lines]
$??? \simeq \mathbb{O}_{12}^* \cdot C_+$ (150, 151)
$??? \simeq \mathbb{O}_{12}^* \cdot C_+$ (180)
$N^?_P / \mu_P \simeq \mathbb{O}_{12}^*$ (172)

(194)
$$ \begin{cases} \mathbb{Z}_{\rho,\sigma} \simeq \mathbb{O}_{12}^* \cdot C_+ & (150, 151) \\ \underline{C}_{\rho,\sigma} \simeq \mathbb{O}_{12}^* \cdot \underline{C}^+ & (15) \\ \underline{C}^\natural_{\rho,\sigma} \simeq \mathbb{O}_{12}^* \cdot \underline{C}^\natural & (160) \\ N^\natural_S \simeq \mathbb{O}_{12}^* \cdot \underline{C}^\natural & (180) \\ \text{et enfin} \\ N^?_P / \mu_P \simeq \mathbb{O}_{12}^* & (172) \end{cases} $$

\newpage

%% ===== Page 87 =====

On trouve alors des isom. canoniques
$$
\begin{cases}
M_{\rho,\sigma} / \mu_p \simeq \mathbb{O}_{12}^\ast \cdot ( G_a \times \underline{\mathfrak{S}}^+ ) \footnote{op. de $\mathbb{O}_{12}^\ast$ sur $G_a \times \underline{\mathfrak{S}}^+$ donné par (181) et (151). Sur $G_a \times \underline{\mathcal{M}}$ s'en déduit [unclear: par l'opération qui ...]} \\
M_{\rho} / \mu_p \simeq \mathbb{O}_{12}^\ast \cdot ( G_a \times \underline{\mathcal{M}} ) \\
\Gamma_{\rho,\sigma} \simeq \mathbb{O}_{12}^\ast \cdot G_a
\end{cases} \leqno (195)
$$
$M_{\rho,\sigma}$ et $M_{\rho}$, on s'en souviendra, sont des extensions de $M_{\rho,\sigma} / \mu_p$ et $M_{\rho} / \mu_p$ par $\mu_p$, se déduisant de l'extension
$$
1 \to \mu_p \to N_p^? \to \mathbb{O}_{12}^\ast \to 1 \leqno (196)
$$
de $\mathbb{O}_{12}^\ast$ par $\mu_p$, par image inverse des hom. canoniques des groupes (195) sur le facteur $\mathbb{O}_{12}^\ast$. Il faudrait revenir sur l'extension (196), qui est sûrement l'image inverse d'une ext. de [unclear: $Sp(\mathbf{F}_p)$] par $H_1$ via
$$
\mathbb{O}_{12}^\ast \to (\operatorname{Aut} \mu_p)_{S_0}
$$
(comparer § 49 V 8°). L'opération de $\mathbb{O}_{12}^\ast$ sur $G_a \times \underline{\mathfrak{S}}^+$ (donc sur $G_a \times \underline{\mathfrak{S}} = G_a \times \underline{\mathcal{M}}$) se fait via des opérations sur $G_a$ et $\underline{\mathfrak{S}}^+$ que j'ai explicitées plus haut (en (181) et (151) !).

\newpage

%% ===== Page 88 =====

10) Les sous-groupes $\underline{M}_{\rho, \sigma}[0]$, $\underline{M}_{\rho, \sigma}(0)$, $\underline{M}_{\rho, \sigma}(1/2)$, $\underline{M}_{\rho, \sigma}'(-J)$ $[\underline{M}_{\rho, \sigma}'''(-J) \supset \underline{M}_{\rho, \sigma}'(-J)]$.

On utilise le plongement

(195) $\underline{M}_{\rho, \sigma} = N_\rho^? \times_\gamma N_\sigma^? \times_\gamma \underline{G}_{\rho, \sigma} \subset \mathcal{N}_\rho^? \times \mathcal{N}_\sigma^? \times \underline{G}_{\rho, \sigma}$
$$ \Big\| $$
$$ \{(a, b, u) \in \mathcal{N}_\rho^? \times \mathcal{N}_\sigma^? \times \underline{G}_{\rho, \sigma} \mid \gamma(a) = \gamma(b) = \underline{\chi}(u)\} $$

On définit donc

(196) $\underline{M}_{\rho, \sigma}[0] \subset \underline{M}_{\rho, \sigma}$
[MARGIN: NB on aura $\underline{M}_{\rho, \sigma}[0] \subset \underline{M}_{\rho, \sigma}^S$]
$$ \| $$
$$ \{(a, b, u) \in \underline{M}_{\rho, \sigma} \mid u \in \underline{Z}_{\rho, \sigma}\} $$

donc
$$ \underline{M}_{\rho, \sigma}[0] \simeq N_\rho^? \times_\gamma N_\sigma^? \times_\gamma \underline{Z}_{\rho, \sigma} $$

(197) Soit $\underline{Z}_{\rho, \sigma}(0) \subset \underline{Z}_{\rho, \sigma}$
le facteur $1/2$ direct $\mathbb{Q}_{12}^*$ de (150) (151),
alors, on a $\underline{M}_{\rho, \sigma}(0) \simeq N_\rho^? \times_\gamma N_\sigma^? \times_\gamma \underline{Z}_{\rho, \sigma}(0)$

[MARGIN: NB $\underline{Z}_{\rho, \sigma}(0) \subset \underline{Z}_{\rho, \sigma} \subset \underline{G}_{\rho, \sigma}$ est en fait le fact. $1/2$ direct $\mathbb{Q}_{12}^*$ de $\underline{G}_{\rho, \sigma}$ ds (150), il proviendrait la notation $\underline{G}_{\rho, \sigma}(0)$]

(198)
\begin{tikzcd}
\underline{M}_{\rho, \sigma}(0) \ar[r, hook] \ar[d, "\downarrow \simeq"'] & \underline{M}_{\rho, \sigma}[0] \ar[d, "\downarrow \simeq"] \\
S \underline{\Gamma}_{\rho, \sigma}(0) \ar[r, hook] \ar[d] & S_0 \underline{\Gamma}_{\rho, \sigma} \\
\underline{\Gamma}_{\rho, \sigma}(0) &
\end{tikzcd}

(199) $\underline{M}_{\rho, \sigma}(0) = \{(a, b, u) \in \underline{M}_{\rho, \sigma} \mid u \in \underline{Z}_{\rho, \sigma}(0)\}$
[MARGIN: dans la descript. (195) de $\underline{M}_{\rho, \sigma}$ comme prod. $1/2$ direct.]
$$ = \underset{\cap \atop \Gamma_Q}{\underline{\mu}} \cdot ( \underset{\cap \atop \Gamma_{\underline{Q}^+}}{\underline{\lambda}} , 0 ) $$

\newpage

%% ===== Page 89 =====

\marginpar{NB $\mu_p \subset \underline{M}_{p,\sigma}(0) \subset \underline{M}_{p,\sigma}[0]$ i.e.}
$$
\underline{M}_{p,\sigma}(0)/\mu_p = \mathbb{Q}_{12}^* \cdot G_a \xhookrightarrow{\text{correspondant à l'inclusion canonique de } G_a \text{ dans } G_a \times \underline{\mathfrak{S}}^+} \underline{M}_{p,\sigma} \simeq \mathbb{Q}_{12}^* (G_a \times \underline{\mathfrak{S}}^+) \leqno{(200)}
$$

D'autre part on a vu
$$
\underline{M}_{p,\sigma}[0]/\mu_p = \mathbb{Q}_{12}^* \cdot (G_a \times L_0^{\underline{M}}) \xhookrightarrow{\underline{M}_{p,\sigma} \simeq} \mathbb{Q}_{12}^* (G_a \times \underline{\mathfrak{S}}^+) \leqno{(201)}
$$

On définit d'autre part
$$
\underline{M}_{p,\sigma}(1/2) = \left\{ (a,b,U) \in \underline{M}_{p,\sigma} \mid U = b \right\} \leqno{(202)}
$$
(donc $\underline{M}_{p,\sigma}(1/2) \subset \underline{M}_{p,\sigma}^\natural$)

Donc on a
$$
\begin{cases}
\underline{M}_{p,\sigma}(1/2) \isommap \underline{\Gamma}_{p,\sigma} \\
\underline{M}_{p,\sigma}^\natural \simeq \underline{M}_{p,\sigma}(1/2) \cdot \underline{M} \\
\underline{M}_{p,\sigma} \simeq \underline{M}_{p,\sigma}(1/2) \cdot \underline{\mathfrak{S}}^+
\end{cases} \leqno{(203)}
$$

Dans la représentation de $\underline{M}_{p,\sigma}/\mu_p$, $\underline{M}_{p,\sigma}^\natural/\mu_p$ en produits semi-directs, on a la description

\marginpar{$\mu_p \subset \underline{M}_{p,\sigma}(1/2)$}
$$
\begin{cases}
\underline{M}_{p,\sigma}(1/2)/\mu_p = \left\{ \mu \cdot (\lambda, \lambda(\lambda_1 + \lambda_2)) \mid \lambda \in \Gamma G_a \subset \Gamma \underline{M} \right\} \\
\qquad \wr \mid \\
\mathbb{Q}_{12}^* \cdot G_a \xhookrightarrow[\text{associé à } G_a \to G_a \times \underline{M}, \, \lambda \mapsto (\lambda, \lambda(\lambda_1 + \lambda_2))]{\underline{M}_{p,\sigma} \simeq} \mathbb{Q}_{12}^* (G_a \times \underline{M})
\end{cases} \leqno{(204)}
$$

\newpage

%% ===== Page 90 =====

Enfin on pose
$$
\begin{matrix}
M_{\rho,\sigma}(-j) \subset M_{\rho,\sigma} \\
\parallel \\
\{ (a,b,U) \in \Gamma_{M_{\rho,\sigma}} \mid U = a \}
\end{matrix} \leqno{(205)}
$$
ici on n'aura pas néc. $U \in \Gamma(\underline{\mathfrak{S}}^\pm)$ (par suite $N_\rho^\natural \not\subset G_{\rho,\sigma}^\natural$, donc $M_{\rho,\sigma}^\natural(-j) \not\subset M_{\rho,\sigma}^\natural$, contrairement à ce qui se passait pour $M_{\rho,\sigma}(0), M_{\rho,\sigma}[0], M_{\rho,\sigma}(\frac{1}{2})$. On aura encore
$$
\begin{cases}
M_{\rho,\sigma}^\natural(-j) \isommap \Gamma_{\rho,\sigma} \\
M_{\rho,\sigma} \simeq M_{\rho,\sigma}(-j) \cdot \underline{\mathfrak{S}}^+
\end{cases} \leqno{(206)}
$$

Mais on a
$$
H_\rho \not\subset M_{\rho,\sigma}(-j)
$$
plus précisément
$$
H_\rho \cap M_{\rho,\sigma}(-j) = \{1\}
$$
donc la relation $M_{\rho,\sigma}(-j)$ ne définit pas (comme $M_{\rho,\sigma}(\frac{1}{2})$) une section de $M_{\rho,\sigma}/H_\rho$ au dessus de $\Gamma_{\rho,\sigma}/H_\rho \simeq N_\sigma^\natural \simeq \mathbf{Q}_{12}^* \cdot G_a$. On a cependant
$$
M_{\rho,\sigma}(-j) \cdot H_\rho = M_{\rho,\sigma}(-j) \cdot [unclear: H_G]
$$
et en posant
$$
\bar{M}_{\rho,\sigma}(-j) = M_{\rho,\sigma}(-j) \cdot H_\rho / H_\rho = \operatorname{Im}(M_{\rho,\sigma}(-j) \to M_{\rho,\sigma}/H_\rho) \leqno{(207)}
$$
$$
\bar{\Gamma}_{\rho,\sigma} = \Gamma_{\rho,\sigma}/H_\rho \simeq (M_{\rho,\sigma}/H_\rho)/\underline{\mathfrak{S}}^+ \quad [unclear: canonique]
$$

\newpage

%% ===== Page 91 =====

\footnote{NB On a $\underline{\Gamma}_{p,\sigma}^\sim \to \underline{N}_\sigma^\natural/\mu_p$ épi de noyau $\underline{\Gamma}_{p,\sigma}^{\prime \sim}$ ($\underline{N}_\sigma^\sim = \text{[unclear: Diag]} \dots$) et donc le diag. comm.
\begin{tikzcd}
1 \ar[r] & \underline{\Gamma}_{p,\sigma}^{\prime\sim} \ar[r] & \underline{\Gamma}_{p,\sigma}^\sim \ar[r] & \underline{N}_\sigma^\natural / \mu_p \ar[r] & 1 \\
& \underline{M}_{p,\sigma}^{\prime\sim}(-j) \ar[u, "\simeq"] \ar[r] & \underline{M}_{p,\sigma}^{\sim}(-j) \ar[u]
\end{tikzcd}
Le diag. en pointillé se déduit de lui par [unclear: intersection avec] $\underline{M}_{p,\sigma}^{\sim}(-j) \dots$}
$$ [unclear: contient] \quad 1 \to \dots \to \underline{M}_{p,\sigma}^\sim(-j) \to \underline{\Gamma}_{p,\sigma}^\sim \to 1 \leqno{(208)} $$

i.e. $\underline{M}_{p,\sigma}^\sim(-1/2)$ est une bisection de $\underline{M}_{p,\sigma}^\sim = M_{p,\sigma} / \mu_p$ au dessus de $\underline{\Gamma}_{p,\sigma}^\sim$. Posant
$$ \underline{M}_{p,\sigma}^{\prime \sim}(-j) \defeq \underline{M}_{p,\sigma}^{\sim}(-j) \cap \underbrace{\underline{M}_{p,\sigma}^{\prime \sim}}_{\substack{|| \text{def} \\ M_{p,\sigma}^\prime / \mu_p}} \leqno{(209)} $$

on a
$$ \underline{M}_{p,\sigma}^{\prime \sim}(-j) \isommap \underline{\Gamma}_{p,\sigma}^{\prime \sim} \defeq \\operatorname{Ker}(\underline{\Gamma}_{p,\sigma}^\sim \xrightarrow{\varepsilon_2} \text{[unclear: $\mu_2$]}) $$

$\underline{M}_{p,\sigma}^{\prime \sim}(-j)$ d'indice 4 dans $\underline{M}_{p,\sigma}^\sim(-j)$, section partielle de $\underline{M}_{p,\sigma}^\sim$ au dessus de $\underline{\Gamma}_{p,\sigma}^{\prime \sim}$, d'image [unclear: $\underline{\Gamma}_{p,\sigma}^{\prime \sim}$] d'indice 2 dans $\underline{\Gamma}_{p,\sigma}^\sim$. On aura donc

$$
\begin{cases}
\underline{M}_{p,\sigma}^{\prime \sim} \defeq \\operatorname{Ker}(\underline{M}_{p,\sigma}^{\sim} \xrightarrow{\varepsilon_3} \text{[unclear: $\mu_2$]}) \simeq \underline{M}_{p,\sigma}^{\prime \sim}(-j) \cdot \mathfrak{S}^+ \\
\underline{M}_{p,\sigma}^{\sim} \defeq \\operatorname{Ker}(\underline{M}_{p,\sigma}^{\sim} \xrightarrow{\varepsilon_4} \text{[unclear: $\mu_2$]}) \simeq \underline{M}_{p,\sigma}^{\sim}(-j) \cdot \mathfrak{S}
\end{cases} \leqno{(210)}
$$

Remarque. Je m'arrête là dans les dévissages de sous-groupes [unclear: intérieurs] de $\underline{M}_{p,\sigma}^\sim(-j)$ -- j'y reviendrai sans doute plus tard, dans le contexte des schémas en groupes $\underline{M}_{p,\sigma}$ généraux. Il devient clair qu'il faudra clarifier,

\newpage

%% ===== Page 92 =====

pour la suite, entre les notations possibles respectives pour les groupes notés ici $\underline{M}_{p,\sigma}$ et $\underline{M}_{p,\sigma}/\mu_p \simeq \underline{M}_{p,\sigma}^\sim$ (et idem sans $\sigma$) ce sont ces derniers, il est vrai, qui me semblent les plus intéressants, car c'est eux\footnote{Plutôt, le gr. pro-fini [unclear: que je regrette...]} qui reçoivent le groupe universel 
$$
\mathcal{M} = \mathcal{N}_{1,1}^\sim \quad \text{(extension de } \Gamma_{\mathbf{Q}}^\sim \text{ (} =\mathcal{G}_{\mathbf{Q}}^\sim \text{) par } \\operatorname{Sl}(2, \mathbf{Z})^\wedge \text{).}
$$
Il semblerait que ces groupes mériteraient la notation la plus simple, en l'occurrence $\underline{M}_{p,\sigma}$, et les groupes notés précédemment $\underline{M}_{p,\sigma}$ prendraient des $\sim$, ou un index analogue. (Le $\sim$ est déjà utilisé dans un sens un peu différent dans $\Gamma_{\mathbf{Q}}^\sim$, $M_{0,3}^\sim$ etc — d'autre part les exposants [unclear: biffés], $^\circ, ^+, ^{1/2}, ^\sim$ déjà utilisés dans des sens différents. Le double emploi du $\sim$ risque d'ailleurs d'introduire des conflits de notations, tôt ou tard ; il est peut-être préférable de remplacer le $\sim$ dans $\Gamma_{\mathbf{Q}}^\sim$ par des lettres telles que $\Gamma_{\mathbf{Q}}^{\text{ét}}$ (suggère "étendu") ou $\Gamma_{\mathbf{Q}}^{\text{gén}}$ "généralisé" gén... Il faudra que je résolve ces di-

\newpage

%% ===== Page 93 =====

dilemmes notationnels, pour l'instant je
garde les notations introduites sur la
bande dans la présente section, qui
seront révisées avec soin dans un $\S$ spécial
ultérieur.

Il faudrait encore décrire des
sous-groupes $\underline{M}_{p,\sigma}^\sim (-J)$, $\underline{M}_{p,\sigma}'^\sim (-J)$
dans $\underline{M}_{p,\sigma}^\sim$ décrit comme produit
1/2 direct, trouvant
$$ 
\underline{M}_{p,\sigma}^\sim (-J) \subset \underline{M}_{p,\sigma}^\sim \simeq Q_1^+ ( \Gamma_{\mathbf{Q}} \times \underline{\mathfrak{S}}_{p,\sigma}^+ ) \leqno{(211)} 
$$
$$ 
= \left\{ (\underset{\in \Gamma_{\mathbf{Q}}}{(\lambda, x)}, \underset{\in Q_1^+}{u}) \mid u_\mu \in \underline{G}_{p,\sigma} \right\} 
$$

où $u_\mu$ est défini dans (151), en particulier
$$ 
u_\mu(\rho) = \lambda_1^{-1} u(\rho) \quad , \quad \nu = \nu(\mu) , \quad \text{avec (154)} 
$$

\marginpar{\delta = \lambda_1 \lambda_0^{-1}}
$$ 
u = \sigma \nu_3 \delta^{-\nu} u_\mu \in \Gamma_{N_p} \quad , \quad \chi(u) = \mu 
$$
et posant $\sigma = \omega \tau ' \varepsilon$, la première
s'écrit aussi
$$ 
u = \varepsilon \nu_3^{1/3} \big( \underbrace{\omega^{1/3} \cdot \tau^{1/3} \delta^{-\nu} u_\mu}_{\in \Gamma_{p,\sigma} \text{ car les trois facteurs de la droite sont } \in \Gamma_{N_p} \cap \underline{G}_{p,\sigma} = \Gamma_{\mathbf{Z}_p}^S} \big) \leqno{(212)} 
$$

\newpage

%% ===== Page 94 =====

%% [page 94 missing or failed: 429 RESOURCE_EXHAUSTED. {'error': {'code': 429, 'message': 'You exceeded your current quota, please check your plan and billing details. For more information on this error, head to: https://ai.google.dev/gemini-api/docs/rate-limits. To monitor your current usage, head to: https://ai.dev/rate-limit. \n* Quota exceeded for metric: generativelanguage.googleapis.com/generate_requests_per_model_per_day, limit: 250, model: gemini-3.1-pro\nPlease retry in 5h18m56.010577551s.', 'status': 'RESOURCE_EXHAUSTED', 'details': [{'@type': 'type.googleapis.com/google.rpc.Help', 'links': [{'description': 'Learn more about Gemini API quotas', 'url': 'https://ai.google.dev/gemini-api/docs/rate-limits'}]}, {'@type': 'type.googleapis.com/google.rpc.QuotaFailure', 'violations': [{'quotaMetric': 'generativelanguage.googleapis.com/generate_requests_per_model_per_day', 'quotaId': 'GenerateRequestsPerDayPerProjectPerModel', 'quotaDimensions': {'location': 'global', 'model': 'gemini-3.1-pro'}, 'quotaValue': '250'}]}, {'@type': 'type.googleapis.com/google.rpc.RetryInfo', 'retryDelay': '19136s'}]}}]

\newpage

%% ===== Page 95 =====

11) L'homomorphisme $M \to M_{\rho,\sigma}(\hat{\mathbb{Z}})^{\sim}$, relatif à la tour des courbes de Fermat.

Notons qu'on a pour tout entier $N \in \mathbb{N}^*$

[MARGIN: Le $\Delta$ indique qu'on prend (215) sections de [unclear] qui sont [unclear]]

$$ \mathbb{O}_N^{\Delta}(\hat{\mathbb{Z}}) \simeq \mathbb{Z} \quad \text{d'où} \quad \mathbb{O}_\Gamma^{\Delta}(\hat{\mathbb{Z}})^{\wedge} \simeq \hat{\mathbb{Z}}^* $$

On a donc, compte tenu des décompositions en produits semi-directs (150) et des isomorphismes

(216)
$$ Z_{\rho,\sigma}^\sim \overset{df}{=} Z_{\rho,\sigma}^\Delta(\hat{\mathbb{Z}})^\wedge \simeq \hat{\mathbb{Z}}^* . \hat{L}_0 \quad \text{où} \quad \hat{L}_0 \simeq \hat{\mathbb{Z}} . \lambda_0 \simeq \hat{\mathbb{Z}} $$
$$ G_{\rho,\sigma}^\sim \overset{df}{=} G_{\rho,\sigma}^\Delta(\hat{\mathbb{Z}})^\wedge \simeq \hat{\mathbb{Z}}^* . \underline{G}^{+\wedge} $$
$$ C_{\rho,\sigma}^\sim \overset{df}{=} C_{\rho,\sigma}^\Delta(\hat{\mathbb{Z}})^\wedge \simeq \hat{\mathbb{Z}}^* . M^\wedge $$
$$ N_{\rho}^\sim \overset{df}{=} N_\rho^\sim(\hat{\mathbb{Z}})^\wedge \simeq \mathbb{O}_\Gamma^\Delta(\hat{\mathbb{Z}})^\wedge \simeq \hat{\mathbb{Z}}^* $$
$$ N_{\sigma}^\sim \overset{df}{=} N_{\sigma}^\sim(\hat{\mathbb{Z}})^\wedge \simeq \hat{\mathbb{Z}}^* . \hat{\mathbb{Z}} $$
$$ M_{\rho,\sigma}^\sim \overset{df}{=} M_{\rho,\sigma}^\sim(\hat{\mathbb{Z}})^\wedge \simeq \hat{\mathbb{Z}}^* . (\hat{\mathbb{Z}} \times \underline{G}^{+\wedge}) $$
$$ M_{\bar{\rho},\bar{\sigma}}^\sim \overset{df}{=} M_{\bar{\rho},\bar{\sigma}}^\sim(\hat{\mathbb{Z}})^\wedge \simeq \hat{\mathbb{Z}}^* . (\hat{\mathbb{Z}} \times \hat{M}) $$

où les opérations de $\hat{\mathbb{Z}}^*$ sur les $\hat{\mathbb{Z}}$ facteurs à droite sont explicitées dans les sections précédentes. On notera que, modulo des ajustages de notations et de définitions, les groupes $Z_{\rho,\sigma}^\sim$ etc. précédents sont ceux qui serais construits directement, en termes de $(\hat{\underline{G}}^\wedge, \underline{G}^{+\wedge}, M^\wedge)$ (non indexé de $\rho,\sigma,\tau$), dans §49 VI.

On a des hom. canoniques

(217)
\begin{tikzcd}
M = [unclear] \ar[r] & M_{\rho,\sigma}^\sim \\
M' \ar[r] \ar[u, hook] & M_{\bar{\rho},\bar{\sigma}}^\sim \ar[u, hook]
\end{tikzcd}

qui comme d'habitude se décompose en

\newpage

%% ===== Page 96 =====

%% [page 96 missing or failed: 429 RESOURCE_EXHAUSTED. {'error': {'code': 429, 'message': 'You exceeded your current quota, please check your plan and billing details. For more information on this error, head to: https://ai.google.dev/gemini-api/docs/rate-limits. To monitor your current usage, head to: https://ai.dev/rate-limit. \n* Quota exceeded for metric: generativelanguage.googleapis.com/generate_requests_per_model_per_day, limit: 250, model: gemini-3.1-pro\nPlease retry in 5h17m5.884393717s.', 'status': 'RESOURCE_EXHAUSTED', 'details': [{'@type': 'type.googleapis.com/google.rpc.Help', 'links': [{'description': 'Learn more about Gemini API quotas', 'url': 'https://ai.google.dev/gemini-api/docs/rate-limits'}]}, {'@type': 'type.googleapis.com/google.rpc.QuotaFailure', 'violations': [{'quotaMetric': 'generativelanguage.googleapis.com/generate_requests_per_model_per_day', 'quotaId': 'GenerateRequestsPerDayPerProjectPerModel', 'quotaDimensions': {'location': 'global', 'model': 'gemini-3.1-pro'}, 'quotaValue': '250'}]}, {'@type': 'type.googleapis.com/google.rpc.RetryInfo', 'retryDelay': '19025s'}]}}]

\newpage

%% ===== Page 97 =====

J'ai envie de mettre ces homomorphismes en
relation avec la "tour" des courbes algébriques
définies sur $\mathbb{Q}$, indexées par $n \in \mathbb{N}^\ast$

(220) $X_n : X_0^n + X_1^n + X_\infty^n = 0$ \qquad \qquad (courbe de
$\subset \mathbb{P}^2_{\mathbb{Q}}$ \qquad \qquad \qquad \qquad Fermat d'indice $n$)

Soit $U_n$ l'ouvert des pts où $X_0 X_1 X_\infty \neq 0$

Sur $X_n$ le groupe fini $\mu_n^{\{0,1,\infty\}} / \mu_n$ (diag.)
opère ($\mu_n \subset \mathbb{G}_m$), et on a aussi
que $\mathfrak{S}_3 = \mathfrak{S}_{\{0,1,\infty\}}$, de sorte que le
produit semi-direct

[MARGIN: Notation $F_n$ pas très heureuse - il est possible qu'il y aura un besoin ultérieur de la notation $F_n$ pour les [unclear] de $U_n, U_\infty$ avec un in-dice $n$? (comparer p. 719)]

(221) $F_n = \mathfrak{S}_3 \cdot \mu_n^{\{0,1,\infty\}} / \mu_n = F_n^0$

opère sur $X_n$.
On a

(222) $X_n / F_n^0 \simeq X_1 \simeq U_{0,3}$ ,

isomorphismes compatibles aux opérations de
$\mathfrak{S}_3$. Plus généralement, si
$n = n'd$, on a

(223)
\begin{tikzcd}
X_n \ar[r] & X_{n'} \\
x_i \ar[r, |->] & x_i^d
\end{tikzcd}

et de cette façon $X_n$ devient un revête-
ment principal (= torseur) sur $X_{n'}$, de
groupe $\mu_d^{\{0,1,\infty\}} / \mu_d$, de façon plus précise,
[unclear], (223) est compatible avec
l'épimorphisme

(224)
\begin{tikzcd}
F_n \ar[r, "\text{épi}"] \ar[d, equal] & F_{n'} \ar[d, equal] & g \ar[r, |->] & g \\
\mathfrak{S}_3 \cdot \mu_n^{\{0,1,\infty\}}/\mu_n \ar[r] & \mathfrak{S}_3 \cdot \mu_{n'}^{\{0,1,\infty\}}/\mu_{n'} & (x_0, x_1, x_\infty) \ar[r, |->] & (x_0^d, x_1^d, x_\infty^d)
\end{tikzcd}

\newpage

%% ===== Page 98 =====

%% [page 98 missing or failed: 429 RESOURCE_EXHAUSTED. {'error': {'code': 429, 'message': 'You exceeded your current quota, please check your plan and billing details. For more information on this error, head to: https://ai.google.dev/gemini-api/docs/rate-limits. To monitor your current usage, head to: https://ai.dev/rate-limit. \n* Quota exceeded for metric: generativelanguage.googleapis.com/generate_requests_per_model_per_day, limit: 250, model: gemini-3.1-pro\nPlease retry in 5h9m59.902769437s.', 'status': 'RESOURCE_EXHAUSTED', 'details': [{'@type': 'type.googleapis.com/google.rpc.Help', 'links': [{'description': 'Learn more about Gemini API quotas', 'url': 'https://ai.google.dev/gemini-api/docs/rate-limits'}]}, {'@type': 'type.googleapis.com/google.rpc.QuotaFailure', 'violations': [{'quotaMetric': 'generativelanguage.googleapis.com/generate_requests_per_model_per_day', 'quotaId': 'GenerateRequestsPerDayPerProjectPerModel', 'quotaDimensions': {'location': 'global', 'model': 'gemini-3.1-pro'}, 'quotaValue': '250'}]}, {'@type': 'type.googleapis.com/google.rpc.RetryInfo', 'retryDelay': '18599s'}]}}]

\newpage

%% ===== Page 99 =====

%% [page 99 missing or failed: 429 RESOURCE_EXHAUSTED. {'error': {'code': 429, 'message': 'You exceeded your current quota, please check your plan and billing details. For more information on this error, head to: https://ai.google.dev/gemini-api/docs/rate-limits. To monitor your current usage, head to: https://ai.dev/rate-limit. \n* Quota exceeded for metric: generativelanguage.googleapis.com/generate_requests_per_model_per_day, limit: 250, model: gemini-3.1-pro\nPlease retry in 5h2m9.39707576s.', 'status': 'RESOURCE_EXHAUSTED', 'details': [{'@type': 'type.googleapis.com/google.rpc.Help', 'links': [{'description': 'Learn more about Gemini API quotas', 'url': 'https://ai.google.dev/gemini-api/docs/rate-limits'}]}, {'@type': 'type.googleapis.com/google.rpc.QuotaFailure', 'violations': [{'quotaMetric': 'generativelanguage.googleapis.com/generate_requests_per_model_per_day', 'quotaId': 'GenerateRequestsPerDayPerProjectPerModel', 'quotaDimensions': {'location': 'global', 'model': 'gemini-3.1-pro'}, 'quotaValue': '250'}]}, {'@type': 'type.googleapis.com/google.rpc.RetryInfo', 'retryDelay': '18129s'}]}}]

\newpage

%% ===== Page 100 =====

%% [page 100 missing or failed: 429 RESOURCE_EXHAUSTED. {'error': {'code': 429, 'message': 'You exceeded your current quota, please check your plan and billing details. For more information on this error, head to: https://ai.google.dev/gemini-api/docs/rate-limits. To monitor your current usage, head to: https://ai.dev/rate-limit. \n* Quota exceeded for metric: generativelanguage.googleapis.com/generate_requests_per_model_per_day, limit: 250, model: gemini-3.1-pro\nPlease retry in 4h51m57.190312094s.', 'status': 'RESOURCE_EXHAUSTED', 'details': [{'@type': 'type.googleapis.com/google.rpc.Help', 'links': [{'description': 'Learn more about Gemini API quotas', 'url': 'https://ai.google.dev/gemini-api/docs/rate-limits'}]}, {'@type': 'type.googleapis.com/google.rpc.QuotaFailure', 'violations': [{'quotaMetric': 'generativelanguage.googleapis.com/generate_requests_per_model_per_day', 'quotaId': 'GenerateRequestsPerDayPerProjectPerModel', 'quotaDimensions': {'model': 'gemini-3.1-pro', 'location': 'global'}, 'quotaValue': '250'}]}, {'@type': 'type.googleapis.com/google.rpc.RetryInfo', 'retryDelay': '17517s'}]}}]

\newpage

%% ===== Page 101 =====

698

$$ (234) \qquad 1 \to \Pi_1(U_{0,3 \overline{\mathbb{Q}}}, \overline{y}) \to \Pi_1(U_{0,3 \mathbb{Q}}, \mathfrak{S}_3; \overline{y}) \to \mathfrak{S}_3 \to 1 $$

par le sous-groupe invariant des commutateurs
du groupe $\Pi_1(U_{0,3 \overline{\mathbb{Q}}}, \overline{y})$. Cette structure
[MARGIN: si, elle est triviale, cf. plus bas...]
d'extension (233) de $\mathfrak{S}_3$ par $M^\wedge$ n'est pas
triviale, sauf erreur, bien que l'extension
induite de $\{1, \rho, \rho^2\} \simeq \mathbb{Z}/3\mathbb{Z}$ par $M^\wedge$, - ou
$\{1, \sigma\}$ par $\hat{M}$, le soit. A fortiori (232) n'est
pas triviale (bien qu'elle ait des scindages
canoniques au-dessus de $\hat{\mathbb{Z}}$, au-dessus de
$\{1, \rho, \rho^2\}$ et au-dessus de $\{1, \sigma\}$). Donc
(sous réserve de vérifier ces points tantot...)
le groupe $C_{\rho,\sigma}/\mu_a$ ne peut s'identifier
à un sous-groupe du groupe 
$$ \boxed{ Aff_{\hat{\mathbb{Z}}}(\hat{M}) \simeq \operatorname{Aut}_{\hat{\mathbb{Z}}}(\hat{M}). \hat{M} } $$
comprenant le sous-groupe $\hat{\mathbb{Z}}^\times \times \mathfrak{S}_3$ de $\operatorname{Aut}_{\hat{\mathbb{Z}}}(\hat{M})$.

Considérons le sous-groupe
$$ (235) \qquad Aff \simeq (\hat{\mathbb{Z}}^\times \times \mathfrak{S}_3) . \hat{M} $$
On a un diagr.

\begin{tikzcd}
(237) & 1 \ar[r] & \hat{M} \ar[r] & Aff \ar[r] & \hat{\mathbb{Z}}^\times \times \mathfrak{S}_3 \ar[r] & 1 \\
& & & & \mathcal{M}_{0,3} \ar[u]
\end{tikzcd}

: ligne exacte. Le donné de la structure
(230) équivaut à isoler (à un conjugué près, à l' [unclear]

\newpage

%% ===== Page 102 =====

%% [page 102 missing or failed: 429 RESOURCE_EXHAUSTED. {'error': {'code': 429, 'message': 'You exceeded your current quota, please check your plan and billing details. For more information on this error, head to: https://ai.google.dev/gemini-api/docs/rate-limits. To monitor your current usage, head to: https://ai.dev/rate-limit. \n* Quota exceeded for metric: generativelanguage.googleapis.com/generate_requests_per_model_per_day, limit: 250, model: gemini-3.1-pro\nPlease retry in 4h31m51.835343072s.', 'status': 'RESOURCE_EXHAUSTED', 'details': [{'@type': 'type.googleapis.com/google.rpc.Help', 'links': [{'description': 'Learn more about Gemini API quotas', 'url': 'https://ai.google.dev/gemini-api/docs/rate-limits'}]}, {'@type': 'type.googleapis.com/google.rpc.QuotaFailure', 'violations': [{'quotaMetric': 'generativelanguage.googleapis.com/generate_requests_per_model_per_day', 'quotaId': 'GenerateRequestsPerDayPerProjectPerModel', 'quotaDimensions': {'location': 'global', 'model': 'gemini-3.1-pro'}, 'quotaValue': '250'}]}, {'@type': 'type.googleapis.com/google.rpc.RetryInfo', 'retryDelay': '16311s'}]}}]

\newpage

%% ===== Page 103 =====

699

(239)
\begin{tikzcd}
1 \ar[r] & M^\wedge \ar[r] & G_{p\Gamma}/H_G \ar[r] & \hat{\mathbb{Z}}^* \times \mathfrak{S}_3 \ar[r] & 1 \\
& & \hat{\mathbb{Z}}^* \times G^{+ \wedge} / H_G \ar[u, "\simeq"'] & & \\
& & & \mathcal{M}_{0,3} = \pi_1(\mathcal{U}_{0,3 \overline{\mathbb{Q}}, \mathfrak{S}_3} ; -\vec{1}) \ar[uu] \ar[ul, dashed]
\end{tikzcd}

Sur $G_{0,3}^{+ \wedge} \simeq \pi_1(\mathcal{U}_{0,3 \overline{\mathbb{Q}}, \mathfrak{S}_3}, -\vec{1}) \subset \mathcal{M}_{0,3}$ , cet hom. induit l'hom. canonique du passage au quotient de $G_{0,3}^{+ \wedge}$ obtenu en rendant $\Pi_{0,3}^\wedge$ abélien , soit $G_{p\Gamma}/H_G \simeq G^{+ \wedge}/H_G$ , extension de $\mathfrak{S}_3$ par $M^\wedge$.

Reprenons l'hom. $\mathcal{M}_{0,3} \to \text{Aff}$ (238) , il s'insère dans un diag. commutatif

(240)
\begin{tikzcd}
1 \ar[r] & \Pi_{0,3}^\wedge \ar[r] \ar[d] & \mathcal{M}_{0,3} \ar[r] \ar[d] & \Gamma_{\mathbb{Q}} \times \mathfrak{S}_3 \ar[r] \ar[d, "\chi \times \text{id}"] & 1 \\
1 \ar[r] & \hat{M} \ar[r] & \text{Aff} \ar[r] & \hat{\mathbb{Z}}^* \times \mathfrak{S}_3 \ar[r] & 1
\end{tikzcd}

où $\Gamma_{\mathbb{Q}} \times \mathfrak{S}_3 \to \hat{\mathbb{Z}}^* \times \mathfrak{S}_3$ est l'hom. de (237) , et
où $\Pi_{0,3}^\wedge \to \hat{M}$ est l'hom. canonique

(241) $$ \Pi_{0,3}^\wedge \to (\Pi_{0,3}^\wedge)_{ab} \simeq ((\Pi_{0,3})_{ab})^\wedge \simeq \hat{M} \ . $$

Si on désigne par $\mathcal{M}_1$ le quotient de $\mathcal{M}_{0,3}$ par $[\Pi_{0,3}^\wedge, \Pi_{0,3}^\wedge]$ , on trouve de même

\newpage

%% ===== Page 104 =====

un hom. canonique de suites exactes

(242)
\begin{tikzcd}
1 \ar[r] & M^\wedge \ar[r] \ar[d, equal] & M_1 \ar[r] \ar[d] & N_{0,3} \ar[r] \ar[d] & 1 \\
1 \ar[r] & M^\wedge \ar[r] & Aff \ar[r] & \hat{\mathbb{Z}}^\times \times \mathfrak{S}_3 \ar[r] & 1
\end{tikzcd}

(dépends du choix d'un pt de $\mathcal{U}_{0,3}(\overline{\mathbb{Q}})$ au-dessus de $-\vec{j}$)
qui montre que l'extension $M_1$ de $\hat{\Gamma}_Q$ par $M^\wedge$
est [unclear: isome] : l'image inv. de l'extension
triviale $Aff$ de $\hat{\mathbb{Z}}^\times \times \mathfrak{S}_3$ par $M^\wedge$, donc est
triviale. Plus précisément encore, considérons
le quotient $M_2$ de $M_1$ par le noyau de $M_1 \to Aff$,
on trouve un diagramme

(242 bis)
\begin{tikzcd}
1 \ar[r] & M^\wedge \ar[r] \ar[d, "\simeq"'] & M_2 \ar[r] \ar[d, hook] & N_2 \ar[r] \ar[d, hook] & 1 \\
1 \ar[r] & M^\wedge \ar[r] & Aff \ar[r] & \hat{\mathbb{Z}}^\times \times \mathfrak{S}_3 \ar[r] & 1
\end{tikzcd}

(où $N_2 \overset{def}{=} M_2 / \text{Im } M^\wedge$). On voit donc
que l'extension de $N_2$ par $M^\wedge$ est triviale
- [unclear: avec un choix de] $M^\wedge$ - conjugaison -
de scindages de cette extension - les
scindages correspondant aux pts de
$\mathcal{U}_{0,3}(\overline{\mathbb{Q}})$ au-dessus du pt $-\vec{j} \in \mathcal{U}_{0,3}(\mathbb{Q})$

\newpage

%% ===== Page 105 =====

On aurait vraiment envie de montrer que
cet homomorphisme $M_{0,3} \to Aff$ est isomorphe
à l'hom. $M_{0,3} \to G_{\rho,0}/\mu_{\hat{\mathbb{Z}}}$ de (239) , moyennant
un isomorphisme canonique de
$Aff$ avec $G_{\rho,0}/\mu_{\hat{\mathbb{Z}}}$. Cela impliquerait que
cette dernière extension est scindée. Les
isomorphismes de cette extension avec $Aff$
correspondraient aux scindages de cette
extension.

12°) Scindages de l'extension $G_{\rho,0}/\mu_{\hat{\mathbb{Z}}}$ de $\hat{\mathbb{Z}}^* \times \mathfrak{S}_3$ par $M$.
Comme on a déjà un relèvement
$$ \hat{\mathbb{Z}}^* \hookrightarrow G_{\rho,0} $$
(150, 151), je vais regarder les scindages
de l'extension (231) qui prolongent celui-ci -
ce qui revient à, les
scindages de l'extension induite sur $\mathfrak{S}_3$ :

[MARGIN: partout : il faut se donner ce scindage qui $\mathfrak{S}_3 \hookrightarrow \underline{S}^+ / \mu_{\hat{\mathbb{Z}}}$ ($G_{\rho,0} \simeq \hat{\mathbb{Z}}^* \dots$)]
[MARGIN: On regarde les scindages [unclear] correspondent bien uniquement aux repères]

(241) $$ 1 \to M \to \underline{S}^+ / \mu_{\hat{\mathbb{Z}}} \to (\mathfrak{S}_3)_{S_0} \to 1 $$

\newpage

%% ===== Page 106 =====

%% [page 106 missing or failed: no entry]

\newpage

%% ===== Page 107 =====

%% [page 107 missing or failed: no entry]

\newpage

%% ===== Page 108 =====

%% [page 108 missing or failed: no entry]

\newpage

%% ===== Page 109 =====

%% [page 109 missing or failed: no entry]

\newpage

%% ===== Page 110 =====

%% [page 110 missing or failed: no entry]

\newpage

%% ===== Page 111 =====

%% [page 111 missing or failed: no entry]

\newpage

%% ===== Page 112 =====

%% [page 112 missing or failed: no entry]

\newpage

%% ===== Page 113 =====

%% [page 113 missing or failed: no entry]

\newpage

%% ===== Page 114 =====

%% [page 114 missing or failed: no entry]

\newpage

%% ===== Page 115 =====

%% [page 115 missing or failed: no entry]

\newpage

%% ===== Page 116 =====

%% [page 116 missing or failed: no entry]

\newpage

%% ===== Page 117 =====

%% [page 117 missing or failed: no entry]

\newpage

%% ===== Page 118 =====

%% [page 118 missing or failed: no entry]

\newpage

%% ===== Page 119 =====

%% [page 119 missing or failed: no entry]

\newpage

%% ===== Page 120 =====

%% [page 120 missing or failed: no entry]

\newpage

%% ===== Page 121 =====

70bis

13°) Retour sur la tour de Fermat.

Commençons par dégager une suite topologique générale, qui était un peu embrouillée entre les lignes dans 110).

Soit $\Gamma$ (ici $\Gamma_{\mathbb{Q}} \times \mathfrak{S}_3 = \mathcal{N}_{0,3}$) opérant sur un topos $\mathcal{X}$ (ici $\underline{M}_{0,3} \simeq M_{0,4 \overline{\mathbb{Q}}}$) conn. loc. connexe, muni d'un "pt" (géométrique) $x_0$ (ici : le point $\vec{1}$ à valeurs dans $\mathbb{Q}$), soient

[MARGIN: en fait on discute [unclear: localement] ($\mathcal{X}$ loc. simplemt [unclear: connexe] ou profini, il faudra y revenir...)]

$$ (278) \begin{cases} \Pi = \Pi_1(\mathcal{X}, x_0), \mathcal{M} = \Pi_1(\mathcal{X}, \Gamma; x_0) \\ M = \Pi_{ab} = \Pi/[\Pi,\Pi] \quad, \tilde{\mathcal{M}} = \mathcal{M}/[\Pi,\Pi] \end{cases} $$

d'où des suites exactes et hom. de suites exactes

(279)
\begin{tikzcd}
1 \ar[r] & \Pi \ar[r] \ar[d, "epi"'] & \mathcal{M} \ar[r] \ar[d, "epi"'] & \Gamma \ar[r] \ar[d, "\sim"'] & 1 \\
1 \ar[r] & M \ar[r] & \tilde{\mathcal{M}} \ar[r] & \Gamma \ar[r] & 1
\end{tikzcd}

(Au lieu de prendre $M = \Pi/[\Pi,\Pi] = \Pi_{ab}$ on aurait pu prendre un quotient de $\Pi_{ab}$ invariant par action de $\Gamma$, mais peu importe...). Considérons "le rev. universel abélien" $\tilde{\mathcal{X}}_{ab}$ de $\mathcal{X}$ en $x_0$, i.e. le quotient par $[\Pi,\Pi]$ du rev. universel de $\mathcal{X}$ en $x_0$ - c'est donc un rev. principal de $\mathcal{X}$, de groupe $M$. On

\newpage

%% ===== Page 122 =====

%% [page 122 missing or failed: no entry]

\newpage

%% ===== Page 123 =====

%% [page 123 missing or failed: no entry]

\newpage

%% ===== Page 124 =====

Donc si on introduit le groupe $Aff(E)$
on aura un homomorphisme de
suites exactes

(281)
\begin{tikzcd}
1 \ar[r] & M \ar[r] \ar[d, "id_M"] & \tilde{\Gamma} \ar[r] \ar[d, "\theta_{\tilde{\Gamma}}"] & \Gamma \ar[r] \ar[d, "\Theta_\Gamma"] & 1 \\
1 \ar[r] & M \ar[r] & Aff(E) \ar[r] & \operatorname{Aut}(M) \ar[r] & 1
\end{tikzcd}

où les flèches verticales extrêmes $id_M$
et $\Theta_\Gamma$ sont connues, et où la flèche médiane
est à déterminer. On peut dire
que la catégorie des $\tilde{X}'$ au dessus
d'un [unclear: l'é.p.] de $\Gamma$ dessus satisfaisant
aux autres [unclear: axiomes] ) équivaut
à celle des couples $(E, \theta_{\tilde{\Gamma}})$, où $E$
est un $M$-torseur, et $\theta_{\tilde{\Gamma}}$ un
hom. de groupes $\tilde{\Gamma} \to Aff(E)$
compatible avec $id_M$ et $\Theta_\Gamma$ i.e. qui
rend commutatif (281). On
peut dire aussi une fois $E$ choisi,
(et ce choix est unique à isom. [unclear: unique] près -)
qu'on a défini une extension
$Aff(E)$ de $\operatorname{Aut}(M)$ par $M$, d'où
[unclear: tirée par] $\Gamma \to \operatorname{Aut}(M)$ une extension

\newpage

%% ===== Page 125 =====

%% [page 125 missing or failed: no entry]

\newpage

%% ===== Page 126 =====

%% [page 126 missing or failed: no entry]

\newpage

%% ===== Page 127 =====

711

(28bis)
\begin{tikzcd}
1 \ar[r] & \hat{M} \ar[r] \ar[d] & \tilde{\mathcal{M}} \ar[r] \ar[d] & \Gamma \ar[r] \ar[d, "N_{0,3}=\Gamma_{\mathbb{Q}} \times \mathfrak{S}_3" description, "\hat{\mathbb{Z}}^\times \times id_{\mathfrak{S}_3}"] & 1 \\
1 \ar[r] & \hat{M} \ar[r] & G_{\tilde{\rho}} \ar[r] & \hat{\mathbb{Z}}^\times \times \mathfrak{S}_3 \ar[r] & 1
\end{tikzcd}

D'ailleurs comme l'on attend - on a trouvé que
l'extension $G_{\tilde{\rho}}$ de $\hat{\mathbb{Z}}^\times \times \mathfrak{S}_3$ par $\hat{M}$ était
scindée, d'où il résulte que
l'ext. $\tilde{\mathcal{M}}$ de $\Gamma$ par $\hat{M}$ l'est aussi
(ce qui au demeurant est aussi conséquence de
l'existence de la tour de Fermat, en
y appliquant les réflexions générales
des pages précédentes). Bien mieux,
on a vu qu'il y a une unique classe
de $\hat{M}$-conj. de scindages de l'extension
$G_{\tilde{\rho}}$, ce qui induit une classe de
$\hat{M}$-conjugaison canonique (mais non
unique!) de scindages de $\tilde{\mathcal{M}} \to$ 
de $\Gamma$. Cela définit donc une
$\hat{M}$-conjugaison canonique de relèvements
sur $\tilde{X} = \widetilde{U_{0,3 \overline{\mathbb{Q}}}}$
de l'action de $\Gamma$ (à $\simeq$ près, ou ce
qui revient au m., un foncteur sur
$(U_{0,3 \mathbb{Q}}, \mathfrak{S}_3)$, de fibre un groupe profini

[MARGIN: NB $H^1(\Gamma, \hat{M})$ est un groupe fort gros! Mais voir cf th. p. 716]

\newpage

%% ===== Page 128 =====

%% [page 128 missing or failed: no entry]

\newpage

%% ===== Page 129 =====

%% [page 129 missing or failed: no entry]

\newpage

%% ===== Page 130 =====

%% [page 130 missing or failed: no entry]

\newpage

%% ===== Page 131 =====

suivante : après l'oubli des opérations de $\Gamma$ on trouve un revêtement [unclear: connexe & abélien] galoisien de $\mathfrak{X}$, correspondant à un groupe quotient $M$ de $\pi = \pi_1(\mathfrak{X}, x_0)$ - de sorte que $\mathfrak{X}'$ devient de façon naturelle un torseur sous $M_{\mathfrak{X}}$.

Le foncteur
$$ \mathfrak{X}' \longmapsto \mathfrak{X}'(x_0) \quad (\text{fibre en } x_0) $$
de la catégorie de ces $\mathfrak{X}'$ dans la catégorie des ensembles, s'en va en fait vers les ensembles $E$ munis des structures suivantes

a) $E$ est un $M$-torseur, et en tant que tel donne naissance à une extension
$$ 1 \to M \to \text{Aff}(E) \to \operatorname{Aut}(M) \to 1 $$
d'où une sous-extension induite

(288)
$$ 1 \to M \to \text{Aff}_\Gamma(E) \to \tilde{\Gamma} \to 1 $$

où $\tilde{\Gamma} = \operatorname{Im}(\Gamma \to \operatorname{Aut}(M))$

b) On a un hom d'extensions

(289)
\begin{tikzcd}
1 \ar[r] & M \ar[r] \ar[d, "id"'] & \tilde{M} \ar[r] \ar[d, "\theta_{\tilde{M}}"'] & \Gamma \ar[r] \ar[d, "can"] & 1 \\
1 \ar[r] & M \ar[r] & \text{Aff}_\Gamma(E) \ar[r] & \tilde{\Gamma} \ar[r] & 1
\end{tikzcd}

La catégorie des $\mathfrak{X}'$ (à un iso...) équivaut à celle des couples $(E, \theta_{\tilde{M}})$, où $E$ est un $M$-torseur, et $\theta_{\tilde{M}}$ un hom. d'extensions (289).

\newpage

%% ===== Page 132 =====

%% [page 132 missing or failed: no entry]

\newpage

%% ===== Page 133 =====

%% [page 133 missing or failed: no entry]

\newpage

%% ===== Page 134 =====

%% [page 134 missing or failed: no entry]

\newpage

%% ===== Page 135 =====

%% [page 135 missing or failed: no entry]

\newpage

%% ===== Page 136 =====

je n'avais pas eu à ma disposition
cette tour de Fermat, je n'aurais
certes pas suspecté que [crossed out: l'extension] $\widetilde{C}^{+\wedge} = \widetilde{C}^{+\wedge}_{H}$ de
$H^+ \simeq \widehat{\mathfrak{S}}_3$ par $M^\wedge$ soit scindée !

La conjecture qui se présente tout de
suite, c'est que ces deux revêtements
pro-étales sont isomorphes - l'isomor-
phisme entre les deux est alors
unique.

Pour l'autre voie de détermination
des classes d'isomorphie [unclear: (c.a.d. $\tilde{M}$ sur $M_{0,3} / [\Gamma_{\mathbb{Q}}, \Gamma_{\mathbb{Q}}]$)]
de scindages de $\Gamma = \Gamma_{\mathbb{Q}} \times \mathfrak{S}_3$ par
$M^\wedge$, on a par la suite exacte associée
à la suite spectrale de Hochschild

$$ (294) \quad 0 \to H^1(\Gamma_{\mathbb{Q}}, H^0(\widehat{\mathfrak{S}}_3, \widehat{M})) \to H^1(\Gamma, \widehat{M}) \to H^0(\Gamma_{\mathbb{Q}}, H^1(\widehat{\mathfrak{S}}_3, \widehat{M})) $$

Or on a vu que
$$ (295) \quad H^0(\widehat{\mathfrak{S}}_3, \widehat{M}) = 0 \quad (\text{[unclear: car] } H^0(L_f, \widehat{M}) = 0) $$

[crossed out: Et] d'autre part
$$ H^1(\widehat{\mathfrak{S}}_3, \widehat{M}) = H^1(\mathfrak{S}_3, M) \otimes_{\mathbb{Z}} \hat{\mathbb{Z}} \quad , $$

[MARGIN: donc $H^1(\Gamma, \widehat{M}) \simeq H^0(\Gamma_{\mathbb{Q}}, H^1(\widehat{\mathfrak{S}}_3, \widehat{M}))$]

\newpage

%% ===== Page 137 =====

%% [page 137 missing or failed: no entry]

\newpage

%% ===== Page 138 =====

%% [page 138 missing or failed: no entry]

\newpage

%% ===== Page 139 =====

%% [page 139 missing or failed: no entry]

\newpage

%% ===== Page 140 =====

Ceci, joint aux relations
$$H^\circ(\hat{\mathbb{Z}}^* \times \mathfrak{S}_3, \hat{M}) = H^1(\hat{\mathbb{Z}}^* \times \mathfrak{S}_3, \hat{M}) = 0 \quad,$$
permettait d'interpréter le groupe
$G_{\rho, \tilde{\sigma}}$ comme un groupe de transformations
sur l'âme $E$ de ce [unclear: pro-jeu singulier],
ce qui permettait d'interpréter
l'hom $\mathcal{M} \to G_{\rho, \tilde{\sigma}}$ (ou $\mathcal{M}^\sim \to G_{\rho, \tilde{\sigma}}$)
comme une spécification de $\mathcal{M}$
(ou $\tilde{\mathcal{M}} = \Pi_1(U_{0,3 \mathbb{Q}}, \underline{\mathfrak{S}}_3) = \Pi_1(M_{1,1 \mathbb{Q}}, \dots)$
comme, [unclear], sur $E$ [unclear: correspondant]
à un revêtement de $(U_{0,3 \mathbb{Q}}, \underline{\mathfrak{S}}_3)$
$\simeq M_{1,1 \mathbb{Q}}$ - et ce revêtement
n'est autre que le prorevêtement
de Fermat.

Il resterait à interpréter de façon
analogue l'hom $\mathcal{M} = \mathcal{N}_{1,1} \to G_{\rho, \sigma}$,
s'insérant dans l'hom. de suites exactes

(298)
\begin{tikzcd}
1 \ar[r] & \mu_M \ar[r] \ar[dd] & \mathcal{M} \ar[r] \ar[dd] & \mathcal{M}^\sim \ar[r] \ar[dd] & 1 \\
& (\pm 1, \omega_f) & (\mathcal{N}_{1,1}) & (\tilde{\mathcal{N}}_{1,1} = \mathcal{M}_{0,3}) & \\
1 \ar[r] & \mu_G \ar[r] & G_{\rho, \sigma} \ar[r] & G_{\rho, \tilde{\sigma}} \ar[r] & 1 \\
& \pm 1, \omega_f & & &
\end{tikzcd}

Pour [unclear] d'interpréter d'ailleurs
$\mathcal{M}$, en tant qu'extension centrale de $\mathcal{M}^\sim$

\newpage

%% ===== Page 141 =====

%% [page 141 missing or failed: no entry]

\newpage

%% ===== Page 142 =====

%% [page 142 missing or failed: no entry]

\newpage

%% ===== Page 143 =====

%% [page 143 missing or failed: no entry]

\newpage

%% ===== Page 144 =====

%% [page 144 missing or failed: no entry]

\newpage

%% ===== Page 145 =====

%% [page 145 missing or failed: no entry]

\newpage

%% ===== Page 146 =====

%% [page 146 missing or failed: no entry]

\newpage

%% ===== Page 147 =====

%% [page 147 missing or failed: no entry]

\newpage

%% ===== Page 148 =====

%% [page 148 missing or failed: no entry]

\newpage

%% ===== Page 149 =====

%% [page 149 missing or failed: no entry]

\newpage

%% ===== Page 150 =====

%% [page 150 missing or failed: no entry]

\newpage

%% ===== Page 151 =====

%% [page 151 missing or failed: no entry]

\newpage

%% ===== Page 152 =====

%% [page 152 missing or failed: no entry]

\newpage

%% ===== Page 153 =====

%% [page 153 missing or failed: no entry]

\newpage

%% ===== Page 154 =====

%% [page 154 missing or failed: no entry]

\newpage

%% ===== Page 155 =====

%% [page 155 missing or failed: no entry]

\newpage

%% ===== Page 156 =====

%% [page 156 missing or failed: no entry]

\newpage

%% ===== Page 157 =====

%% [page 157 missing or failed: no entry]

\newpage

%% ===== Page 158 =====

(328)
\begin{tikzcd}
{[\Gamma_\mathbb{Q} \hookrightarrow]} \Gamma_\mathbb{Q}^{et} \ar[r] & \mathcal{N}_\sigma^\S \simeq \hat{\mathbb{Z}}^* \cdot \hat{\mathbb{Z}} \\
{[\underset{\Pi_1(U_{0,3\mathbb{Q}})}{M_{0,3}^!} \hookrightarrow] M_{0,3}^{! et}} \ar[u, hook, "K_{1/2}"] \ar[r, "\text{''Fermat''}"', "(32\bar{5})"] & G_{p,\sigma}^\S \simeq \hat{\mathbb{Z}}^* \cdot \hat{M} \ar[u, hook]
\end{tikzcd}

Cela signifie donc que l'on trouve
l'hom. induit par (326)

(329) $$\Gamma_{\mathbb{Q}} \longrightarrow \mathcal{N}_\sigma^\S \simeq \hat{\mathbb{Z}}^* \cdot \hat{\mathbb{Z}}$$

en regardant l'image de l'hom.

(330) $$\Gamma_{\mathbb{Q}} \longrightarrow \hat{\mathbb{Z}}^* \cdot \hat{M} = Aff(\hat{M}) \ ,$$

dans $\mathcal{N}_\sigma^\S \subset \hat{\mathbb{Z}}^* \cdot \hat{M}$ , déduit de la
tour des courbes de Fermat en
regardant la fibre de la dite tour
au point $1/2$ , ce qui donne un
torseur sous $T_a(\hat{\mathbb{Z}}_{(m)}) \otimes_{\hat{\mathbb{Z}}} \hat{M}$ , s'explici-
tant galoisiennement justement par
l'hom. (326). Notons que le $K_{1/2}$
de $U_{0,3\mathbb{Q}}$ corresp. par l'iso. $\varphi$ (317)
de $U_{0,3\mathbb{Q}}$ avec $U_{1,\mathbb{Q}}$, au point [unclear: $(\xi_{1/2}, 1)$]

\newpage

%% ===== Page 159 =====

%% [page 159 missing or failed: no entry]

\newpage

%% ===== Page 160 =====

%% [page 160 missing or failed: no entry]

\newpage

%% ===== Page 161 =====

%% [page 161 missing or failed: no entry]

\newpage

%% ===== Page 162 =====

%% [page 162 missing or failed: no entry]

\newpage

%% ===== Page 163 =====

74

où l'image de $G \to \hat{\mathbb{Z}}^*$ est formée des $\mu \in \hat{\mathbb{Z}}^*$ tels que l'homothétie $\mu_T$ ne change pas la classe de l'ext. (342) de $A$ par $T$. L'hom (341) peut lors s'interpréter comme s'insérant dans un diagramme

(345)
\begin{tikzcd}
1 \ar[r] & \operatorname{Hom}_\mathbb{Z}(A, T_\infty(\bar{k}^*)) \ar[r] & G \ar[r] & \hat{\mathbb{Z}}^* \\
& & \Pi_1 \ar[u] \ar[ur, "\text{"caractère cyclotomique"}"']
\end{tikzcd}

[MARGIN: NB.] Si $S$ est sch. de type fini sur $\mathbb{Q}$, alors l'image de $\Pi_1 \to \hat{\mathbb{Z}}^*$ est un sous-groupe ouvert de $\hat{\mathbb{Z}}^*$, donc il en est a fortiori ainsi de l'image de $G \to \hat{\mathbb{Z}}^*$ - qui sera surjectif si le car. cyclotomique sur $\Pi_1$ est surjectif, i.e. si les seules racines de 1 sur $S$ sont $\pm 1$ (p. ex. si $S = Spec \mathbb{Q} \dots$).

Le diagramme (345) a
une propriété de fonctorialité évidente
pour $A$ variable (disons), i.e. pour un

\newpage

%% ===== Page 164 =====

hom. de groupes abéliens
(347) $A' \xrightarrow{f} A$ ,

l'extension $E_{\infty}(A')_S$ s'identifie
à l'image inverse de l'ext. $E_{\infty}(A)_S$
par (347), d'où un hom.
$G_A \to G_{A'}$ (injectif si $f$ est surjectif)
s'insérant dans un diagramme
commutatif

(349)
\begin{tikzcd}
1 \ar[r] & \operatorname{Hom}_{\mathbb{Z}}(A', T) \ar[r] & G_{A'} \ar[r] & \hat{\mathbb{Z}}^* \\
1 \ar[r] & \operatorname{Hom}_{\mathbb{Z}}(A, T) \ar[r] \ar[u] & G_A \ar[r] \ar[u] & \hat{\mathbb{Z}}^* \ar[u, "1"'] \\
& & \Pi_1 \ar[u, "\gamma_A"'] \ar[uu, bend left=45, "\gamma_{A'}"] \ar[ur, "\text{car. cyclot.}"']
\end{tikzcd}

où la flèche sur les $\operatorname{Hom}_{\mathbb{Z}}$ est trans-
posée de (347). Cela montre que
pour étudier la structure de l'image de $\gamma_A$ (345),
il est commode, par fonctorialité, de se

\newpage

%% ===== Page 165 =====

730

étudier le cas où $j$ est injectif
(350) $$j : A \hookrightarrow \mathbb{G}_m(S) \ .$$

On peut d'ailleurs "diviser" encore un peu le diagramme (345), en introduisant

(351) $$\Pi_1^+ = \Pi_1^+(S, \bar{s}) \overset{df}{=} \operatorname{Ker}\Big(\Pi_1(S, \bar{s}) \overset{cyclot.}{\longrightarrow} \hat{\mathbb{Z}}^*\Big)$$

et l'hom. de suites exactes

(352)
\begin{tikzcd}
1 \ar[r] & \Pi_1^+ \ar[r] \ar[d, "\gamma_j^+"'] & \Pi_1 \ar[r, "\chi_{\Pi}"] \ar[d, "\gamma_j"'] & \hat{\mathbb{Z}}^* \ar[d] \\
1 \ar[r] & \operatorname{Hom}_{\hat{\mathbb{Z}}}(A, T) \ar[r] & G_j \ar[r] & \hat{\mathbb{Z}}^*
\end{tikzcd}
.

On peut considérer que l'étude de la "taille" de $\operatorname{Im}(\gamma_j : \Pi_1 \to G_j)$ est essentiellement ramenée à celle de la taille de $Im \chi_{\Pi}$ et de celle de $Im \gamma_j^+$; de façon précise, si on pose

(353) $$Im \gamma_j = \tilde{\Pi}_1 \ , \ Im \gamma_j^+ = \tilde{\Pi}_1^+ \ , \ Im \chi_{\Pi} = \Gamma$$

on a une suite exacte

(354) $$1 \to \tilde{\Pi}_1^+ \to \tilde{\Pi}_1 \to \Gamma \to 1 \ .$$

Je voudrais essayer de donner des

\newpage

%% ===== Page 166 =====

%% [page 166 missing or failed: no entry]

\newpage

%% ===== Page 167 =====

%% [page 167 missing or failed: no entry]

\newpage

%% ===== Page 168 =====

%% [page 168 missing or failed: no entry]

\newpage

%% ===== Page 169 =====

%% [page 169 missing or failed: no entry]

\newpage

%% ===== Page 170 =====

%% [page 170 missing or failed: no entry]

\newpage

%% ===== Page 171 =====

%% [page 171 missing or failed: no entry]

\newpage

%% ===== Page 172 =====

%% [page 172 missing or failed: no entry]

\newpage

%% ===== Page 173 =====

%% [page 173 missing or failed: no entry]

\newpage

%% ===== Page 174 =====

Il s'impose de partir d'approximations de
Lie au moins aussi fines que celle du
présent $\S$, qui apparait comme l'approximation
"minima" dans laquelle on puisse exprimer
"algébriquement" la tour des courbes de
Fermat, ~~via~~ comme on l'a vu dans 13°),
à partir de $M_{0,3} \subset \hat{M}_{0,3}$.

Tenant compte des opérations de $\mathfrak{S}_3$
sur $U_n$, la suite ex.te (371) se prolonge
d'ailleurs dans une suite ex.te plus
riche

$$(372) \quad 1 \longrightarrow \pi_1(U_{n,\overline{\mathbb{Q}}}) \longrightarrow \underset{\overset{\wedge}{\pi_1(U_1,\mathfrak{S}_3) \simeq M_{0,3}}}{\pi_1(U_n,\mathfrak{S}_3)} \longrightarrow \Gamma_{\mathbb{Q}} \times \mathfrak{S}_3 \longrightarrow 1$$

et les calculs faits dans le présent $\S$ semblent
suggérer que cette structure d'extension plus
riche jouera un rôle important.

Du point de vue de l'approximation
de Lie de $M_{0,3}$ développée dans le présent
$\S$, il y a lieu de remplacer la structure
d'extension (372) par une structure quotient
moins fine, en divisant par un sous-groupe
invariant de $\pi_1(U_n, \mathfrak{S}_3)$ contenu dans

\newpage

%% ===== Page 175 =====

$73r$

$\Pi_1(\mathcal{U}_n, \overline{\mathbb{Q}})$, savoir l'intersection dans $M_{0,3}$ de ce
groupe avec $[\overline{\pi_0, \hat{j}}, \overline{\pi_0, \hat{j}}]$. On trouve un
diagramme commutatif

(373)
\begin{tikzcd}
1 \ar[r] & M^\wedge[n] \ar[r] \ar[d, equal] & \widetilde{\Pi_1(\mathcal{U}_n, \mathfrak{S}_3)} \ar[r] \ar[d, "p_n"'] & \Gamma_{\mathbb{Q}} \times \mathfrak{S}_3 \ar[r] \ar[d, "\chi \times id"] & 1 \\
1 \ar[r] & M^\wedge[n] \ar[r] \ar[d] & \widetilde{G}_{\rho, \sigma}[n] \ar[r] \ar[d] & \hat{\mathbb{Z}}^* \times \mathfrak{S}_3 \ar[r] \ar[d, equal] & 1 \\
1 \ar[r] & M^\wedge \ar[r] & \widetilde{G}_{\rho, \sigma} \ar[r] & \hat{\mathbb{Z}}^* \times \mathfrak{S}_3 \ar[r] & 1
\end{tikzcd}

où l'image de $M^\wedge[n] \to M^\wedge$ est $n M^\wedge$.

Voici une façon de construire "algébriquement"
ces diagrammes (373) (pour $n$ variable)
à partir du diagramme

(374)
\begin{tikzcd}
1 \ar[r] & M^\wedge \ar[r] \ar[d, equal] & \overset{\mathcal{E}_1}{\widetilde{\Pi_1(\mathcal{U}_1, \mathfrak{S}_3)}} \ar[r] \ar[d] & \Gamma_{\mathbb{Q}} \times \mathfrak{S}_3 \ar[r] \ar[d, "\chi \times id"] & 1 \\
& & (M_{0,3}/[\overline{\pi_0, \hat{j}}, \overline{\pi_0, \hat{j}}]) \ar[d] & & \\
1 \ar[r] & M^\wedge \ar[r] & \widetilde{G}_{\rho, \sigma} \ar[r] \ar[d, "\simeq"'] & \hat{\mathbb{Z}}^* \times \mathfrak{S}_3 \ar[r] & 1 \\
& & (\hat{\mathbb{Z}}^* \times \mathfrak{S}_3).M^\wedge & & 
\end{tikzcd}

On choisit un scindage de la première,
qui donne [unclear] un scindage de la seconde
(scindage [unclear: par] un [unclear: sous-groupe] fini [unclear]
[unclear: engendré par] un unique élément de $M^\wedge \dots$). [unclear: En prenant]
[unclear] on peut déduire $U_n$ en-dessous de [unclear]

[MARGIN: On a vu qu'en un sens un peu tautologique les scindages de l'ext. $\mathcal{E}_1$ [unclear: équivalent] à ceux de l'ext. [unclear] (dans 12°), i.e. à un [unclear: ensemble] de choix tels [unclear: canoniques]]

\newpage

%% ===== Page 176 =====

%% [page 176 missing or failed: no entry]

\newpage

%% ===== Page 177 =====

%% [page 177 missing or failed: no entry]

\newpage

%% ===== Page 178 =====

%% [page 178 missing or failed: no entry]

\newpage

%% ===== Page 179 =====

%% [page 179 missing or failed: no entry]

\newpage

%% ===== Page 180 =====

%% [page 180 missing or failed: no entry]

\newpage

%% ===== Page 181 =====

%% [page 181 missing or failed: no entry]

\newpage

%% ===== Page 182 =====

"invariantes sous $\pi$" dans un sens évident.
En particulier, si $\pi$ est une extension
d'un groupe fini par un groupe
discret $\pi_0$, alors les $\mathcal{L}$-structures sur
$\pi$ correspondent aux $\mathcal{L}$-structures
sur $\pi_0$ invariantes par $\pi / \pi_0 \to$ [unclear: $Out(\pi_0)$]
opérant sur l'ensemble des $\mathcal{L}$-structures de $\pi_0$
de façon évidente. [NB $\pi \mapsto \mathcal{L}(\pi)$ est
un foncteur en les groupes extérieurs
$\pi$, avec les isomorphismes extérieurs
comme flèches ...]

Exemple standard : soient $X$ une
surface 0-conn. "de présentation finie",
$G$ un groupe fini opérant sur $X$,

(5) $\pi = \pi_1(X, G ; -)$

où $-$ est un pt base de $X$, de
sorte qu'on a une suite exacte

(6) 
\begin{tikzcd}[row sep=small]
1 \ar[r] & \pi_0 \ar[r] & \pi \ar[r] & G \ar[r] & 1 \\
& \pi_1(X,-) \ar[u, equal]
\end{tikzcd}

On a donc sur $\pi_0$ une $\mathcal{L}$-structure : [unclear: locale]

\newpage

%% ===== Page 183 =====

%% [page 183 missing or failed: no entry]

\newpage

%% ===== Page 184 =====

%% [page 184 missing or failed: no entry]

\newpage

%% ===== Page 185 =====

%% [page 185 missing or failed: no entry]

\newpage

%% ===== Page 186 =====

%% [page 186 missing or failed: no entry]

\newpage

%% ===== Page 187 =====

%% [page 187 missing or failed: no entry]

\newpage

%% ===== Page 188 =====

%% [page 188 missing or failed: no entry]

\newpage

%% ===== Page 189 =====

%% [page 189 missing or failed: no entry]

\newpage

%% ===== Page 190 =====

%% [page 190 missing or failed: no entry]

\newpage

%% ===== Page 191 =====

%% [page 191 missing or failed: no entry]

\newpage

%% ===== Page 192 =====

%% [page 192 missing or failed: no entry]

\newpage

%% ===== Page 193 =====

%% [page 193 missing or failed: no entry]

\newpage

%% ===== Page 194 =====

Plus intrinsèquement, il doit être
vrai qu'il existe un rev. étale fini
de $M$ (qu'on peut supposer galoisien,
de groupe $G$) qui soit un
[crossed out: schéma algébrique] - lequel sera néc.
quasi-projectif, lisse, de dim relative 1
sur $k$. Supposons en tous cas
qu'il en soit ainsi, choisissons
un pt base géom. $a$ de $M$, et
considérons

(10) $$\pi = \pi_1(M, a) \simeq \pi_1(U, G; a) \ .$$

Nous supposerons désormais $k$ de
car. zéro. Supposons d'abord $k$
alg clos. Utilisons la suite exacte

(11) $$1 \longrightarrow \underset{\pi_0}{\pi_1(U_{\tilde{a}})} \longrightarrow \underset{\pi}{\pi_1(M, G; a)} \longrightarrow G \longrightarrow 1$$

et la structure : [unclear: toute] naturelle
sur $\pi_0$, on trouve une $L$-structure
sur $\pi = \pi_1(M, a)$, qui [unclear: manifeste] sur

\newpage

%% ===== Page 195 =====

%% [page 195 missing or failed: no entry]

\newpage

%% ===== Page 196 =====

%% [page 196 missing or failed: no entry]

\newpage

%% ===== Page 197 =====

%% [page 197 missing or failed: no entry]

\newpage

%% ===== Page 198 =====

%% [page 198 missing or failed: no entry]

\newpage

%% ===== Page 199 =====

%% [page 199 missing or failed: no entry]

\newpage

%% ===== Page 200 =====

%% [page 200 missing or failed: no entry]

\newpage

%% ===== Page 201 =====

%% [page 201 missing or failed: no entry]

\newpage

%% ===== Page 202 =====

%% [page 202 missing or failed: no entry]

\newpage

%% ===== Page 203 =====

741

ReNB Sauf erreur, la propriété 1) signifie qu'
on peut trouver un schéma en groupes $G$
affine de t.f. sur $\mathbb{Z}$, qui définit $G_\mathbb{Q}$, et
tel que le (27) se factorise par $G(\hat{\mathbb{Z}})$.
(29) $$M \longrightarrow G(\hat{\mathbb{Z}})$$

2°) Compatibilité avec $\chi$ :

(30)
\begin{tikzcd}
M \ar[r] \ar[d, "\chi_M"'] & G(\mathbb{A}_\mathbb{Q}) \ar[d, "\chi_G"] \\
\hat{\mathbb{Z}} \ar[r, hook] & \mathbb{G}_m(\mathbb{A}_\mathbb{Q}) \simeq \mathbb{A}_\mathbb{Q}^* \subset \prod_{l \in P} \mathbb{Q}_l^*
\end{tikzcd}

est commutatif.

3°) Compatibilité avec la [unclear: désintégration] $\pi_0$

(31)
\begin{tikzcd}
\Pi'_0 \ar[r, hook] \ar[d, "\text{factorisation}"'] & \Pi_0 \ar[r, hook] & M \ar[d] \\
G(\mathbb{Q}) \ar[rr, hook] & & G(\mathbb{A}_\mathbb{Q})
\end{tikzcd}
,

pour un sous-groupe d'indice fini
convenable $\Pi'_0$ de $\Pi_0$.

NB Quitte à changer les notations, [unclear: on peut supposer] $\Pi'_0 = \Pi_0$.

4°) Caractère minimal (irrédundance) :
Si $G'$ est un sous-schéma en groupes de

\newpage

%% ===== Page 204 =====

%% [page 204 missing or failed: no entry]

\newpage

%% ===== Page 205 =====

250

Réflexion faite, je préfère poser les choses de
façon légèrement différente, en introduisant,
dans le LT-gr. profini $M$, un ~~[unclear]~~
distingué $\pi$ avec une L-structure (par
ex. une structure : toute pièce,
si par ex $\pi = \operatorname{Sl}(2, \mathbb{Z})$) définissant la
L T-structure de $M$, une discrétisation
$\pi_0$ de $\pi$ (dans un sens évident, pour les
L-structures) un sous-groupe discret
$\mathfrak{S}_0 \supset \pi_0$ normalisant $\pi_0$ (NB. dans mon petit
exemple [unclear] on a $\mathfrak{S}_0 = \pi_0$), tel que
[MARGIN: son adh.] $\mathfrak{S} = \widehat{\mathfrak{S}}_0$ dans $M$ soit distingué dans $M$.

On a donc un diag. de groupes

(32)
\begin{tikzcd}
& \tau \ar[d, no head] & \Gamma \ar[d, no head] \\
1 \ar[r, phantom, "\subset"] & \pi \ar[r, phantom, "\subset"] \ar[d, equal] & \mathfrak{S} \ar[r, phantom, "\subset"] \ar[d, equal] & M \\
& \widehat{\pi}_0 & \widehat{\mathfrak{S}}_0 \\
& \uparrow & \uparrow \\
1 \ar[r, phantom, "\subset"] & \pi_0 \ar[r, phantom, "\subset"] \ar[uu] & \mathfrak{S}_0 \ar[uu]
\end{tikzcd}

où la première ligne est formée de
groupes profinis, la deuxième de groupes
discrets. On notera que $\mathfrak{S}$ s'identifie : au
quotient du complété profini de $\mathfrak{S}_0$, mais

\newpage

%% ===== Page 206 =====

%% [page 206 missing or failed: no entry]

\newpage

%% ===== Page 207 =====

%% [page 207 missing or failed: no entry]

\newpage

%% ===== Page 208 =====

tion $\mathbb{G}_0 \to G(\mathbb{Q})$, rendant commutatif le diagramme

(35)
\begin{tikzcd}
\mathbb{G}_0 \ar[r, hook] \ar[d] & \mathcal{M} \ar[d] \\
G(\mathbb{Q}) \ar[r, hook] & G(\mathbb{A}_{\mathbb{Q}})
\end{tikzcd}

Les structures de la forme (32) 
[MARGIN: Quitte par la suite à introduire des variantes sur ces structures]
méritent sans doute un nom, je vais les appeler - ne serait-ce que provisoirement - structures de Galois-Teichmüller.

Si $G$ est un groupe algébrique affine sur $\mathbb{Q}$, un hom $\mathcal{M} \to G(\mathbb{A}_{\mathbb{Q}})$ satisfaisant les conditions 1) 2) $3^{\circ \text{forte}}$) 4) est appelé un "hom de Lie" de la structure de Galois-Teichmüller.

Je vais désigner par la suite par (en abrégé $M$) $M_{\mathbb{Q}}$ le schéma en groupes sur $\mathbb{Q}$ affine.
Je désigne par $\mathbb{G}_{\mathbb{Q}}$ le sous-schéma en groupes de $M_{\mathbb{Q}}$ engendré par l'

\newpage

%% ===== Page 209 =====

%% [page 209 missing or failed: no entry]

\newpage

%% ===== Page 210 =====

%% [page 210 missing or failed: no entry]

\newpage

%% ===== Page 211 =====

%% [page 211 missing or failed: no entry]

\newpage

%% ===== Page 212 =====

$\underline{\Pi}$, $\underline{\mathfrak{S}}$ les adhérences schématiques
dans $\underline{M}$ de $\underline{\Pi}_\mathbb{Q}$, $\underline{\mathfrak{S}}_\mathbb{Q}$, d'où un diagramme
de schémas en groupes sur $\mathbb{Z}$

(46) $$ \underbrace{\underline{\Pi} \subset \underline{\mathfrak{S}}}_{\underline{\Gamma}} \subset \underline{M} $$
[MARGIN: with another brace under $\underline{\mathfrak{S}} \subset \underline{M}$ labeled $\underline{N}$]

(sous réserve de représentabilité des
quotients $\underline{\tau}, \underline{\Gamma}, \underline{N}$ - mais je crois que
pour Raynaud il n'y a jamais de pb...).

[unclear: On a des factorisations] [crossed out]

(47)
\begin{tikzcd}
\underline{\Pi} \ar[r, hook] \ar[d] & \underline{\mathfrak{S}} \ar[r, hook] \ar[d] & \underline{M} \ar[d] \\
\underline{\Pi}(\hat{\mathbb{Z}}) \ar[r, hook] & \underline{\mathfrak{S}}(\hat{\mathbb{Z}}) \ar[r, hook] & \underline{M}(\hat{\mathbb{Z}})
\end{tikzcd}

On suppose bien sûr également que
l'hom $\chi_{\underline{M}} : \underline{M}_\mathbb{Q} \to \mathbb{G}_{m \mathbb{Q}}$ se
prolonge en

(48) $$ \chi_{\underline{M}} : \underline{M} \to \mathbb{G}_{m \mathbb{Z}} $$

lequel nécessairement se factorise par

(49) $$ \chi_{\underline{\Gamma}} : \underline{\Gamma} \to \mathbb{G}_{m \mathbb{Z}} . $$

\newpage

%% ===== Page 213 =====

%% [page 213 missing or failed: no entry]

\newpage

%% ===== Page 214 =====

%% [page 214 missing or failed: no entry]

\newpage

%% ===== Page 215 =====

%% [page 215 missing or failed: no entry]

\newpage

%% ===== Page 216 =====

%% [page 216 missing or failed: no entry]

\newpage

%% ===== Page 217 =====

756

Le schéma en groupes $\underline{\Pi}^\circ$ sur $S_0 = Spec \mathbb{Z}$, vu en tant
que faisceau étale sur le "gros topos étale" de
$S_0$, est engendré par ses sous-faisceaux en
sous-groupes
$\underline{k}'_!(\underline{G}_\circ)$ (avec répétition possible d'une même
classe de conjugaison), et encore $\underline{\Pi}$ est
(Et il y a [unclear] pratiquement quant aux
les $\underline{k}'_!(\underline{G}_\circ)$ lors d'un changement de base) et
le groupe discret $(\Pi_0)_{S_0}$. (Et il se pose
la question naturellement d'expliciter les
relations... suffit-il d'écrire les relations de
$\Pi_0$, plus les commutativités de

(57)
\begin{tikzcd}
T_0 \ar[r, "k'_!"] \ar[d, "can"'] & \Pi_0 \ar[d, "can"] \\
\underline{T}_0 \ar[r, "\underline{k}'_!"'] & \underline{\Pi}
\end{tikzcd}
? )

Autres conditions concernant les schémas
en groupes

(58)
$$ \gamma_{\underline{\Pi}} = \underline{\Pi} / \underline{\Pi}^\circ, \gamma_{\underline{\mathfrak{S}}} = \underline{\mathfrak{S}} / \underline{\mathfrak{S}}^\circ, \gamma_{\underline{M}} = \underline{M} / \underline{M}^\circ $$

- on voudrait que ces schémas soient
en groupes constants, correspondant aux
schémas groupes finis discrets qui décrivent les
quotients correspondants des groupes sur $\mathbb{Q}$,
lesquels groupes discrets sont des quotients

\newpage

%% ===== Page 218 =====

%% [page 218 missing or failed: no entry]

\newpage

%% ===== Page 219 =====

%% [page 219 missing or failed: no entry]

\newpage

%% ===== Page 220 =====

%% [page 220 missing or failed: no entry]

\newpage

%% ===== Page 221 =====

%% [page 221 missing or failed: no entry]

\newpage

%% ===== Page 222 =====

donnant lieu à des compatibilités

(68)
\begin{tikzcd}
T_0 \ar[r, "K_i"] \ar[d, "can"'] & \underline{\mathfrak{S}}_0 \ar[d, "(65)"] \\
T_{0\,Q}(Q) \ar[r, "\underline{K}_i(Q)"'] & \underline{\mathfrak{S}}(Q) \ .
\end{tikzcd}

Je considère alors le schéma en groupes [MARGIN: extérieurs (resp. à un...)] sur $Q$ des (automorphismes de $\underline{\mathfrak{S}}_Q$, formé des $\mu \in \underline{Aut}(\underline{\mathfrak{S}}_Q)$, tels que pour tt $i \in I$, il existe loc. une section $f_i$ de $\underline{\Pi}_Q$ (engendré par l'image de $\Pi_0$ dans $\underline{\mathfrak{S}}(Q)$), telle que l'on ait

(69) $$ \mu(\underline{K}_i(\lambda)) = int(f_i)^{-1}(\mu \lambda) \ . $$

Je désigne par $\underline{M}_Q^\varphi$ ce schéma en groupes sur $Q$, s'envoyant vers $\underline{G}_Q$ par $(u,\mu) \mapsto \mu$, par un hom. qui est trivial sur $\operatorname{Im}(\underline{\mathfrak{S}}_Q \to \underline{M}_Q^\varphi)$, d'où se factorise par un hom.

(70) $$ \underline{\Gamma}_Q^\varphi \overset{df}{=} \underline{M}_Q^\varphi / Im \underline{\mathfrak{S}}_Q \xrightarrow{\lambda^\varphi} \underline{G}_Q \ . $$

Si maintenant $(\underline{\mathfrak{S}}_0, \Pi_0)$ fait partie d'un système de [unclear]-Teichmüller

\newpage

%% ===== Page 223 =====

$(\Pi_0 \subset \mathbb{G}_0 \subset \mathcal{M})$, et si $\mathbb{G}_0 \xrightarrow{\varphi} \underline{\mathbb{G}}^\varphi_\mathbb{Q}$ provient d'un
hom de Lie de ce système, $\mathcal{M} \to M(\mathbb{A}_\mathbb{Q})$
par la méthode vue la section précédente,
alors on trouve un hom. naturel
via op. de $\underline{M}_\mathbb{Q}$ sur $\underline{\mathbb{G}}_\mathbb{Q}$, et $\lambda_M$
(71) $$ \underline{M}_\mathbb{Q} \xrightarrow{\varphi} \underline{M}^\varphi_\mathbb{Q} $$
rendant commutatif le diagramme

(72)
\begin{tikzcd}
\mathbb{G}_0 \ar[r] \ar[d] & \mathcal{M} \ar[d] & \underline{\mathbb{G}}_\mathbb{Q} \ar[r] \ar[dr, "via\ \text{can.}\ int."'] & \underline{M}_\mathbb{Q} \ar[d, "\varphi"] \\
\underline{\mathbb{G}}(\mathbb{Q}) \ar[r] \ar[d, "int."'] & M(\mathbb{A}_\mathbb{Q}) \ar[d] & & \underline{M}^\varphi_\mathbb{Q} \\
\underline{M}^\varphi_\mathbb{Q}(\mathbb{Q}) \ar[r, hook] & \underline{M}^\varphi_\mathbb{Q}(\mathbb{A}_\mathbb{Q}) & &
\end{tikzcd}

Si le centre de $\underline{\mathbb{G}}_\mathbb{Q}$ est trivial (ce qui est
une hyp. "crédible" surtout si le centre de
$\mathbb{G}_0$ l'est ...) lors $\underline{\mathbb{G}}_\mathbb{Q} \to \underline{M}^\varphi_\mathbb{Q}$ est
injectif, et $\underline{\mathbb{G}}_\mathbb{Q}$ s'interprète comme
"le $\underline{\mathbb{G}}_\mathbb{Q}$" "procréé" par un hom. de
Lie
$$ (\Pi \subset \mathbb{G} \subset \mathcal{M}) \longrightarrow (\Pi \subset \underline{\mathbb{G}}_\mathbb{Q} \subset \underline{M}^\varphi_\mathbb{Q}) , $$
qui a un certain caractère universel
p.r. à l'hom. donné (65) $\varphi : \mathbb{G}_0 \to \underline{\mathbb{G}}(\mathbb{Q})$.

\newpage

%% ===== Page 224 =====

Quand on ne suppose pas d'avance l'exis-
tence d'un tel hom. de Lie dans lequel
s'insère $\underline{\mathfrak{S}}_\mathbb{Q}$ , il n'est pas automatique
qu'on puisse trouver un hom.

(73) $$\varphi_M : \mathcal{M} \longrightarrow \underline{\mathcal{M}}_\mathbb{Q}^\varphi (A_\mathbb{Q})$$

qui rende compte des opérations de
$\mathcal{M}$ sur $\underline{\mathfrak{S}}_\mathbb{Q}$ , par la commutativité
d'un diagramme , pour tout $g \in \mathcal{M}$

(74)
\begin{tikzcd}
\underline{\mathfrak{S}} \ar[r, "int(g)|\underline{\mathfrak{S}}"] \ar[d, "\hat{\varphi}"'] & \underline{\mathfrak{S}} \ar[d, "\hat{\varphi}"] \\
\underline{\mathfrak{S}}(A_\mathbb{Q}) \ar[r, "int(\varphi_M(g))|\underline{\mathfrak{S}}(A_\mathbb{Q})"'] & \underline{\mathfrak{S}}(A_\mathbb{Q})
\end{tikzcd}

On s'intéresse seulement aux $\varphi$ (65)
pour lesquels (73) existe - au moins
quitte à se restreindre à un sous-groupe
ouvert de $\mathcal{M}$...

\newpage

%% ===== Page 225 =====

%% [page 225 missing or failed: no entry]

\newpage

%% ===== Page 226 =====

%% [page 226 missing or failed: no entry]

\newpage

%% ===== Page 227 =====

761

et considérons l'hom. canonique
$$ G \longrightarrow G/G^\circ \xrightarrow{\sim} \mathcal{T}(X,S) \quad \text{(groupe de Teichmüller)} $$
on a alors un hom.
canonique d'extensions

[MARGIN: Pour la démonstration de la dite proposition voir [unclear] (p. 755 ff, p. 764 ff)]

(7)
\begin{tikzcd}
1 \ar[r] & T_0^I \ar[r] \ar[d] & \tilde{G} \ar[r] \ar[d] & G \ar[r] \ar[d] & 1 \\
1 \ar[r] & T_0^I \ar[r] & Sp\mathcal{T}(X,S) \ar[r] & \mathcal{T}(X,S) \ar[r] & 1 \ ,
\end{tikzcd}

$Sp\mathcal{T}(X,S)$ est le groupe de Teichmüller
spécial de $(X,S)$.

NB On notera que si pour tout $i \in I$ on
fixe une origine $e_i$ sur $S_i$, il
résulte de la propriété universelle de
$\tilde{S}$ que la donnée d'un relèvement
$\tilde{u}$ de $u \in G$ à $\tilde{S}$ revient à la donnée,
pour tout $i \in I$, d'un relèvement de
$u(e_i)$ en un point $\tilde{a}_i$ de $\tilde{S}_{\sigma(i)}$ (i.e., de $\tilde{S}_{u(i)}$).

\newpage

%% ===== Page 228 =====

%% [page 228 missing or failed: no entry]

\newpage

%% ===== Page 229 =====

%% [page 229 missing or failed: no entry]

\newpage

%% ===== Page 230 =====

%% [page 230 missing or failed: no entry]

\newpage

%% ===== Page 231 =====

%% [page 231 missing or failed: no entry]

\newpage

%% ===== Page 232 =====

%% [page 232 missing or failed: no entry]

\newpage

%% ===== Page 233 =====

%% [page 233 missing or failed: no entry]

\newpage

%% ===== Page 234 =====

%% [page 234 missing or failed: no entry]

\newpage

%% ===== Page 235 =====

%% [page 235 missing or failed: no entry]

\newpage

%% ===== Page 236 =====

$$
(29) \quad \begin{cases}
\Gamma_\infty = \tau_\infty = \begin{pmatrix} -1 & 0 \\ 0 & 1 \end{pmatrix} \\
\Gamma_0 = \tau'_\infty = \begin{pmatrix} 0 & 1 \\ 1 & 0 \end{pmatrix} \\
\Gamma_1 = -\tau'_0 = \begin{pmatrix} -1 & 1 \\ 0 & 1 \end{pmatrix}
\end{cases}
\qquad
(30) \quad \begin{cases}
\Gamma_0(\rho) = \rho^{-1}, \Gamma_0(\sigma) = \sigma^{-1} \\
\Gamma_1 = \rho^{-1} \Gamma_0 = \Gamma_0 \rho \\
\Gamma_\infty = \sigma \Gamma_0 = \Gamma_0 \sigma^{-1}
\end{cases}
$$

Notons que $(\tilde{\rho}, \tilde{\tau})$ engendre un groupe
isomorphe à $D_\infty$, que je note $D_{\tilde{\rho}}$,
et que $(\tilde{\sigma}, \tilde{\tau})$ engendre un groupe $D_{\tilde{\sigma}}$,
isomorphe à $D_\infty$. Enfin $(\omega_0, \tilde{\tau})$ enge-
ndre un groupe $D_{\omega_0}$, isomorphe encore
à $D_\infty$. On a un diagramme

$$
(31) \quad 
\begin{tikzcd}
& D_{\omega_0} \ar[dl, "i"'] \ar[dr, "j"] & \\
D_{\tilde{\rho}} & & D_{\tilde{\sigma}}
\end{tikzcd}
\qquad
\begin{array}{l}
\begin{cases} i(\omega_0) = \tilde{\rho}^3 \\ i(\tilde{\tau}) = \tilde{\tau} \end{cases} \\
\begin{cases} j(\omega_0) = \tilde{\sigma}^2 \\ j(\tilde{\tau}) = \tilde{\tau} \end{cases}
\end{array}
$$

et le cor. 1 s'écrit sous la forme

$$
(32) \quad \operatorname{Gl}(2, \mathbb{Z})^\sim \simeq D_{\tilde{\rho}} *_{D_{\omega_0}} D_{\tilde{\sigma}}
$$

(somme amalgamée de groupes).

Or on peut prendre dans $D_{\tilde{\rho}}$, $D_{\tilde{\sigma}}$ les
générateurs involutifs

$$
(33) \quad \begin{cases}
\tilde{\Gamma}_0 = \tilde{\tau} = \tilde{\tau}'_\infty \\
\tilde{\Gamma}_\rho = .......
\end{cases}
\text{et} \quad
\begin{matrix}
\tilde{\Gamma}_1 \overset{df}{=} \tilde{\rho}^{-1} \tilde{\Gamma}_0 = \tilde{\Gamma}_0 \tilde{\rho} \text{ dans } D_{\tilde{\rho}} \\
\text{et } \tilde{\Gamma}_\infty \overset{df}{=} \tilde{\sigma} \tilde{\Gamma}_0 = \tilde{\Gamma}_0 \tilde{\sigma}^{-1}
\end{matrix}
$$

\newpage

%% ===== Page 237 =====

%% [page 237 missing or failed: no entry]

\newpage

%% ===== Page 238 =====

%% [page 238 missing or failed: no entry]

\newpage

%% ===== Page 239 =====

%% [page 239 missing or failed: no entry]

\newpage

%% ===== Page 240 =====

%% [page 240 missing or failed: no entry]

\newpage

%% ===== Page 241 =====

[MARGIN: 268 crossed out]

$a \in \mathbb{R}^2 \setminus \mathbb{Z}^2$,
au dessus de $a_{1,1}$ , on a donc un iso. can.

(42) $$ \Pi_{1,1} \simeq \Pi_1(\mathbb{R}^2 \setminus \mathbb{Z}^2, \mathbb{Z}^2; a) $$

d'où une suite exacte

(43)
\begin{tikzcd}
1 \ar[r] & \Pi_1(\mathbb{R}^2 \setminus \mathbb{Z}^2) \ar[r] & \Pi_{1,1} \ar[r] \ar[d, equal] & \mathbb{Z}^2 \ar[r] \ar[d, "\simeq"] & 1 \\
& & \Pi_1(U_{1,1}=X_1 \setminus \{a_{1,1}\}) & \Pi_1(X_1) \ar[d, "\simeq"] \\
& & & (\Pi_{1,1})_{ab}
\end{tikzcd}

Les éléments du $\Pi_1$ mixte s'explicitent par le calcul de Jean Malgoire à l'aide de couples

(44) $$ (\gamma, l) \quad (\gamma \in \mathbb{Z}^2, l \text{ classe de chemin de } \mathbb{R}^2 \setminus \mathbb{Z}^2 \text{ de } a \text{ à } a+\gamma) $$

qu'on "compose comme les fonctions".

Nous prendrons, au dessus des générateurs $e_0, e_1$ de $\mathbb{Z}^2$, les éléments $x, y$ définis par les chemins rectilignes de $a$ à $a+e_0$, resp. $a+e_1$, à partir de l'origine

(45) $$ a = \tilde{a}_{1,1} = (+\frac{1}{2}, +\frac{1}{2}) $$

Alors l'élément $xyx^{-1}y^{-1}$ est figuré par le chemin ci-dessus, qui est visiblement opposé au lacet "rectiligne" direct de $a$ autour de l'origine $0$.

\newpage

%% ===== Page 242 =====

%% [page 242 missing or failed: no entry]

\newpage

%% ===== Page 243 =====

%% [page 243 missing or failed: no entry]

\newpage

%% ===== Page 244 =====

%% [page 244 missing or failed: no entry]

\newpage

%% ===== Page 245 =====

%% [page 245 missing or failed: no entry]

\newpage

%% ===== Page 246 =====

%% [page 246 missing or failed: no entry]

\newpage

%% ===== Page 247 =====

%% [page 247 missing or failed: no entry]

\newpage

%% ===== Page 248 =====

%% [page 248 missing or failed: no entry]

\newpage

%% ===== Page 249 =====

%% [page 249 missing or failed: no entry]

\newpage

%% ===== Page 250 =====

induit par l'inclusion
$$D_x^* \hookrightarrow \mathcal{U}$$
ou encore, pour en termes d'automorph. d'un
foncteur fibre, comme l'hom. opération de
$T_x$ sur $\Gamma(D_x^{*\sim}, \mathcal{U}' \times_\mathcal{U} D_x^{*\sim})$ provenant

\begin{tikzcd}
& \mathcal{U}'_{D_x^{*\sim}} \ar[dl] \ar[d] \\
D_x^{*\sim} \ar[dr] & \mathcal{U}' \ar[d] \\
& \mathcal{U}
\end{tikzcd}

de l'opération de $T_x$ sur
$D_x^{*\sim}$. [unclear: text crossed out]

Supposons qu'on ait déjà
[MARGIN: $\pi_1(\mathcal{U},\tilde{\mathcal{U}}) = \pi$]
choisi (par ailleurs) un [unclear: foncteur fibre sur] $\mathcal{U}$,
par la donnée d' un pt base, ..., ou plus
généralement, d'un rev. universel $\tilde{\mathcal{U}}$ de $\mathcal{U}$

Pour le comparer avec $\pi_1(\mathcal{U}, D_x^{*\sim})$,
- il suffira de donner une "classe de chemins"
de $D_x^{*\sim} / \mathcal{U}$ à $\tilde{\mathcal{U}} / \mathcal{U}$ , i.e. un
isomorphisme

(62) $$ \tilde{\mathcal{U}} \xrightarrow{\sim} [unclear: Rev](\mathcal{U}, D_x^{*\sim}) $$

[MARGIN: l'image dans $\mathcal{U}$ du [unclear] de (62) donne dans $\mathcal{U}$ un point $y$ : [unclear]]

[unclear: crossed out paragraph at the bottom]
NB si $D_x^{*\sim}$ défini
par un rev. universel $\widetilde{D_x^*}$
de $D_x^*$, et si $\tilde{\mathcal{U}}$
défini par un pt base $y \in \mathcal{U}$,
alors [unclear]

\newpage

%% ===== Page 251 =====

773

( l'indétermination est un autom.
int. de $\Pi_1$, ou de $\Pi_1(\mathcal{U}, D_x^{* \sim})$ ) qui
induit un isomorphisme

(63) $$ \pi = \Pi_1(\mathcal{U}, \tilde{\mathcal{U}}) \overset{\sim}{\longleftrightarrow} \Pi_1(\mathcal{U}, D_x^{* \sim}) $$

(qui est changé par un autom. intérieur,
quand on change (62) par un él. de $\pi$).
Composant (63) et (60), on trouve donc
un hom.

(64) $$ T_n \longrightarrow \pi = \Pi_1(\mathcal{U}, \tilde{\mathcal{U}}) , $$

qui dépend du choix de (62) — et
qui, indépendamment de ce choix, définit
un hom. ext. bien déterminé —
qui n'est autre que l'hom. can.

(65) 
\begin{tikzcd}
\Pi_1(D_x^*) \ar[r] \ar[d, "\simeq"'] & \Pi_1(\mathcal{U}) \ar[d, equal] \\
T_n & \pi
\end{tikzcd}

[MARGIN: C) Relation avec
germes d'arcs (approche
originale "par pts base")]
Je vais décrire une façon commode
de fixer un rev. universel $\tilde{D}_x^*$ de $D_x^*$,
par un germe d'arc de courbe [unclear: analyt.]

\newpage

%% ===== Page 252 =====

%% [page 252 missing or failed: no entry]

\newpage

%% ===== Page 253 =====

774

d'où un rev. universel
(70) $$D_x^{*\sim} \overset{\text{déf}}{=} Rev_{univ}(D_x^*, [I], \varphi) \ .$$

Bien entendu, on aura une factorisation

\begin{tikzcd}
{[I]} \ar[r] \ar[rr, bend right, "\varphi"'] & D_x^{*\sim} \ar[r] & D_x^* \ (\hookrightarrow U)
\end{tikzcd}

donc un iso can

(71) $$Rev_{univ}(U, D_x^{*\sim}) \simeq Rev_{univ}(U, [I]) \ ,$$

où $I$ est considéré comme un germe
1-connexe sur $U$ par le morphisme composé
$$[I] \longrightarrow D_x^* \hookrightarrow U \ .$$

Donc . . .

(72) $$\pi_1(U, D_x^{*\sim}) = \pi_1(U, [I]) \overset{K_x}{\longleftarrow} T_x \ .$$

[MARGIN: D) Relation avec structure différentiable]
Lorsque $X$ est muni d'une structure
différentiable au voisi-nage de $x$, on dit
qu'un germe de courbe de $X$ sortant
de $x$ est différentiable, si l'application
continue
$$I \longrightarrow X$$
qui le définit est différentiable, et de plus

\newpage

%% ===== Page 254 =====

%% [page 254 missing or failed: no entry]

\newpage

%% ===== Page 255 =====

Soit d'autre part, dans le cas où $X$ est diff.
en $x$,

(76) $$ \underline{Couronne}_{diff}(X,-) $$

la sous-catégorie pleine de $\underline{Couronne}(X,x)$
formée des germes d'arcs différentiables. Soit
$S_x$
l'espace des demi-droites de
l'espace tangent $Tang_x(X)$

(77) $$ S_x = \text{espace des demi-droites de } Tang_x(X) \simeq \mathbb{S}^1 $$

et considérons le groupoïde fondamental

(78) $$ \Pi_1(S_x) \overset{\simeq}{\longrightarrow} \underline{Revunis}(S_x) $$ .

Je vais définir une équivalence de
catégories

(79)
\begin{tikzcd}
\underline{Couronne}_{diff}(X,x) \ar[r, "\simeq"] \ar[d, "\simeq"'] & \Pi_1(S_x) \ar[d, "\simeq"] \\
\underline{Revunis}(D_x^{*\sim}) & \underline{Revunis}(S_x)
\end{tikzcd}

qui, sur les ensembles d'objets, soit
l'application associant à tout germe d'arc
différentiable en $x$ sa demi-tangente. Il
s'agit donc, si $\gamma, \gamma'$ sont deux tels
arcs, correspondants aux demi-tangentes $\Delta, \Delta'$
de définir une bijection canonique

\newpage

%% ===== Page 256 =====

(80) 
\begin{tikzcd}
Isom(\gamma, \gamma') \ar[r, "\sim"] \ar[d, equal, "\text{déf}"'] & Isom_{\Pi_1 \mathbb{S}_n}(\Delta, \Delta') \ar[d, equal, "\text{déf}"] \\
Isom(D^{*\sim}_{n,\gamma}, D^{*\sim}_{n,\gamma'}) & \text{classes de chemins dans } \mathbb{S}_n \text{ de } \Delta \text{ à } \Delta'
\end{tikzcd}

Au lieu de définir directement la bijection
(80), je vais définir plutôt un
foncteur (qui sera une équivalence de
catégories)

(81) $$ \Phi : \Pi_1(\mathbb{S}_n) \xrightarrow{(\approx)} Revouv(D^*_n) \ , $$

Soit

(82) $A_n$ ens des parties fermées (1-connexes) de $\mathbb{S}_n$
($i.e.$ des $\underline{arcs}$ de cercle de $\mathbb{S}_n$,
mais $\neq \emptyset$,)

c'est un ensemble ordonné, $\underline{donc}$
lequel définit une catégorie ordonnée $\underline{A}_n$
(les objets sont les $a \in A_n$, les flèches
sont les inclusions) -- et on sait que
le groupoïde enveloppant $Gr(\underline{A}_n)$ de cette catégorie
ordonnée est canoniquement équivalent à $\Pi_1(\mathbb{S}_n)$,

\newpage

%% ===== Page 257 =====

%% [page 257 missing or failed: no entry]

\newpage

%% ===== Page 258 =====

%% [page 258 missing or failed: no entry]

\newpage

%% ===== Page 259 =====

%% [page 259 missing or failed: no entry]

\newpage

%% ===== Page 260 =====

%% [page 260 missing or failed: no entry]

\newpage

%% ===== Page 261 =====

%% [page 261 missing or failed: no entry]

\newpage

%% ===== Page 262 =====

(93)
\begin{tikzcd}
germ\, revdiff(X,x) \ar[r, "\simeq"', "\gamma \mapsto \Delta(\gamma)"] \ar[d, "\gamma \mapsto Revunis(D_x^*,\gamma) = D_{x,\gamma}^{*\sim}"'] & \Pi_1(S_x) \ar[d, "\simeq", "a \mapsto Revunis(S_x,-)"'] \\
Rev.unis(D_x^*) \ar[r, "\simeq"'] \ar[ur, dashed, "\Theta \simeq"'] & Revunis(S_x)
\end{tikzcd}

[MARGIN: Bien sûr, la difficulté en fait provenait de ce qu'on n'avait point d'interpr. géométrique explicite entre [unclear] l'espace $S_x$ ...]

Le seul foncteur qui n'ait pas une
explication est le foncteur
(94)
\begin{tikzcd}
Rev.unis.(D_x^*) \ar[r, "\simeq"] \ar[dr, "\Theta^{-1}"'] & Revunis(S_x) \\
& \Pi_1(S_x) \ar[u, "can"']
\end{tikzcd}

défini comme composé de $\Theta^{-1}$ avec
le foncteur can. $\Pi_1(S_x) \to Revunis(S_x)$.

Finalement, dans ces considérations
de géométrie topologico-différentielle, la
construction-clef est celle du foncteur
$\Theta$ (81). Nous avons pris quelques peines
à la définir de façon invariante. Dans
le cas où on dispose d'une iso locale
(95) $$f \colon (T_x(X),0) \to (X,x),$$ la
construction de $\Theta$ peut s'expliciter comme
une construction canonique dans un
espace vectoriel $V$. Toute demi-droite [unclear].

\newpage

%% ===== Page 263 =====

%% [page 263 missing or failed: no entry]

\newpage

%% ===== Page 264 =====

%% [page 264 missing or failed: no entry]

\newpage

%% ===== Page 265 =====

(97) 
\begin{tikzcd}
\Gamma(A, \tilde{S}_x/S_x) \ar[d, equal] & , & \Gamma(E_p(A), \tilde{D}_x^*/D_x^*) \ar[d, equal] \\
\operatorname{Hom}_{S_x}(A, \tilde{S}_x) \ar[d, equal, "\text{déf}"'] & & \operatorname{Hom}_{D_x^*}(E_p(A), \tilde{D}_x^*) \ar[d, equal, "\text{déf}"'] \\
\tilde{S}_x(A) & & \tilde{D}_x^*(E_p(A))
\end{tikzcd}

qui sont l'un et l'autre des Torseurs sur le groupe $T_{X,x} = T_x$ le $\mathbb{Z}$-module des orientations de $X$ en $x$ (pour $\sim$). Considérons l'ens

(98) $$H(A) = Isom_{T_x}(\tilde{S}_x(A), \tilde{D}_x^*(E_p(A)))$$

qui est lui-même un $T_x$-torseur. Soit en effet une inclusion

$$A \subset A'$$

on en déduit, par restriction, des isomorph. de $T_x$-torseurs

(99) $$\tilde{S}_x(A') \simeq \tilde{S}_x(A), \quad \tilde{D}_x^*(E_p(A')) \simeq \tilde{D}_x^*(E_p(A))$$

d'où un isom. canonique

(100) $$H(A) \simeq H(A')$$

de sorte que 

(101) $$A \mapsto H(A), \text{ foncteur } A_x \xrightarrow{H} \underline{Tors}(T_x)$$

où $A_x$ est la catégorie ordonnée [unclear] associée à l'ens. ordonné $A_x$. On trouve donc par factorisation un foncteur de la catégorie groupoïde associé $Gr(A_x)$ de $A_x$ dans $\underline{Tors}(T_x)$,

(102) 
\begin{tikzcd}
A_x \ar[d] \ar[r, "H"] & \underline{Tors}(T_x) \\
Gr(A_x) \ar[ru, "Gr(H)"'] & 
\end{tikzcd}

\newpage

%% ===== Page 266 =====

%% [page 266 missing or failed: no entry]

\newpage

%% ===== Page 267 =====

781

-ment par la commutation des diagrammes

(107)
\begin{tikzcd}
H(A) \ar[r, "i_A", "\sim"'] \ar[d, "(100)"'] & H(\tilde{S}_\alpha, \tilde{D}_\alpha^*) \\
H(A') \ar[ur, "\sim"', "i_{A'}"] &
\end{tikzcd}

Pour vérifier le fait que le foncteur $Gr(H)$
de (102) est [unclear: bien consistant], je considère les
deux foncteurs (dépendant du choix de
$\tilde{S}_\alpha, \tilde{D}_\alpha^*$)

(108)
$$ \begin{cases} 
A \longmapsto \tilde{S}_{\alpha,A} \quad , \quad A \longmapsto \tilde{D}_{\alpha,A}^* \\ 
\underline{A}_\alpha \longrightarrow \underline{Rev}_{univ}(S_\alpha) \quad \underline{A}_\alpha \longrightarrow \underline{Rev}_{univ}(D_\alpha^*) 
\end{cases} $$

qui se factorisent en des équivalences

(109)
$$ Gr(\underline{A}_\alpha) \xrightarrow{\sim} \underline{Rev}_{univ}(S_\alpha) \quad , \quad Gr(\underline{A}_\alpha) \xrightarrow{\sim} \underline{Rev}_{univ}(D_\alpha^*) $$
utilisant $\tilde{S}_\alpha, \tilde{D}_\alpha^*$.

Ceci dit, pour un $A \in Ob \, \underline{A}_\alpha = A_\alpha$ [unclear: donné],
des isom. fonctoriels en $A$

(110)
$$ \begin{cases} 
[unclear](A) \simeq \underline{Isom}_{\underline{Rev}_{univ}(S_\alpha)}(\underbrace{F(A)}_{\tilde{S}_{\alpha,A}}, \tilde{S}) \\ 
[unclear](E,A) \simeq \underline{Isom}_{\underline{Rev}_{univ}(D_\alpha^*)}(\underbrace{G(A)}_{\tilde{D}_{\alpha,A}^*}, \tilde{D}_\alpha^*) 
\end{cases} $$

et par la propriété univ. de $Gr(\underline{A}_\alpha)$,
ces isom. s'étendent sur la
catégorie $Gr(\underline{A}_\alpha)$. Dans la situation où
elle se réduit à un groupe profini $T (= T_\alpha)$
trois $T$-groupoïdes connexes $g, g', g''$

\newpage

%% ===== Page 268 =====

%% [page 268 missing or failed: no entry]

\newpage

%% ===== Page 269 =====

7P2

Définition. — On dit qu'on a une corrélation entre un rev. univ. $\tilde{S}_x$ de $S_x$, et un rev. univ. $\tilde{D}_x^*$ de $D_x^*$, quand on s'est donné un élément de $H(\tilde{S}_x, \tilde{D}_x^*)$.

Donc, si $A \in A_a$, la donnée d'une corrélation revient à la donnée d'un élément de $H(A)$, i.e. d'un isomorphisme entre les $T_a$-torseurs $\tilde{S}_x^\sim(A)$ et $\tilde{D}_x^{*\sim}(E_p(A))$. Cette identification est compatible avec les flèches d'inclusion entre $A \subset A'$. Les corrélations entre $\tilde{S}_x$ et $\tilde{D}_x^*$ forment un torseur sous $T_a$. Pour $\tilde{S}_x$ fixé

Soit $Corr_x(X) = Cour_x(X)$ la catégorie des triplets $(\tilde{S}_x, \tilde{D}_x^*, h)$ de rev. univ. corrélés de $S_x$ et de $D_x^*$, on a des foncteurs $(\tilde{S}_x, \tilde{D}_x^*, h) \rightrightarrows {\tilde{S}_x \atop \tilde{D}_x^*}$

(117)
\begin{tikzcd}
& Corr_x(X) \ar[dl, "\simeq"'] \ar[dr, "\simeq"] & \\
Revuniv(S_x) & & Revuniv(D_x^*)
\end{tikzcd}

qui sont des équivalences de catégories.

\newpage

%% ===== Page 270 =====

%% [page 270 missing or failed: no entry]

\newpage

%% ===== Page 271 =====

78s

images inverses par $u_S$ et par $u_{D^*}$ de $\tilde{S}_{x'}^\sim$ et de $\tilde{D}_{x'}^{* \sim}$
soient $\tilde{S}_x^{\prime \sim}$, $\tilde{D}_x^{* \prime \sim}$, sont munis d'une covariation
que je désigne par $h_x^\prime$. Ceci posé, les theoremes
d'iso. de $T_x$-torseurs canoniques

$$ (121) \qquad F(u) \simeq Isom_{\tilde{S}_x}(\tilde{S}_x^\sim, \tilde{S}_x^{\prime \sim}) $$
$$ G(u) \simeq Isom_{\tilde{D}_x^*}(\tilde{D}_x^{*\sim}, \tilde{D}_x^{* \prime \sim}) $$

Or par le diagramme d'équivalence (117) on
a bien un iso de $T_x$-torseurs
$$ Isom((\tilde{S}_x^\sim, \tilde{D}_x^{*\sim}, h_x), (\tilde{S}_x^{\prime \sim}, \tilde{D}_x^{* \prime \sim}, h_x^\prime)) $$

(122)
\begin{tikzcd}
& Isom((\tilde{S}_x^\sim, \tilde{D}_x^{*\sim}, h_x), (\tilde{S}_x^{\prime \sim}, \tilde{D}_x^{* \prime \sim}, h_x^\prime)) \ar[dl] \ar[dr] & \\
F(u) \simeq Isom(\tilde{S}_x^\sim, \tilde{S}_x^{\prime \sim}) & & G(u) \simeq Isom(\tilde{D}_x^{* \sim}, \tilde{D}_x^{* \prime \sim})
\end{tikzcd}

d'où la bijection canonique cherchée (120) —
[unclear: des qui donne un justification (96).]
La bijection canonique (120) a la
propriété $\pm$ évidente suivante, qui la caractérise:

(123) $\begin{cases}
\text{Soient } u_{\tilde{S}} \text{ et } u_{\tilde{D}^*} \text{ des prolongements se correspondants} \\
\tilde{S}_x^\sim \xrightarrow{\sim} \tilde{S}_{x'}^\sim, \ \tilde{D}_x^{*\sim} \xrightarrow{\sim} \tilde{D}_{x'}^{* \sim}, \text{ associés par (120), soit } A \\
\in S_x, \ A' \in S_{x'} \text{ son image par } u_S \text{, soit } \alpha \text{ une} \\
\text{section de } \tilde{S}_x^\sim \text{ sur } A, \ \beta \text{ la section de } \tilde{D}_x^{*\sim} \text{ sur } Ep(A) \\
\text{associé par la covariation } h_x, \text{ soit } \alpha' \text{ la section} \\
\text{de } \tilde{S}_{x'}^\sim \text{ sur } A' \text{ déduite de } \alpha \text{ en appliquant } u_{\tilde{S}}, \\
\beta' \text{ la section de } \tilde{D}_{x'}^{* \sim} \text{ sur } Ep(A') \text{ déduite de } \beta \\
\text{en appliquant } u_{\tilde{D}^*}, \text{ alors } \alpha' \text{ et } \beta' \text{ sont associés} \\
\text{par la covariation } h'.
\end{cases}$

\newpage

%% ===== Page 272 =====

%% [page 272 missing or failed: no entry]

\newpage

%% ===== Page 273 =====

%% [page 273 missing or failed: no entry]

\newpage

%% ===== Page 274 =====

%% [page 274 missing or failed: no entry]

\newpage

%% ===== Page 275 =====

%% [page 275 missing or failed: no entry]

\newpage

%% ===== Page 276 =====

%% [page 276 missing or failed: no entry]

\newpage

%% ===== Page 277 =====

%% [page 277 missing or failed: no entry]

\newpage

%% ===== Page 278 =====

qui prend ainsi le nom de groupe de
Teichmüller spécial de $(X,S)$. Pour
$G$ quelconque et action de $G$ sur
$(X,S)$ orientée : la donnée d'un hom.
de groupes topologiques
$$G \to \mathcal{A}$$
s'intègre dans un diagramme de
groupes top.

(141)
\begin{tikzcd}
G \ar[r] \ar[d] & \mathcal{A} \ar[d] \\
\pi_0(G) \ar[r] & \tau = \pi_0(\mathcal{A})
\end{tikzcd}

et on trouve alors, par fonctorialité,
que le $\tilde{G}$ construit précédemment
est l'image inverse de l'ext. $\tilde{\tau}$ de $\tau$
par $\mathbb{Z}^I$, via $G \to \tau$.

Ces préliminaires étant posés,
revenons à la situation où $X$
admet une structure diff. $C^\infty$ au voisin-
age de $S$, et soit pour tout $x \in S$,

\newpage

%% ===== Page 279 =====

%% [page 279 missing or failed: no entry]

\newpage

%% ===== Page 280 =====

%% [page 280 missing or failed: no entry]
